% 本文件由 LaTeX 排版 Agent 根据有效 translation.json 自动生成。
% author=Tertullian; work_id=0310; title=论灵魂

\chapter{论灵魂}

% structure: p0001 source=h1 target=omitted reason=page-title-already-rendered-by-parent

% structure: p0002 source=h2 target=section reason=relative-heading-rank; base=h2

\section{第一章 关于灵魂，我们不应求助于哲学家，而应求助于天主}

在与\Person{赫尔谟根尼}讨论了灵魂起源的单一论点之后，按其假设引导我，即灵魂更多地在于对物质的适应，而非天主的嘘气，我现在转而讨论与此主题相关的其他问题；（在讨论这些问题时）我显然将主要与哲学家们争辩。在\Person{苏格拉底}的监狱里，他们就灵魂的状态进行过小规模论战。我当即怀疑，那时间是否适宜于他们（伟大）的老师——（更不必说地点），尽管\Emph{那}也许并不太重要。因为那时\Person{苏格拉底}的灵魂能清晰而平静地默观什么呢？圣船已（从\Place{提洛}）返回，他被判处的毒芹汁已经饮下，死亡已在他面前：（他的心灵）可以想见，因每一情绪而自然激动；或者，若本性已丧失其影响，那它必定已被剥夺了一切思维力量。或者，任凭你设想它多么平静安详，毫不动摇，不顾自然情感的呼唤，无视那即将成为遗孀者的眼泪，以及从此成为孤儿的子女的景象，但他的灵魂甚至在其极力抑制情感的努力中，也必定受到触动；他的坚毅本身也必定动摇，当他与周围激动的干扰抗争时。此外，一个被不义判处死刑的人，除了那些能抚慰他所受伤害的想法，还能有什么别的念头？尤其对那荣耀的生灵——哲学家而言，更是如此，对之而言，不公的待遇不会引发对安慰的渴望，而是激起怨恨和愤怒的感情。因此，在他被判决后，当他的妻子带着她软弱的哭喊来到他面前：\Quote{啊，\Person{苏格拉底}，你被不义地判处了死刑！}他仿佛已经欢喜地回答说：\Quote{那你倒希望我被公义地判处死刑吗？}因此，若即使在监狱里，出于挣脱\Person{阿尼图斯}和\Person{墨勒图斯}肮脏之手的渴望，他在死亡面前，以强烈的假设断言灵魂不死，正如需要这样的假设来挫败他们加于他的不义，这并不足为奇。所以，\Person{苏格拉底}在那时刻的所有智慧，都来自一种故作镇定的做作，而非对已确证真理的坚定信念。因为，不藉着天主，谁曾发现过真理？不藉着基督，谁曾找到过天主？不藉着圣神，谁曾探究过基督？不藉着信德的神秘恩赐，谁曾获得过圣神？\Person{苏格拉底}，无人能怀疑，是被另一种神所驱动。因为据说，从童年起就有一个魔鬼依附于他——诚然是最坏的老师，尽管诗人和哲学家赋予它崇高的地位——甚至仅次于（不，与诸神本身并列）。基督德能的教训尚未赐下——唯有这德能能够驳斥这最有害的邪恶影响，它毫无良善，反而是一切错误的始作俑者，一切真理的诱惑者。现在，如果\Person{苏格拉底}被皮提亚魔鬼的神谕宣告为最智慧的人——你可以肯定，这魔鬼为其朋友巧妙地安排了此事——那么，基督徒的智慧宣称具有多么大的尊严和坚毅，在其气息面前，众魔鬼的整个军团便四散！这天国学派的智慧坦率而毫无保留地否认今世的众神，并不显示这样的矛盾，即命令\Quote{将一只公鸡祭献给\Person{埃斯科拉庇俄斯}:}：它不引入新的神和魔鬼，而是驱逐旧有的；它不腐蚀青年，而是教导他们一切良善和节制；因此，为了那真理的缘故——那真理越是完满，就招致更大的仇恨——它承受的不仅是一个城市的不义判决，而是整个世界的；它品尝死亡，并非几乎以欢愉的方式从（毒）杯中饮下；而是在绞架和火祭中，以各种残酷的苦痛，饮尽死亡。与此同时，在世界上更阴郁的监狱里，在你们的\Person{克贝斯}们和\Person{斐多}们中间，在关于（人的）灵魂的每一次探究中，它依照天主的法则进行探询。无论如何，你们不能给我们指出比灵魂的作者更权威的解释者。从天主那里，你可以学到你所拥有自天主的；但若你不从天主那里得到，你绝不可能从别处得到这知识。因为谁能揭示天主所隐藏的呢？我们在探询中必须求助于那一位，即使我们的无知也由此最安全地获得。因为，对于一件事，因为天主没有启示给我们而不知道，比起因为人的智慧（\Emph{他}竟敢假设知道）而知道，对我们实在更好。\par

% structure: p0004 source=h2 target=section reason=relative-heading-rank; base=h2

\section{第二章 基督徒对当前主题拥有确定而简单的知识}

当然我们不否认哲学家有时思考与我们相同的事物。真理的见证就是其结果。即使在风暴中，当天与海的界限在混乱中消失时，有时也出于某种幸运的偶然，某港口（历经辛苦的船）偶然发现；有时在深夜里，仅凭借盲目的运气，一个人找到进入某地或离开某地的通道。然而，在自然中，大多数结论可以说是由那种共同智慧提示的，天主乐于用那智慧赋予人的灵魂。这种智慧被哲学拾起，并且为了荣耀她自己的艺术，被膨胀了（我用这种语言并不奇怪），她竭力追求那种语言的便利，这种便利在一切事物的建立与摧毁中被运用，并且更擅长用言语说服人而非教导。她给事物赋予形式和条件；有时使它们公共和公开，有时将它们据为私用；她在确定之物上任意打上不确定性的印记；她援引先例，仿佛一切事物都可以相互比较；她用规则和定义描述一切事物，甚至为相似的对象分配不同属性；她将任何事都不归因于天主的许可，而是将自然法则假定为她的原则。我可以容忍她的虚饰，只要她自己忠于自然，并且向我证明她因参与自然的创造而拥有对自然的掌控。她无疑认为她是从神圣的源泉中获取她的奥秘，如同人们认为的那样，因为在古代，大多作者被视为（我不说神一般的，而是）实际上为神：例如，埃及的\Person{默丘利}，\Person{柏拉图}对他极为敬重；弗里吉亚的\Person{西勒诺斯}，\Person{米达斯}曾把长耳借给他，当牧人带他来见他时；\Person{赫莫提穆斯}，\Place{克拉佐美奈}的善良人民在他死后为他建了一座神庙；\Person{俄耳甫斯}；\Person{穆赛俄斯}；以及\Person{毕达哥拉斯}的老师\Person{斐瑞居得}。但我们何必在意，既然这些哲学家也攻击那些被我们称为伪经的著作，我们确信，凡与现今时代兴起的真确先知体系不符的，都不应接受；因为我们没有忘记曾有假先知，而在他们之前很久就有堕落的神灵，他们以这种（\Emph{哲学的}）狡猾知识教导了整个世界的基调与面貌？确实，任何寻求智慧的人出于好奇甚至去咨询先知本人，也并非不可信；（\Emph{但无论怎样}），如果你看哲学家，你会发现他们中相异多于一致，因为即便在他们的一致中也可发现相异。无论什么在他们\Emph{的体系中}是真实的且符合先知智慧的，他们要么推荐为源于其他来源，要么歪曲地应用于其他意义。当他们假装真理要么由虚假所助，要么虚假从中得到支持时，这个过程给真理带来极大损害。以下情形必然使我们与哲学家争执，尤其在本问题上，即他们有时将双方共有的情感包裹在自身特有的论证中，这些论证在某些点上与我们的信德规则和标准相悖；另一些时候则用双方都承认有效的、偶尔符合其信仰体系的论证来辩护特别属于他们自己的意见。真理就这样几乎被哲学家排除了，因为他们用毒药感染了它；因此，如果我们考虑\Emph{我们刚才提到的}两种结合模式，它们同样敌视真理，我们感到迫切的需要：一方面，将我们与他们共有的情感从哲学家的论证中解放出来；另一方面，将双方使用的论证与同一个哲学家的意见分离开来。\Emph{而这我们可以做到}，通过将一切问题回溯到天主默感的标准，明显排除那些不受任何先入为主之见纠缠的简单情况，人可以单纯根据\Emph{人的}见证合理地接受；因为这类明证我们有时必须从对手那里借用，当对手从中一无所获时。我并非不知道哲学家们关于眼前这个主题，在他们自己的注释中积累了何等大量的文献——有何等不同的原则学派、何等意见冲突、何等丰富的发问来源、何等令人困惑的解决之道。此外，我还查看了医学，即哲学（如他们所说）的姊妹，她声称其职能是治疗身体，并借此对灵魂有特别的认识。由此她与她的姊妹有巨大分歧，因为后者假装通过对灵魂在其\Emph{身体的}寓所中的更明显处理而更了解灵魂。但不必介意它们之间的这番争胜！为了延伸各自对灵魂的研究，一方面，哲学尽享其天才的全部范围；另一方面，医学则拥有其技艺与实践的严苛要求。人们对不确定性的探求是广泛的；对推测的争论则更为广泛。无论举证的困难有多大，产生信服的辛劳也丝毫不减；所以忧郁的\Person{赫拉克利特}说得完全正确，当他观察到笼罩灵魂探寻者研究的浓密黑暗，并厌倦于他们无休止的问题时，他宣布他确实没有探索到灵魂的界限，尽管他已走遍了每一条\Emph{在她的领域内的}道路。然而，对于基督徒而言，只需寥寥数语就能清楚理解整个主题。但在这寥寥数语中为他总能产生确定性；他也不被允许将探求扩展到超出解决范围的程度；因为宗徒禁止\Quote{无休止的问题}。\BibleRef{弟前 1:4}然而必须补充，人只能找到从天主所学到的解决之道；而从天主所学到的，便是整件事情的总结与实质。\par

% structure: p0006 source=h2 target=section reason=relative-heading-rank; base=h2

\section{第三章 根据圣经的简明语句确定灵魂的起源}

\Quote{但愿从未有过异端是必要的，好叫那些被认可的人得以显明出来！}\BibleRef{格前 10:19} 这样，我们就永远不必为灵魂的问题与哲学家们——这些可被恰当地称为异端之父的人——进行较量。宗徒早在自己的时代就已预见到，哲学将严重损害真理。\BibleRef{哥 2:8} 他在雅典之后，结识了那座\Emph{好辩}的城市，并在那里尝到了那些贩卖智慧的聪明人和空谈家的滋味后，才受到激励而发出\Emph{关于虚假哲学}的劝诫。同样，按照人的诡辩学说\Quote{将水掺入酒中}\BibleRef{依 1:22} 来处理灵魂问题也是如此。他们中有人否认灵魂的不朽；有人肯定灵魂不朽，甚至更多；有人争论其实体；有人争论其形式；还有人争论其各种能力。一个哲学派别从不同来源推导其状态，另一个则将其归宿归于不同之处。\Emph{各学派反映了其大师的特征}，正如他们从\Person{柏拉图}的尊贵、\Person{芝诺}的力量、\Person{亚里士多德}的镇定、\Person{伊壁鸠鲁}的愚蠢、\Person{赫拉克利特}的忧郁或\Person{恩培多克勒}的疯狂中接受了印象。我猜想，神圣教义的过失在于它来自\Place{犹太}\BibleRef{依 2:3}，而非\Place{希腊}。基督也犯了错误，他派遣渔夫而非诡辩家去传道。因此，无论哲学散发出何种有害的雾气，遮蔽了真理的清新健康氛围，基督徒都应当清除它们，既要粉碎从万物原理——即哲学家的原理——中得出的论证，又要以天上智慧的原则——即主所启示的——来对抗它们，以便既消除哲学俘虏外教人的陷阱，又压制异端动摇基督徒信仰的手段。在本论文的开头，我们与\Person{赫尔默根尼}的争论中已经决定了一个要点，即我们主张灵魂是由天主的嘘气形成的，而非出自物质。在那里，我们甚至依靠圣经启示的明确指示，即\BibleQuote{上主天主向人的鼻孔吹了一口气，人就成了一个有灵的生物}\BibleRef{创 2:7}——当然是藉着天主的嘘气。因此，这一点上我们无需进一步探讨或阐述。它有自己的专论和自己的异端。我将把它作为我进入主题其他部分的引子。\par

% structure: p0008 source=h2 target=section reason=relative-heading-rank; base=h2

\section{第四章 反对柏拉图：灵魂是在出生时被创造和产生的}

确定了灵魂的起源后，接下来是其状态或状况。因为我们承认灵魂源于天主的嘘气，所以自然要为其赋予一个开端。\Person{柏拉图}确实拒绝赋予灵魂开端，因为他认为灵魂是无生无造的。然而，我们恰恰从其有开端以及其性质出发，教导灵魂既有出生也有受造。当我们同时赋予灵魂出生和受造时，我们并没有犯错误：因为\Emph{出生}和\Emph{受造}是不同的——前者是最适用于生物的术语。然而，当区分各有其时间和地点时，它们之间偶尔也具有互换性。因此，受造可以被理解为出生；因为任何以任何方式获得\Emph{存在}或\Emph{实有}的事物，实际上都是被产生的。因为造物主确实可以被称作受造物的父亲：在此意义上，\Person{柏拉图}也使用这种措辞。因此，就我们相信灵魂受造或出生而言，哲学家的意见已被预言本身的权威所推翻。\par

% structure: p0010 source=h2 target=section reason=relative-heading-rank; base=h2

\section{第五章 斯多葛派的可能观点：灵魂具有物质性本质}

假設某人召來一位\Person{歐布洛斯}、一位\Person{克裡托勞斯}、一位\Person{芝諾克拉底}，此次還有\Person{柏拉圖}的朋友\Person{亞里士多德}作為幫手。他們很可能已準備好剝奪靈魂的形體性，除非他們碰巧看到其他哲學家——而且人數更多——與他們的目的相反，主張靈魂具有形體本質。我這裡不僅指那些將靈魂模塑自明顯有形體質料的人，如\Person{希帕庫斯}和\Person{赫拉克利特}（以火），\Person{希朋}和\Person{泰勒斯}（以水），\Person{恩培多克勒}和\Person{克里提亞斯}（以血），\Person{伊壁鳩魯}（以原子），因為原子也因其結合形成有形體的團塊；\Person{克裡托勞斯}及其逍遙學派（以某種難以描述的\OriginalTerm{第五元素}{quintessence}），如果那可以稱為形體，其實它更包含並包裹有形體質料；——我也召喚斯多葛派來幫助我，他們雖然幾乎用與我們相同的術語宣稱靈魂是一種精神的本質（因為氣息與精神在本性上彼此非常接近），但仍會毫不困難地說服我們靈魂是有形體的本質。的確，\Person{芝諾}將靈魂定義為與（身體）一同產生的精神，其論證方式是：那因其離去而導致生物死亡的本質是有形體的。現在，生物因與（身體）一同產生的精神之離去而死亡；因此，與（身體）一同產生的精神是有形體的本質。但這與（身體）一同產生的精神就是靈魂：由此得出，靈魂是有形體的本質。\Person{克萊安西斯}也主張，父母與子女之間的家族相似不僅體現在身體特徵上，也體現在靈魂特質上；彷彿透過（人的）行為、能力和情感的一面鏡子，身體的相似與不相似也被靈魂捕捉和反映。因此，正因靈魂是有形體的，它才容易接受相似與不相似。此外，有形體與無形體的事物在感受性上沒有任何共同之處。但靈魂確實與身體共鳴，並在身體被挫傷、創傷和潰瘍所傷時分擔其痛苦；身體也與靈魂一同受苦，並在靈魂因焦慮、痛苦或愛而失去活力時與之結合，並通過自己的羞紅和蒼白來見證靈魂的羞恥和恐懼。因此，從這種感受性的互通可以證明靈魂是有形體的。\Person{克呂西波}也與\Person{克萊安西斯}攜手，他認定，有形體的事物完全不可能與無形體的事物分離，因為它們沒有相互接觸或凝聚的關係。相應地，\Person{盧克萊修}說：\par

\Quote{Tangere enim et tangi nisi corpus nulla potest res.}\par

\Quote{因為除了形體，沒有任何東西能夠觸摸或被觸摸。}\par

（然而，這種分離在靈魂與身體之間是相當自然的）；因為當身體被靈魂拋棄時，它就被死亡征服。因此，靈魂被賦予了形體；因為如果它不是有形體的，就不能離開身體。\par

% structure: p0015 source=h2 target=section reason=relative-heading-rank; base=h2

\section{第六章　\Person{柏拉圖}派為靈魂非形體性所作之論證，或許是膚淺的}

柏拉图主义者更多地以精妙而非真理来搅扰这些结论。他们声称，每一个形体必定具有有灵的本性或无灵的本性。若它具有无灵的本性，则其运动来自外部；若具有有灵的本性，则来自内部。然而，灵魂既不从外部也不从内部接受运动：非从外部，因为它不具有无灵的本性；亦非从内部，因为它本身是赋予身体运动者。那么，显而易见，它不是有形实体，因为根据有形实体的本性和法则，它两种方式都不接受运动。这里首先令我们惊讶的是，这一定义诉诸与灵魂毫无亲缘关系的对象，这是不恰当的。因为灵魂不可能被称为有灵的身体或无灵的身体，因为正是灵魂本身使身体成为有灵的（若它临在于身体），或成为无灵的（若它离开身体）。因此，产生结果者本身不能是结果，从而有资格被称为有灵之物或无灵之物。灵魂是就其自身的实体而被如此称呼的。那么，若灵魂本身不能被称作有灵的身体或无灵的身体，它如何能与有灵和无灵身体的本性与法则相提并论呢？再者，既然身体的特征是被其他事物从外部推动，而且我们已证明灵魂受到先知性的影响或疯狂的影响时，是从其他事物接受运动（当然是从外部由其他事物推动），因此，根据我们所举的例子，凡被外部其他对象推动的，我必正确地视之为有形实体。那么，若从其他事物接受运动是身体的特性，那么将运动赋予他物更是如此！然而灵魂推动身体，身体的一切活动显然都是外部的、来自外部的。正是灵魂赋予脚以行走的运动，赋予手以触觉，赋予眼以视觉，赋予舌以言语——一种内在的形象，推动并活化表面。若灵魂是无形体，它从何处获得如此能力？一个无实质之物如何能推动固体对象？但人的感官如何似乎可分为有形体的和理智的类别呢？他们告诉我们，有形事物（如土与火）的性质由身体的感官（触觉和视觉）指示；而无形体事物（例如仁慈与恶行）的性质则由理智能力发现。由此他们推断出（在他们看来）明显的结论：灵魂是无形体的，其特性不是由身体器官的知觉，而是由理智能力来理解。好吧，我若不立刻削去他们论点所立足的根基，那才奇怪呢。因为我向他们指明，无形体的事物如何通常被提交给身体的感觉——例如，声音提交给听觉器官；颜色提交给视觉器官；气味提交给嗅觉器官。并且，正如这些情形，灵魂同样与身体接触；更不用说，无形体对象之所以通过身体器官报告给我们，正是因为它们与这些器官接触。因此，既然显然甚至无形体对象也能为有形对象所包容和理解，那么为什么灵魂（是有形体的）不能同样地被无形体的能力所理解和认识呢？由此可见，他们的论点确实失败了。在他们更显著的论点中，有这一条：他们断定每个有形实体都由有形实体滋养；而灵魂作为无形体本质，则由无形体的滋养品——例如智慧的研究——来滋养。但这一根基也不稳固，因为\Person{Soranus}——一位在医学上极有造诣的权威——给我们提供了答案，他断言灵魂甚至由有形食物滋养；事实上，当灵魂衰败虚弱时，常常确实被食物所复原。确实，若被剥夺一切食物，灵魂岂不彻底离开身体吗？于是，\Person{Soranus}在极其详尽地论述了灵魂（写了四卷专题论著）并仔细权衡了所有哲学家的意见之后，捍卫了灵魂的有形体性，尽管在此过程中他剥夺了灵魂的不朽性。因为并非所有人都被赐予相信基督徒所享有的真理。因此，既然\Person{Soranus}从事实向我们表明灵魂是由有形食物滋养的，那么就让哲学家（采用类似的证明方式）表明灵魂是由无形食物维持的。但事实是，从来没有人能用\Person{Plato}精妙辞藻的蜜水来止息这个人对灵魂状况的疑惑和困难，也没有人能用\Person{Aristotle}琐碎秘方的碎屑使他满足。但那些强壮野蛮人的灵魂又将如何呢？他们确实未曾受过哲学家学问的滋养，却在未受教化的实践智慧中强壮，虽然他们在哲学上极为匮乏，没有你们的雅典学园和柱廊，甚至没有\Person{Socrates}的监狱，却仍然设法活着。因为受学问研究的滋养受益的并非灵魂的实际实体，而仅仅是其行为和训练；此种滋养并不增加其体积，而仅增添其优雅。此外，斯多葛派断言甚至连技艺也有形体性，这是一件幸运的事；因为照此推论，灵魂也必为有形体的，因为人们通常认为灵魂是由技艺滋养的。然而，哲学家的思想如此巨大的执迷，以至于通常不能直视眼前之物。因此有\Person{Thales}掉进井里的轶事。它也常常由于甚至不理解自己的意见而怀疑自己健康的失灵。因此有\Person{Chrysippus}与藜芦的故事。我想，当他断言两个身体不可能被包含在一个之中时，他必定发生了某种这样的幻觉：他必定忽略了那些怀孕妇女的情况，她们日复一日，在一个子宫的拥抱中，生产的不是一个身体，而是两个甚至三个。同样，在民法记载中，我们发现一位希腊妇女的例子，她一次分娩生下了五个孩子，所有这些孩子是一位母亲一次分娩所生，同一胎的众多父母，来自同一子宫的丰饶产物，她被这么多身体护卫——我几乎要说，一个人群——她自己至少是第六个人！整个受造界都证明，那些自然注定从身体中产出的身体，已经（被包含）在它们所出自的那个身体中。出自其他事物者，必居其次。然而，除非通过生育的过程，没有什么从其他事物中产出；但即便如此，它们是两个（事物）。\par

% structure: p0017 source=h2 target=section reason=relative-heading-rank; base=h2

\section{第七章。灵魂的形体性由福音所证明。}

关于哲学家，我们已经说得够多了。至于我们自己的导师，确实，我们对他们的引用是\OriginalTerm{超出必要}{ex abundanti}——权威的冗余：在福音本身中，他们会发现关于灵魂形体性的最清楚证据。在地狱里，一个人的灵魂在受苦，被火焰惩罚，遭受极度的口渴，并从一个更幸福之人的指头，为他的舌头恳求一滴水的安慰。\BibleRef{路 16:23}你以为这有福的穷人和可怜的富人的结局只是想象的吗？那么，为什么在这个叙述中有拉匝禄的名字，如果这情况不是（在）真实事件（的范畴）内？但即使它被视为想象，它仍然是对真理和现实的见证。因为除非灵魂拥有形体性，否则一个灵魂的形象不可能包含一个具有形体实质的手指；圣经也不会虚构关于身体肢体的陈述，如果这些都是不存在的。但在身体分离后被送到阴府的是什么东西？在那里被拘留的是什么？被保留到审判之日的又是什么？基督在死时也下降到那里？我想那是先祖们的灵魂。但为什么（这一切），如果灵魂在地下居所中毫无所有？因为它肯定是一无所有，如果不是有形的实质。因为任何无形的东西都无法以任何方式被保留和看守；它也不受惩罚或安慰。那一定是身体，通过它惩罚和安慰才能被经验。关于这一点，我将在更合适的地方更详细地讨论。因此，无论灵魂在阴府中，在它的监狱或寓所里，在火中或在亚巴郎的怀里，尝到多少惩罚或安慰，都因此证明它自己的形体性。因为无形的东西不受任何苦，没有使之能够受苦的东西；否则，如果它有这种能力，它一定是有形的实质。因为每一个有形的东西都能受苦，那么能够受苦的也是形体。\par

% structure: p0019 source=h2 target=section reason=relative-heading-rank; base=h2

\section{第八章。其他柏拉图主义论点之考量。}

此外，以任何事物不完全像有形类别的其他构成成分为由，将其排除在有形存在之外，将是一种严厉且荒谬的做法。而当个体受造物拥有各种特性时，难道同类别作品中的这种多样性不显明造物主的伟大，使它们同时不同却又相似，彼此友好却又竞争？事实上，哲学家们自己也同意宇宙由和谐的对立构成，根据恩培多克勒的（理论）友谊与敌意。因此，虽然有形本质与无形本质相对，但它们彼此之间的差异却使它们通过多样性扩大自己的种类，而不改变其类属，仍然都保持形体；以其多样性在多种存在中归光荣于天主；因其差异而如此多样；如此不同，以至于它们中一些拥有一种知觉，另一些拥有另一种；一些以一种食粮为生，另一些以另一种为生；此外，一些具有可见性，而另一些不可见；一些沉重，另一些轻盈。他们习惯说：灵魂必须因此被宣称为无形的，因为死人身体在灵魂离开后变得更重，而它们本应更轻，因为失去了一个身体的重量的灵魂是一个有形的实质。但是，索拉努斯（在回答这一论点时）说：如果人们否认海是有形的，因为船出水后变成沉重不动的块状物，那又怎样？那么，灵魂的形体本质就更真实、更强有力了，它携带身体，最终以最敏捷的运动承担如此大的重量！再者，即使灵魂是不可见的，这也恰恰符合其自身形体性的条件，符合其自身本质的特性，以及那些其命运使之看不见灵魂者的性质。猫头鹰的眼睛无法忍受太阳，而老鹰却能面对它的荣耀，以致其幼崽的崇高品格由它们凝视的坚定力量决定；而鹰雏如果避开太阳光，就被作为退化的生物赶出巢穴！因此，对一只眼睛不可见的事物，对另一只眼睛却可能清晰可见——这并不意味着那没有被赋予同样强大力量（视觉）的东西是无形体。太阳确实是有形实体，因为它是（由）火构成的；然而，鹰雏立即承认其存在的物体，猫头鹰却否认，但这并不损害老鹰的见证。对于灵魂的形体性，也有同样的差别，它（也许）对肉体不可见，但对精神却完全可见。于是，若望在\Quote{圣神内}，\BibleRef{默 1:10}清楚地看见了殉道者的灵魂。\BibleRef{默 6:9}\par

% structure: p0021 source=h2 target=section reason=relative-heading-rank; base=h2

\section{第九章。所谓与一位蒙丹派姐妹沟通的细节。}

当我们断言灵魂拥有一种独特性质和种类的形体时，在这种特殊状态中，我们将已经获得关于其形体性的所有其他偶性的决定：它们如何属于它，因为我们已证明它是一个形体，但就连它们也有一种与它们所属形体的特殊本性相称的独特性质；否则，如果（形体的）某些偶性在此实例中因其缺失而引人注目，那么这也是灵魂形体性状态的独特性所致，它缺失了所有其他有形存在者所具备的种种性质。然而，即便如此，我们若宣称形体更常见的特征——它们总是归于形体状态——也属于灵魂，如形状与限制，以及那三度——我指长、宽、高——哲学家用以度量所有形体，也毫不矛盾。现在剩下的不就是给灵魂一个形状吗？\Person{柏拉图}拒绝这样做，仿佛这会危及灵魂的不朽。因为根据他，一切有形状的东西都是复合的，由部分构成；而灵魂是不朽的；不朽的，因此是不可分解的；不可分解的，因此是无形状的：因为相反，如果它有形状，它将是复合的、有结构的形成物。他却在另一种方式下为灵魂塑造了一个理智形式的肖像，因其公正的对称与哲学的训导而美丽，却被某些相反的性质扭曲。至于我们，我们确实将形体性的轮廓刻印在灵魂上，不只是基于理性对其形体本性的肯定，而且也基于天主的恩宠藉启示刻印在我们心中的坚定信念。因为，既然我们承认属神的\Emph{\OriginalTerm{神恩}{charismata}}，或恩赐，我们也配得到了先知之恩，虽然是在\Person{若翰}（洗者）之后。现在我们中间有一位姊妹，她有幸蒙受各种启示之恩，她在教会主日的神圣礼仪中，在圣神内通过神魂超拔的经历体验这些恩赐：她与天使交谈，有时甚至与主交谈；她看见并听见神秘的讯息；她了解一些人的心，并向有需要的人分施疗愈。无论是在读经、咏唱圣咏、讲道，还是献上祈祷中，所有这些宗教礼仪都为她提供了看见异象的题材和机会。或许发生在我们身上，当我们的这位姊妹在圣神中神魂超拔时，我们曾以某种不可言喻的方式谈论灵魂。在神圣礼仪结束众人散去之后，她习惯向我们报告她在异象中所见的一切（因为她的一切讯息都经过最严谨的查证，以检验其真实性）。\Quote{除了其他事，}她说，\Quote{有一个具有形体形状的灵魂向我显示，一个灵常向我显现；并非虚无空洞的幻象，而是甚至可伸手抓握的，柔软、透明，带着以太颜色，形状在各方面皆肖似人。}这就是她的异象，为她作证的是天主；宗徒确实预言过，教会中将有\Quote{属神的恩赐}。现在，纵然每一项确凿的证据都摆在眼前令你信服，你还能拒绝相信吗？既然灵魂是有形质体，无疑它拥有我们刚才提到的那些性质，其中就有\Emph{颜色}的特性。除了以太般的透明颜色，你还能将什么颜色归于灵魂呢？并非说它的实体真的是以太或空气（尽管这是\Person{埃奈西德摩}和\Person{阿那克西美尼}的看法，我想\Person{赫拉克利特}也是如此，有人这样说），也非透明的光（尽管\Person{本都的赫拉克利德}认为它是光）。\Quote{雷石}确实并非火质，因为它们闪耀红彤之色；绿柱石也非水所构成，因为它们呈现纯白波浪之色。此外，还有多少事物因颜色而被归入同一类，但自然却将它们远远分开！然而，由于一切极为稀薄透明之物都与空气极为相似，灵魂也是如此，因为在其物质本性上，它是风和气息（或神）；因此，由于其本质的极端稀薄与微妙，对其形体性质的信仰受到威胁。同样，关于人类灵魂的形状，你凭自己的概念可以想见，它不外乎人形；实际上，不外乎每个灵魂所赋予生命并使之活动的那个身体的形状。我们默想人的原初形成便可立即被引向承认这一点。因为只须仔细思量：在天主向人的面孔嘘入生命的气息之后，人便成了一个有灵的生物，那气息必然立即穿过面孔进入内部结构，并散布到身体的各个空间；一旦藉着天主的灵感而凝结，它必定将自身印在每个内部特征上——这些特征已为凝结所充满——因而可以说形状被固定（或定型）。因此，通过这种致密化过程，灵魂的形体性得以固定；通过印模，其形状得以形成和塑造。这就是内在的人，与外在的人不同，却在双重状态中为一。它也有自己的眼睛和耳朵，\Person{保禄}必是藉此听见并看见主；\BibleRef{格后 12:2}此外，它还有身体的其他所有肢体，藉助它们，它实现思考的一切过程和梦中的一切活动。因此，地狱中的富人有一条舌头，穷人（\Person{拉匝禄}）有一根手指，\Person{亚巴郎}有一个怀抱。\BibleRef{路 16:23}祭台下的殉道者灵魂也通过这些特征被识别和认识。事实上，起初与\Person{亚当}身体结合的灵魂，随着身体的成长而成长，并按其形状塑造，结果证明它既是整个（人类灵魂）实体的胚芽，也是那（部分）受造物的胚芽。\par

% structure: p0023 source=h2 target=section reason=relative-heading-rank; base=h2

\section{第10章 灵魂的单纯本性随柏拉图而被确认。精神与灵魂的同一性。}

坚定信仰的必要条件是随同\Person{柏拉图}宣认灵魂是单纯的；换言之，是一致而不复合的；也就是说，就其本体而言是单纯的。不必理会人为的见解和理论，让异端的捏造见鬼去吧！有些人主张，灵魂内有一种与它不同的自然本体——精神；仿佛拥有生命——灵魂的功能——是一回事，而呼气——精神所谓的功能——是另一回事。但并非所有动物都同时具有这两种功能；因为有许多动物只活而不呼吸，因为它们不具备呼吸器官——肺和气管。但在考察人的灵魂时，从蚊虫或蚂蚁借用证据又有何用？伟大的造物主在其神圣安排中，已赋予每一种动物适合其本性与气质的生命器官，因此我们不应从这类比较中汲取任何猜测。人虽然器官完备地拥有肺和气管，但不会因此被证明是用一种过程\Emph{呼吸}，用另一种过程\Emph{活}；蚂蚁虽然缺乏这些器官，但不会因此被认为没有呼吸，仿佛它只是活着而已。因为究竟有谁真正获得了对天主作为如此清晰的洞察，以致有权假定任何生物都缺乏这些器官资源？那位著名的外科医生，或我几乎可以称之为屠夫的\Person{赫罗菲路斯}，为了探究自然的秘密，解剖了无数人，无情地处理人类的受造物以发现（其形体和构造）：我怀疑他是否成功清晰地探查了所有内部结构，因为死亡本身会改变并扰乱生命的自然功能，尤其是当死亡非自然发生，而必定在解剖过程中引起不规则和错误时。哲学家们断言，蚊虫、蚂蚁和飞蛾没有肺或动脉器官是一个确定的事实。那么请问，你这好奇而精密的奥秘研究者，它们有用于观看的眼睛吗？然而它们却可以到达任何想去的地方，并通过视觉过程躲避和瞄准各种对象：请给我指出它们的眼睛，让我看看它们的瞳孔。飞蛾也啃咬和进食：向我展示它们的下颌，揭示它们的臼齿。再者，蚊虫嗡嗡作响，即使在黑暗中也能找到通往我们耳朵的路：那么就向我指出，不仅是那发出噪音的管子，还有它们口中那刺人的长矛。随便取一种生物，哪怕是最微小的，它也必须以某种食物来喂养和维持：那么请向我展示它们摄取、消化和排泄食物的器官。因此，我们该说什么呢？如果生命是靠这些工具维持的，那么这些工具当然必须存在于所有将要生活的物体中，即使它们因微小而不为眼睛或理解力所察觉。你更容易相信这一点，如果你记得天主在小物件中与在最大物件中同样彰显祂的创造伟大。然而，如果你认为天主的智慧没有能力形成如此微小的粒子，你仍可以认识到祂的伟大，因为祂甚至赋予了最小动物生命的功能，尽管缺少合适的器官——使它们即使没有眼睛也能拥有视力，没有牙齿也能进食，没有胃也能消化。有些动物还能在没有脚的情况下前进，如蛇的滑行运动，或蠕虫的垂直努力，或蜗牛和蛞蝓的粘滑爬行。那么，你为何不相信同样可以在没有肺的风箱和动脉管道的情况下进行呼吸呢？这样你就会为自己提供一个强有力的证据，证明精神或呼吸是人的灵魂的附属品，正是因为有些生物缺乏呼吸，而它们缺乏呼吸是因为不具备呼吸器官。你认为一个东西没有呼吸也能活着；那么为什么不设想一个东西没有肺也能呼吸呢？请问，呼吸意味着什么？我想，意思是从你自身呼出气息。不活着意味着什么？我想，意思是不从你自身呼出气息。如果\Quote{呼吸}跟\Quote{活着}不是一回事，我就必须这样回答。然而，不呼吸一定是死人的特征：因此呼吸是活人的特征。但呼吸也是呼吸者的特征：所以呼吸也是活人的特征。现在，如果两者都可能在没有灵魂的情况下完成，那么呼吸可能不是灵魂的功能，而仅仅是活着。但事实上，活着就是呼吸，呼吸就是活着。因此，整个呼吸和活着的过程都属于那个活着所属于的东西——即灵魂。那么，既然你把精神（或气息）和灵魂分开，也把它们的运作分开吧。让它们各自彼此分开完成某种行为——灵魂分开，精神分开。让灵魂在没有精神的情况下活着；让精神在没有灵魂的情况下呼吸。让其中一个离开人的身体，让另一个留下；让死亡和生命相遇并调和。如果灵魂和精神确实是两个，它们就可以被分开；这样，通过离开的那个与留下的那个的分离，就会产生生命和死亡的结合与相遇。但这种结合永远不会产生：因此它们不是两个，也不能被分开；但如果它们是两个，本来是可以被分开的。诚然，两个事物在生长中肯定会合并。但这里的两个永远不会合并，因为活着是一回事，呼吸是另一回事。实体由它们的运作来区分。既然你没有给灵魂和精神赋予任何区别，那么你相信灵魂和精神是一体的理由就更充分了；这样，灵魂本身就是精神，呼吸是那个被赋予生命的东西的功能！但如果你坚持认为白天是一回事，而日光是白天偶然的属性是另一回事，而白天仅仅是光本身呢？当然，也有不同种类的光，正如从火的使用中（所显现的）那样。同样，也有不同种类的精神，根据它们来自天主还是魔鬼。事实上，每当问题涉及灵魂和精神时，灵魂将被（理解为）本身就是精神，正如白天本身就是光一样。因为一个事物与它藉以存在的东西是同一的。\par

% structure: p0025 source=h2 target=section reason=relative-heading-rank; base=h2

\section{第十一章 神灵——指灵魂行动而非本性的术语。须与天主之神仔细区分。}

但我目前研究的性质迫使我将灵魂称为神灵或气息，因为呼吸被归属于另一种实体。然而，我们却将这一（行动）归于灵魂，我们承认灵魂是不可分的单纯实体，因此我们必须以确定的意义称它为神灵——不是因为它的状态，而是因为它的行动；不是就其本性而言，而是就其运作而言；因为它\Emph{呼吸}，而不是因为它在任何特殊意义上是神灵。因为吹气或呼吸就是呼吸。所以我们不得不通过（表示这一呼吸的术语，即）\Emph{神灵}来描述灵魂，我们因灵魂行动的特性而认为它即是气息。此外，我们正确且特别地坚持称它为气息（或神灵），以反对\Person{赫尔莫根}，他主张灵魂来自物质，而非来自天主的\Emph{吹气}或气息。他确实公然反对圣经的证言，并为了这一观点将气息转变为神灵，因为他无法相信（被吹入）天主之神的气息的受造物会陷入罪恶，然后陷入定罪；因此他断定灵魂来自物质，而非来自天主之神或气息。为此，我们这一方甚至从那段经文起，就坚持灵魂是气息而非神灵，这是就圣经中独特的神灵意义而言；在此，我们遗憾地以较低的意义使用神灵一词，因为呼吸和吹气的行动相同。在那段经文中，唯一的问题关乎自然实体；呼吸本是一个本性的行动。如果不是为了那些将某种我不理解的神性胚芽引入灵魂的异端者，我本不会在此停留片刻：（他们声称）这胚芽是由灵魂之母\OriginalTerm{智慧}{Sophia}（\Emph{智慧}）在造物主不知情的情况下秘密慷慨地赐予灵魂的。但（神圣的）圣经更了解灵魂的创造者，更确切地说就是天主，它只告诉我们天主向人的脸上吹入生命的气息，人就成了有生命的灵魂，借此他得以生活和呼吸；同时，圣经在以下这类经文中对神灵与灵魂作了足够清楚的区分，天主亲自宣告：\Quote{我的神从我发出，我造了各人的气息。我神的气息成了灵魂。} 又说：\BibleQuote{祂将气息赐给地上的人民，将神灵赐给行走其上的人。} \BibleRef{依 42:5} 首先到来的是（本性的）灵魂，即气息，给地上的人民——换句话说，给那些在肉身内按本性行事的人；然后后来神灵降临给行走其上的人——即那些制服肉身工作的人；因为宗徒也说：\BibleQuote{那不是首先属神的，而是首先属本性的（即拥有本性灵魂的），然后才是属神的。} \BibleRef{格前 15:46} 因为，亚当随即预言了\BibleQuote{基督与教会的大奥秘，} \BibleRef{弗 5:31}，当他说：\BibleQuote{这才是我骨中的骨，肉中的肉；因此人要离开父亲和母亲，依附他的妻子，二人成为一体} \BibleRef{创 2:24}，他经历了神灵的影响。因为那时\OriginalTerm{神魂超拔}{ecstasy}降临在他身上，那是圣神运作的预言德能。甚至恶神也是一种降临于人身上的影响力。的确，天主之神并不比后来恶神将他转变为另一个人——即一个背教者——更为真实地\BibleQuote{将撒乌耳转变为另一个人，} \BibleRef{撒上 10:6}，即变为一位先知，当时\BibleQuote{人们彼此说：克士的儿子遇到了什么？撒乌耳也列在先知中吗？} \BibleRef{撒上 10:11} 犹达斯同样被长期算在选民（宗徒）之列，甚至被任命为他们的司库；他尚未成为叛徒，尽管已变得欺诈；但后来魔鬼进入了他。因此，既然天主之神与魔鬼之神都不是与人的灵魂在出生时自然栽种的，那么灵魂显然必须独立且单独存在，在二者之一加入之前；如果这样独立且单独，它在其本体上也必然是单纯而非复合的；因此它除了来自其自身本体的实际状态外，不可能从其他原因呼吸。\par

% structure: p0027 source=h2 target=section reason=relative-heading-rank; base=h2

\section{第十二章 理智与灵魂的区别及其相互关系}

同理，心智（或 animus，希腊人称之为 \GreekText{ΝΟΥΣ}）在我们看来，也仅仅是指那内在固植于灵魂、天然属于灵魂的功能或官能，灵魂藉此运作、获取知识，并因拥有它而能在自身内自发运动，从而看似被心智所推动，仿佛它是另一种实体一般——正如那些认为灵魂是宇宙运动本原者所主张的那样：即苏格拉底之神、\Person{瓦伦提诺}所说的其父 \OriginalTerm{拜托斯}{Bythus} 的\Quote{独生子}及其母 \OriginalTerm{西格}{Sige}。\Person{阿那克萨戈拉}的观点多么混乱！因为他想象心智是万物的始基，宇宙的天平以它为轴心；并且断言心智是单一、纯粹、不可混合的本原，正是基于这一点，他主要将心智与灵魂的一切混合分开；然而在另一处，他又实际将心智与灵魂结合起来。\Person{亚里士多德}也注意到了这种（矛盾）；但他究竟是想作建设性的批评，以充实自己的体系，而非破坏他人的原理，我实在难以断定。至于他本人，虽然推迟了对心智的定义，但他一开始就提到心智的两种自然成分之一，即他猜想为不动情或不受情感影响的神圣本原，从而将心智从与灵魂的一切关联中移除。因为既然灵魂显然会承受那些它自然应受的情感，那么它必然要么由心智承受，要么与心智一同承受。如果灵魂在本性上与心智相连，就不能推出心智是不动情的；反之，如果灵魂不由心智或与心智承受，那么它就不可能在本性上与心智相连，因为与它一同承受的心智本身什么也不承受。此外，如果灵魂不由心智或与心智承受任何东西，它就不会有任何感觉，不会获得任何知识，也不会通过心智经历任何情感，正如他们主张的那样。因为\Person{亚里士多德}甚至将感官视为情感或情绪状态。这很正确。因为运用感官就是承受情感，因为承受就是感觉。同样，获取知识就是运用感官；经历情感就是运用感官；这一切都是承受状态。但我们看到，灵魂在这些方面并没有以这样一种方式经历任何东西，以至于心智也被情绪所影响，尽管一切都是由情绪、并与情绪一起完成的。因此，心智是可以混合的，与\Person{阿那克萨戈拉}相反；并且是可动情或易受情感的，与\Person{亚里士多德}的意见相反。此外，如果承认灵魂与心智是分离的状态，以至于它们在实体上是两个东西，那么其中之一将具有情感、感觉、一切品味、所有行动和运动作为特征；而另一个的自然状态将是宁静、安息和麻木。因此别无选择：要么心智无用且空虚，要么灵魂无用且空虚。但如果这些情感肯定都可以归于两者，那么两者就是一回事，\Person{德谟克利特}就会达成他的主张，即消除两者之间的一切区别。问题将是两者如何成为一——是通过两种实体的混淆，还是通过一种实体的安排？然而，我们断言心智与灵魂结合——并不是在实体上与之分离，而是作为它的自然功能和活动。\par

% structure: p0029 source=h2 target=section reason=relative-heading-rank; base=h2

\section{第十三章 灵魂的至上性}

接下来需要考察至上性何在；换言之，两者之中哪一个优于另一个，这样，至上性显然所在的那一个便是本质上更优的实体；而那被本质上更优的实体所统辖的，则应被视为更优实体的自然功能执行者。谁还会犹豫将这全部权威归于灵魂呢？因为从灵魂的名称，整个人在普通用语中得到了自己的称呼。富人说：我养活多少\Emph{灵魂}？而不是多少\Emph{心智}。船夫的愿望也是从船难中拯救那么多\Emph{灵魂}，而不是那么多心智；劳工在劳作中，士兵在战场上，都说他交出他的灵魂（或生命），而不是他的心智。两者之中，哪一个的危险、许愿或愿望更常挂在人们嘴边——心智还是灵魂？哪一个是临终者据说要面对的——心智还是灵魂？简而言之，哲学家们和医学家们，即使他们有意讨论心智，也总是在标题和目录上写着：\Quote{\WorkTitle{De Anima}}（\Quote{\WorkTitle{论灵魂}}）。而且为了你也拥有天主在这事上的见证，祂所说话的对象是灵魂；祂劝诫和嘱咐的也是灵魂，叫它将心智和理智转向祂。基督来拯救的是灵魂；祂威胁要在地狱中毁灭的也是灵魂；祂禁止我们过分珍视的是灵魂（或生命）；好牧人自己为羊舍去的也是祂的灵魂（或生命）。因此，你当将至上性归于灵魂；在\Emph{它}里面，你也拥有那实体的统一，你看出心智是这统一体的工具，而非主宰力量。\par

% structure: p0031 source=h2 target=section reason=relative-heading-rank; base=h2

\section{第十四章 哲学家对灵魂的不同划分；这种划分并非物质性的解剖}

灵魂既如此单一、单纯、自足，便既不能由外在元素构成和组合，也不能在自身内分裂，因为它是不分解的。因为若灵魂能被构造和毁灭，它便不再是永生的了。然而，既然它并非有死的，它也就不能被分解和分裂。分裂意味着分解，分解意味着死亡。然而（哲学家们）却将灵魂分为部分：例如，柏拉图分为两部分；芝诺分为三部分；帕奈提乌斯分为五或六部分；索拉努斯分为七部分；克吕西普分为多达八部分；阿波罗芬分为多达九部分；而某些斯多亚哲学家则在灵魂中发现了多达十二部分。波西多尼乌斯甚至在此之上又加了两部分：他从灵魂的两个主导官能开始——他们称为\OriginalTerm{主导官能}{\GreekText{ἡγεμονικόν}}的\Emph{指导}官能，以及他们称为\OriginalTerm{理性官能}{\GreekText{λογικόν}}的\Emph{理性}官能——最终又将它们细分为十七部分。各家学派就这样以各种方式解剖灵魂。然而，这些划分不应被视为灵魂的部分，而应视作其能力、官能或运作——甚至亚里士多德本人也曾视其中一些为如此。因为它们不是灵魂实质的片段或有机部分，而是灵魂的功能——如运动、行动、思想的功能，以及他们以这种方式划分的其他功能；同样，还有众所周知的五种感官——视、听、尝、触、嗅。虽然他们分别为这些感官分配了身体中的特定部位作为其特殊居所，但这并不意味着类似的分配也适用于灵魂的各部分；因为即使身体本身也不容许像他们想让灵魂经受的那种划分。而是由全部肢体构成一个身体，所以这种安排更是一种凝聚而非分裂。请看阿基米德那件非常奇妙的机械装置——我指的是他的水力风琴，其众多部件、零件、带子、音孔、出音口、和声构造以及管道的排列；然而所有这些细节只构成一件乐器。同样，在水力引擎的推动下吹遍这架风琴的风，并不因其散布于乐器中使其运作而被分成若干部分：它在实质上仍是完整而整全的，尽管在运作上是分散的。这个例子与斯特拉托、埃奈西德穆和赫拉克利特的（比喻）相差不远：因为这些哲学家主张灵魂是统一的，散布于整个身体，却在每一部分都是相同的。恰如吹过风琴各管的风，灵魂通过感官以多种方式展现其能力，并非被分割，而是按自然秩序分布。至于这些能力应以何种名称来认识，应依据它们自身的何种划分来分类，以及它们在身体中应分别局限于哪些特殊职务和功能，应由医生和哲学家去考虑和决定：对我们而言，只作几点简短的评论便已足够。\par

% structure: p0033 source=h2 target=section reason=relative-heading-rank; base=h2

\section{第十五章 灵魂的活力与理智。其在人身上的特性与居所}

首先，（我们必须确定）灵魂中是否存在某种至高的生机与理智的原则，他们称之为\Quote{灵魂的统治力}——\OriginalTerm{灵魂的统治力}{\GreekText{τὸ} \GreekText{ἡγεμονικόν}}。因为如果不承认这一点，灵魂的整个状况就岌岌可危。的确，那些说没有这种指导能力的人，从一开始就假定灵魂本身只是一个虚无。一位麦西尼亚人\Person{狄凯亚尔库斯}，以及在医学界中的\Person{安德烈亚斯}和\Person{阿斯克勒皮亚德斯}，就这样摧毁了（灵魂的）指导能力，因为他们实际上将感官置于理智中，并为感官主张统治能力。\Person{阿斯克勒皮亚德斯}甚至用这样的论据粗暴地对待我们：许多动物在失去其身体中一般认为灵魂的生机与感觉原则主要寓居的部分之后，仍然在很大程度上保留生命和感觉；例如苍蝇、黄蜂和蝗虫，当砍掉它们的头时；以及母山羊、乌龟和鳗鱼，当取出它们的心脏时。因此（他得出结论），灵魂没有特殊的原理或能力；因为如果有的话，当灵魂随其居所（或身体器官）被移除时，它的活力和力量就无法持续。然而，\Person{狄凯亚尔库斯}有几位权威人士反对他——也是哲学家——\Person{柏拉图}、\Person{斯特拉托}、\Person{伊壁鸠鲁}、\Person{德谟克利特}、\Person{恩培多克勒}、\Person{苏格拉底}、\Person{亚里士多德}；而反对\Person{安德烈亚斯}和\Person{阿斯克勒皮亚德斯}（可以列出他们的同行）的医生有\Person{希罗菲卢斯}、\Person{埃拉西斯特拉图斯}、\Person{狄奥克莱斯}、\Person{希波克拉底}，以及\Person{索拉努斯}本人；而且胜过所有其他人，还有我们基督宗教的权威。关于这两个问题，我们由天主教导：即灵魂中有一个统治能力，并且它被供奉在身体的一个特定隐密处。因为，当人读到天主是\BibleQuote{心灵的探索者与见证者；}\BibleRef{智 1:6}；当祂的先知因祂向他揭露心灵隐秘而受责时\BibleRef{箴 24:12}；当天主自己在祂子民中预知他们心中的思念说：\BibleQuote{你们为什么心里思念恶事？}\BibleRef{玛 9:4}；当\Person{达味}祈祷说\BibleQuote{天主，求祢在我内造一颗纯洁的心}，\Person{保禄}宣告说\BibleQuote{心里相信，就可成义，}\BibleRef{罗 10:10}；\Person{若望}说\BibleQuote{各人在自己心里被定罪；}\BibleRef{若一 3:20}；最后，\BibleQuote{凡注视妇女，有意贪恋她的，他已在心里与她犯了奸淫，}\BibleRef{玛 5:28}——那时，这两点就完全清楚了：灵魂有一个指导官能，天主的旨意可与它相符；换句话说，理智与生机的一个至高原则（因为哪里有理智，哪里就必有生机），并且它寄居在我们身体中最宝贵的那部分，天主特别地看顾它。因此，你不要与\Person{赫拉克利特}一样，以为我们正在讨论的这个至高官能是被某种外力驱动的；也不要像\Person{莫斯奇翁}那样，以为它飘浮于全身；也不要像\Person{柏拉图}那样，以为它封闭在头部；也不要像\Person{克塞诺芬}那样，以为它集中于头顶；也不要如\Person{希波克拉底}所认为的那样，安息于大脑中；也不要如\Person{希罗菲卢斯}所想，在大脑底部周围；也不要如\Person{斯特拉托}和\Person{埃拉西斯特拉图斯}所说，在脑膜中；也不要如医生\Person{斯特拉托}所持，在双眉之间；也不要如\Person{伊壁鸠鲁}所说，在胸腔的围护内；而是如埃及人一直所教导的，尤其是他们中被视为神圣真理诠释者的人所教导的；同时也符合\Person{俄耳甫斯}或\Person{恩培多克勒}的诗句：\par

\Quote{因为人的感知在于心脏周围的血液。}\par

\Quote{人的（最高）感知在于他心脏周围的血液。}\par

同样地，\Person{普罗泰戈拉}、\Person{阿波罗多罗斯}和\Person{克吕西波}也持有相同的观点，因此（我们的朋友）\Person{阿斯克勒皮亚德斯}可以去寻找他的无心脏而咩咩叫的山羊，去追猎他的无头苍蝇；也让所有那些从野兽的状态预先断定人类灵魂特性的人（这些贤哲们）确信，反而是他们自己活在一种无心无脑的状态。\par

% structure: p0038 source=h2 target=section reason=relative-heading-rank; base=h2

\section{第十六章 灵魂的各个部分。理性灵魂的要素。}

\Person{柏拉图}的这个观点也与信仰相当一致，他将灵魂分为两部分：理性的和非理性的。对于这一定义，我们并无异议，但我们不会将这双重的区分归诸（灵魂的）本性。我们必须相信理性的元素是灵魂的自然状态，从其最初被创造时就由它的创造主印刻在它上面，而创造主自己本质上就是理性的。因为，天主出于祂自己的推动所产生的，又何况祂明明借着祂自己的\Emph{气息}或呼吸所发出的，怎么会不是理性的呢？然而，非理性的元素，我们必须理解为后来才产生的，源自蛇的诱惑——即（原初）过犯的成就——从此成为灵魂的内在，与灵魂一同成长，此时已经呈现出一种自然发展的方式，因为它确实是在本性的开端就发生了。但是，既然同一位\Person{柏拉图}认为只有理性的元素存在于天主自己的灵魂中，如果我们也将非理性的元素归与我们灵魂从天主所接受的本性，那么非理性的元素也将同样源自天主，作为自然的产物，因为天主是自然之作者。而从魔鬼则生出罪的诱因。然而，所有罪都是非理性的：因此非理性的来自魔鬼，罪也来自他；它对于天主是外在的，非理性的原则对于天主也是异己的。这两个元素的多样性来自于它们作者的差异。因此，当\Person{柏拉图}将（灵魂的）理性元素保留给天主独有，并将它再分为两部分：\Emph{易怒的}（他们称之为\OriginalTerm{易怒的部分}{\GreekText{θυμικόν}}）和\Emph{贪欲的}（他们用\OriginalTerm{贪欲的部分}{\GreekText{ἐπιθυμητικόν}}来称呼），这样的分法使得前者为我们和狮子所共有，后者为人类与苍蝇所共有，而理性元素仅限于我们和天主——我看这一点将必须由我们来处理，因为我们发现在基督内也有这些官能运作。因为你们可以在主身上看到这三种性质。祂有\Emph{理性的}元素，借此教导、讲论、预备救恩之路；此外祂还有\Emph{义怒}，借此斥责经师和法利塞人；还有\Emph{渴望}的原则，借此祂如此切望同门徒一起过逾越节。\BibleRef{路 22:15}因此，在我们的情况中，我们灵魂的易怒和贪欲的元素并不总是要归到非理性的（本性）上，因为我们确定在我们主身上这些元素是完全按照理性运作的。天主会完全有理由向所有配得祂愤怒的人发怒；并且天主也会合理地渴望那些与祂相称的对象和要求。因为祂会对恶人显示义怒，对善人则渴望救恩。连宗徒也容许我们有贪欲的倾向。他说：\BibleQuote{若有人想要主教的职分，就是羡慕善工。}\BibleRef{弟前 3:1}现在，通过说\Quote{善工}，他告诉我们这渴望是合理的。他也容许我们有所义怒。他自己也经历同样的义怒，怎会不容许呢？他说：\BibleQuote{我恨不得那搅乱你们的人把自己割绝了。}\BibleRef{迦 5:12}那义怒完全与理性一致，是由他渴望保持纪律和秩序而产生的。然而，当他说：\BibleQuote{我们从前也都在怒气中生活，}\BibleRef{弗 2:3}他谴责了一种非理性的易怒，这种易怒并非出自天主所生的本性，而是出自魔鬼所引进的，魔鬼自己被称作其阶级的\Quote{主人}，\BibleQuote{你们不能事奉两个\Emph{主人}，}\BibleRef{玛 6:24}并且有\Quote{父亲}的实际称呼：\BibleQuote{你们是出于你们的\Emph{父亲}魔鬼。}\BibleRef{若 6:44}因此，当你读到他是\Quote{撒莠子的}，是夜间毁坏庄稼的祸害时，\BibleRef{玛 13:25}你不必害怕将那第二个、后来的、堕落的本性（我们一直谈到的）的主权和统治归于他。\par

% structure: p0040 source=h2 target=section reason=relative-heading-rank; base=h2

\section{第十七章：感觉的忠实性——\Person{柏拉图}质疑，\Person{基督}亲自辩护}

随后，当我们遇到这样一个问题（关于我们与识字一同学会的那五种感官的真实性；因为连从这里也产生对我们异端者的一些支持）时。这些官能是视、听、嗅、味、触。柏拉图主义者们对这些感官的可靠性提出了过于严厉的指责，据某些人说，\Person{赫拉克利特}、\Person{狄奥克勒斯}和\Person{恩培多克勒}也是如此；无论如何，\Person{柏拉图}在\WorkTitle{蒂迈欧篇}中宣称感官的活动是非理性的，并受我们的意见或信念所败坏。人们归罪于视觉，因为它断言浸在水中的桨是倾斜或弯曲的，尽管确定它们本是直的；因为，它又非常确信那座远处实际是四方轮廓的塔楼是圆的；因为它还会认为那座廊子或拱廊的构造本是真正平行的，却设想它越到尽头越窄；并且它会将高悬于大海之上的天空与大海连接起来。同样，我们的听觉也被指控有谬误：例如，我们以为那是空中的响声，而它不过是马车的隆隆声；或者，如果你更喜欢另一种方式，当雷声在远处轰鸣时，我们很确信那是一辆马车发出的声响。我们的嗅觉和味觉也有缺陷，因为同样的香水和葡萄酒在我们使用一段时间后就失去了价值。基于相同的原理，我们的触觉也受到指责，因为同一块路面，手摸起来粗糙，脚感却足够平滑；而在浴场，一股温水起初被认为很热，后来却又觉得温度适宜。因此，按照他们的说法，我们的感官欺骗了我们，而其实是\Emph{我们}自己（通过改变意见）造成了这些不一致。斯多葛学派的态度更为温和，因为他们并不每个感官、所有时候都背负欺骗的恶名。伊壁鸠鲁学派则表现出更大的一致性，他们坚持所有感官在其见证中同样真实，且始终如此——只是方式不同。并非我们的感觉器官有错，而是我们的意见。感官只经历感觉，并不形成意见；是灵魂在\Emph{意见}。他们将意见与感官分开，将感觉与灵魂分开。然而，意见若不来自感官，又从何而来？除非眼睛没看到那座塔的圆形形状，它不可能有圆形的概念。同样，感觉若不来自灵魂，又从何而来？因为如果灵魂没有身体，就不会有感觉。因此，感觉来自灵魂，意见来自感觉；整个过程都是灵魂。但进一步，可以合理地坚持说，存在某种东西造成了感官报告与事实真相之间的不一致。既然（如我们所见）可能报道那些并不存在于对象中的现象，为什么就不可能同样地报道那些并非由感官引起，而是由介入的事由和条件（本质上出于情况本身）所产生的现象呢？如果是这样，那么恰当的做法就是予以正确认识。真相是，水是导致桨看起来倾斜或弯曲的原因；在水外，它从外观上（以及实际上）是完全直的。那因光泽而形成镜子的物质或介质，根据它受到撞击或摇晃的方式，其振动实际上破坏了一条直线的直度外观。同样，填满我们与塔楼之间距离的开放空间的状态，必然使我们无法注意到塔楼的真实形状；因为周围空气均匀的密度用相似的光覆盖了塔楼的棱角，抹去了它们的轮廓。同样，拱廊的等宽在向末端时变得尖锐或变窄，直到其外观在延长的屋顶下越来越收缩，在朝向最远距离的方向上趋于消失点。天空与大海混合在一起，视觉最终耗尽，它曾适当维持着两种元素的界限，只要其敏锐的目光持续。至于（所谓的）听觉错觉，除了声音的相似性，还能产生什么幻觉呢？如果后来香水对嗅觉不那么强烈，葡萄酒对味觉更平淡，水对触觉不那么热，那么它们最初的力量毕竟在整体上仍然相当完好无损地存在着。然而，关于路面的粗糙与平滑，手脚这样的肢体在柔软与坚硬上如此不同，产生不同印象只是自然而合理的。这样看来，我们的感官不可能在没有充分原因的情况下产生错觉。那么，如果特殊原因（如我们所指出的）误导我们的感官和（通过感官）我们的意见，那么我们就不应再将欺骗归咎于那些遵循错觉特定原因的感官，也不应归咎于我们形成的意见；因为这些意见由我们的感官引发和控制，而感官只是遵循那些原因。患有疯狂或精神错乱的人会将一物误认为另一物。\Person{俄瑞斯忒斯}在他的妹妹身上看到了他的母亲；\Person{埃阿斯}在屠杀的牛群中看到了\Person{尤利西斯}；\Person{阿塔玛斯}和\Person{阿高厄}在他们的孩子身上看到了野兽。对于如此巨大的谬误，你是该责怪他们的眼睛还是他们的疯狂？所有东西对于患黄疸病的人来说，由于胆汁过多，都尝起来发苦。你会将这种生理上的偏离归咎于他们的味觉，还是他们的病态健康？因此，所有感官有时都会紊乱或受骗，但只是以这样一种方式，即它们自己的自然功能完全无过。但进一步说，我们甚至不该把欺骗的指控加给这些特殊的原因和条件本身。因为，既然这些生理偏差出于确定的原因，这些原因就不应被视为欺骗。凡应当以某种方式发生的事就不是欺骗。那么，如果连这些附带的原因也必须免于一切责难和指责，我们岂不更应使感官免于非难——这些原因对感官施加着自由的影响？因此，我们必定十分确定地要求感官具有真理、忠实和完整，因为它们对其印象所给出的报告，除了在一切情况下产生那种出现在感官报告与对象实在之间的不一致的那些特殊原因或条件所命令者外，从未给出过其他。那么，最狂妄的学园啊，你是什么意思？你颠覆了人类生活的整个状况；你扰乱了整个自然秩序；你遮蔽了天主自己的美好眷顾：因为天主指派给祂一切工程之上的人的感官——我们藉以理解、居住、管理并享受它们——（你却指责）是谬误和背信的暴君！难道不是藉着它们，一切受造物才接受我们的服务？难道不是藉着它们，第二个形式甚至印在世界之上？——如此多的技艺，如此多的勤勉资源，如此多的事业，如此多的生意，如此多的职务，如此多的商业，如此多的疗法、商议、安慰、方式、文明和生活的成就！所有这些东西产生了人类存在的真正趣味和滋味；而藉着人的这些感官，他在所有有生命的自然中唯独享有理性动物的特征，具有智力和知识的能力——甚至有能力形成学园本身！但\Person{柏拉图}，为了贬低感官的见证，在\WorkTitle{费德鲁斯篇}中（借\Person{苏格拉底}之口）否认他自己有能力认识自己，尽管有\Place{德尔斐}神谕的诫命；在\WorkTitle{泰阿泰德篇}中他剥夺了自己知识和感觉的官能；又在\WorkTitle{费德鲁斯篇}中他将对真理的所谓死后知识推迟到死后；然而他甚至在死前仍继续做哲学家。我说，我们不可、不可质疑（可怜的、被诋毁的）感官的真理，以免我们甚至在基督自己身上，也对其感觉的真理产生怀疑；以免人们说祂并没有\BibleQuote{看见撒殚如同闪电从天坠落；}\BibleRef{路 10:18}祂并没有真的听见圣父为祂作证的声音；\BibleRef{玛 3:17}或者说祂在抚摸\Person{伯多禄}的岳母时受了欺骗；\BibleRef{玛 8:15}或者说祂后来闻到的香膏的芬芳不同于祂接受为埋葬自己的那一种；\BibleRef{玛 26:7}或者那酒的味道不同于祂为纪念自己的血所祝圣的那一种。基于这个错误原则，\Person{玛尔强}实际选择了相信祂是一个幻影，否认祂具有一个完美身体的真实性。然而，甚至连祂的宗徒们也从未被祂的本性所欺骗。祂在山上是真正被看见和听见的；\BibleRef{玛 17:3}在\Place{加里肋亚}（\Place{加纳}）婚宴上的那酒是真实且实在的；\BibleRef{若 2:1}后来相信了的\Person{多默}的触摸也是真实且实在的。\BibleRef{若 20:27}请读\Person{若望}的见证：\BibleQuote{论到那从起初就有的生命之圣言，就是我们所看见的，我们所听见的，我们亲眼所注视过，亲手摸过的。}\BibleRef{若一 1:1}当然，如果我们的眼睛、耳朵和手的见证在本性上是谎言，那么那个见证必是虚假和欺骗性的。\par

% structure: p0042 source=h2 target=section reason=relative-heading-rank; base=h2

\section{第十八章. \Person{柏拉图}向\OriginalTerm{诺斯替派}{Gnostics}提出的某些错误. 灵魂的职能}

现在我要转向理智官能的部分，正如柏拉图将其传授给异端者，它与身体机能相区别，在死亡之前便获得了对它们的认识。他在\WorkTitle{斐多篇}中问道：那么，你认为关于实际拥有\Emph{知识}这件事又当如何？如果有人将身体接纳为寻求知识中的同伴，它会成为障碍还是不会？我也要问一个类似的问题：视觉和听觉的官能对人类来说有任何真理和实在性吗？难道不是甚至连诗人们也总是在私下对我们嘀咕，我们永远无法确切地听到或看到任何事情吗？他无疑记得喜剧诗人厄庇卡尔莫斯曾说过：\Quote{是心灵在看，是心灵在听——其他一切都是又瞎又聋。}他又说了同样大意的话：人是最有智慧的，谁的理智能力最清晰；他从不运用视觉，也不将任何这类官能的帮助加给他的心灵，而是当他在沉思中，为了获得对事物本性的纯正洞见时，以纯粹宁静的方式运用理智本身；他竭尽全力将自己从眼睛、耳朵和（如必须表达的那样）整个身体中分离出来，理由是身体打扰灵魂，每当灵魂与身体相通时，身体不允许灵魂拥有真理或智慧。那么，我们看到，与肉体感觉相对立，另一种更有用的官能被提供了，即灵魂的力量，它产生对那真理的理解，其实在性既不可触摸，也不为肉体感官所感知，而是远离人们的日常知识，隐藏在秘密中——在至高的高处，在天主本身面前。因为柏拉图主张，存在某些不可见的、非物质的、属天的、神圣的、永恒的本体，他们称之为\Quote{理念}，即（原型的）形式，它们是那些对我们显现、处于我们肉体感官之下的自然对象的模型和原因：前者（根据柏拉图）是实际的真理，后者是它们的形象和肖像。那么，这里不是有诺斯替派和瓦伦廷派的异端原则的闪光吗？正是从这种哲学中，他们迫切地采纳了肉体感觉与理智官能之间的区别——他们实际上将这一区别应用于十个童女的比喻：使五个愚拙的童女象征五种肉体感觉，因为它们是如此愚蠢且容易被欺骗；而智慧的童女则表达理智官能的意义，它们如此智慧，以至于能达到那神秘而超验的、位于圆满中的真理。（因此，我们有了）这些异端者理念的神秘原形。因为在这种哲学中，既有他们的永世，也有他们的谱系。同样，他们也把感官分为两种：一种是从其精神种籽来的理智力量，另一种是从属魂的东西来的感觉官能，后者绝对不能理解属灵之事。从前者萌生出不可见之物；从后者萌生出可见之物，是卑贱而暂时的，且为感官所明显，因为它们置于可触摸的形式中。正是由于这些观点，我们在前文中已作为预备事实陈述过：理智无非是灵魂的一个装置或工具，而精神也不是与灵魂分离的另一种官能，而是呼吸中运作的灵魂本身；尽管那天主一方面或魔鬼一方面吹入灵魂的影响，必须被视为一个附加元素。现在，关于理智力量与感觉官能之间的区别，我们只承认自然差异性所要求的程度。（当然，有区别）在形体的事物与精神的事物之间，在可见的与不可见的存在者之间，在对视线显明的事物与隐藏的事物之间；因为一类归于感觉，另一类归于理智。然而，两者都必须被视为内在于灵魂，且服从于灵魂，正如灵魂借由身体把握形体对象，同样地，它借由心灵构想非形体对象，只不过它在运用理智时也同时进行感觉。因为，难道不是这样吗：运用感觉即是运用理智？而运用理智等于运用感觉？感觉本身若不是对所感觉对象的理解，又能是什么呢？而理智或理解本身若不是对所理解对象的看见，又能是什么呢？为什么要采用这种折磨人的方式来折磨简单的知识并钉死真理呢？谁能给我指出一种感觉，它不理解它所感觉的对象，或者一种理智，它不感知它所理解的对象，足以向我证明两者可以彼此分离？如果形体事物是感觉的对象，而非形体事物是理智的对象，那么不同的是对象的\Emph{类别}，而不是感觉和理智的住所或居所；换言之，不是\OriginalTerm{灵魂}{anima}和\OriginalTerm{理智}{animus}。简言之，形体事物是通过什么来感知的？如果是通过灵魂，那么理智就是一种感觉官能，而不仅仅是理智力量；因为当它理解时，它也感知，因为没有感知就没有理解。然而，如果形体事物是通过灵魂来感知的，那么灵魂的力量就是一种理智力量，而不仅仅是感觉官能；因为当它感知时，它也理解，因为没有理解就没有感知。再者，非形体事物是通过什么来理解的？如果是通过理智，那么灵魂在哪里？如果是通过灵魂，那么理智在哪里？因为不同的事物当它们各自履行其职能和职责时，应当彼此缺席。你必然认为，在某些场合理智从灵魂中缺席；因为你假设我们被造得如此，以致不知道我们曾看见或听见某事，基于理智当时缺席的假设。因此我必须坚持：灵魂本身既没有看见也没有听见，因为在给定时刻它与其能动力量——即理智——一同缺席。事实上，每当一个人失去理智时，是他的灵魂错乱了——不是因为理智缺席，而是因为那时它（与灵魂）一同受苦。确实，正是灵魂主要受这类事故的影响。这一事实从哪里得到确认？从以下考虑得到确认：在灵魂离开后，理智不再在一个人身上找到；它总是跟随灵魂；在死后它最终也不会独自留在灵魂后面。既然它跟随灵魂，它也与灵魂不可分离地连接；正如理解与灵魂连接，理智跟随灵魂，理解与理智不可分离地连接。现在，就算理解高于感觉，是奥秘更好的发现者，那又有什么关系呢？只要它只是灵魂的一种特殊官能，正如感觉本身一样？它根本不影响我的论证，除非理解被认为高于感觉，目的是为了从这种更高的声称中也推导出它的分离状态。在驳斥了它们所谓的区别之后，在我接近（异端所提出的）关于一位更高天主的信仰之前，我还必须反驳这个高低的问题。然而，关于（更高）天主这一点，我们必须与异端者在自己地盘上交锋。我们当前的主题关乎灵魂，重点是防止阴险地将优越性归于理智或理解。现在，尽管理智所触及的对象具有更高的性质，因为它们是属灵的，相比感觉所把握的对象，因为后者是属身体的，但它仍然是\Emph{对象}上的优越性——如同崇高的对比于卑微的——而不是\Emph{官能}的理智相对于感觉的优越性。因为当感觉教育理智发现各种真理时，理智怎能优越于感觉呢？事实上，这些真理是通过可触摸的形式学到的；换句话说，不可见之物借助可见之物被发现，正如宗徒在他的书信中告诉我们的：\BibleQuote{因为祂那看不见的事，从创世以来，借着所造之物，就可以明明地看见，被人领悟；}\BibleRef{罗 1:20}并且柏拉图也可能这样告知我们的异端者：\Quote{显现的事物是隐藏之物的形象，}因此必然得出，这个世界无论如何是另一个世界的形象：所以理智显然利用感觉作为自己的引导、权威和支持；没有感觉，真理无法获得。那么，一个事物怎能优越于对其存在起工具作用、对其不可或缺、且其获得的一切都归功于其帮助的事物呢？因此，从我们所说的得出两个结论：（1）理智不应被置于感觉之上，基于（假设的）理由：事物赖以存在的动因低于事物本身；（2）理智不可与感觉分离，因为维持事物存在的工具与事物本身相关联。\par

% structure: p0044 source=h2 target=section reason=relative-heading-rank; base=h2

\section{第十九章 理智与人灵魂同时存在。来自\Person{亚里士多德}的例子被转化为支持这些观点的证据。}

我们也不可忽略那些剥夺灵魂拥有理智——即使是很短一段时间——的作者。\Emph{他们这样做}是为了为在生命较晚的时期引入理智——以及心智——铺路，甚至是在人的理智出现之时。他们主张婴儿期仅靠灵魂支撑，仅仅是为了促进活力，而并无获取知识的意图，因为并非所有拥有生命的受造物都拥有知识。例如，引用\Person{亚里士多德}的例子，树木拥有生命力，但没有知识；\Emph{与他一致}的是每一位赋予所有有生命之物以动物实体之份额的人——按照我们的观点，这动物实体只存在于人之中，作为他特有的财产——不是因为它是天主的化工（其他所有受造物也都是天主的化工），而是因为它是天主的嘘气，而这（人的灵魂）是唯一拥有这嘘气的，我们声称它出生时即已装备齐全，拥有其特有的官能。好吧，让他们用树木的例子来反驳我们；我们接受他们的挑战，（我们也不会发现这对我们的论证有任何损害）；因为一个不争的事实是：当树木尚为嫩枝和幼芽，甚至尚未达到树苗阶段时，一旦它们从自然的苗床中萌发，它们自身就拥有其特有的生命力官能。但是，随着时间的推移，树木的活力缓慢增长，随着它生长并硬化成木质的树干，直到它的成熟年龄完成了自然为其命定的状态。否则，树木在适当的时候会拥有什么资源用于嫁接、形成叶子、花蕾膨胀、花朵优雅地飘落以及汁液软化？如果不是因为其本性全套装备的悄然生长，以及这生命力分布到所有枝干以实现其成熟？因此，树木拥有能力或知识；它们从哪里获得这些，它们就从哪里获得生命力——也就是说，从那个对它们的本性而言独特的生命力与知识的唯一源泉，并且\Emph{那}是从它们也开始的婴儿期就有的。因为我观察到，即使葡萄树虽然柔软而未成熟，仍然理解它自己的天然事务，并努力攀附某种支撑物，以便倚靠其上，缠绕其间，从而得以生长。事实上，不等农夫训练，没有篱笆，没有支柱，无论其卷须抓住什么，它都会亲切地缠绕，并以真正的更大的坚韧和力量——凭借其自身的倾向而非你的意志——紧紧拥抱。它渴望并急于变得稳固。再看看常春藤，无论多么幼嫩：我观察到它们从最初就试图抓住上方的物体，并且超越一切，挂到最高的东西上，因为它们宁愿用它们的叶状网格铺满墙壁，也不愿在地上爬行，被任何喜欢践踏它们的脚碾碎。另一方面，对于那些因接触建筑物而受损的树木，它们如何在生长时避开并避免伤害它们！你可以看到它们的枝条天生就打算朝向相反方向，并且你可以很好地理解这种树木从其避开墙壁的行为中所显示的生命本能。它满足于（即使只是一棵小灌木）自己微不足道的命运，它从婴儿期就通过其预见的本能完全意识到了这一点，只是它仍然害怕即使是一座废墟的建筑。那么，在我这边，我为何不应为树木这些明智而有远见的本性辩护？让它们拥有生命力，如同哲学家所允许的那样；但让它们也拥有知识，尽管哲学家否认这一点。那么，甚至一根木头的婴儿期也可能拥有（适合它的）理智：何况一个人呢？人的灵魂（可比作一棵树的新生嫩芽）源自\Person{亚当}作为其根，并通过女人（灵魂被托付给她以传递）在\Person{亚当}的后裔中繁殖，从而带着其所有自然的装备——包括理智和感觉——萌发进入生命！如果我大错特错，那么人，甚至从婴儿期起，当他用婴儿的哭声向生命致意时，难道不正是通过他出生的事实来证明他实际拥有感觉和理智的官能吗？同时主张所有感官的使用——视觉通过光，听觉通过声音，味觉通过液体，嗅觉通过空气，触觉通过地面。因此，婴儿期最早的这哭声，是感官的第一次努力，也是心智感知的初始冲动。还有进一步的事实：有些人理解婴儿的这哀哭是对未来泪眼人生的苦难的预兆，借此从出生那一刻起，（灵魂）就被视为具有先见之明，更不用说智力了。因此，凭着这种直觉，婴儿认识他的母亲，辨别奶妈，甚至认出女仆；拒绝另一个女人的乳房，拒绝不是他自己的摇篮，只渴望他习惯的怀抱。那么，他从哪里获得对新颖和习惯的这种辨别力，如果不是来自本能的知识？他如何会烦躁和安静下来，如果不是凭借他最初的理智？如果婴儿天性如此活泼，却没有心智力量；天性如此敏感和易受情感影响，却没有理智，那将是非常奇怪的。但是（我们持相反的见解）：因为基督通过\Quote{从婴孩和哺乳者口中领受赞美}，已经宣告了童年和婴儿期都不是没有感觉的——前者在欢呼迎接祂时，证明了它能够向祂作证；\BibleRef{玛 21:15}而后者因祂的缘故被屠杀，当然知道暴力意味着什么。\BibleRef{玛 2:16}\par

% structure: p0046 source=h2 target=section reason=relative-heading-rank; base=h2

\section{第二十章 灵魂就其本性而言是统一的，但其官能各有发展。各种差异只是偶然的。}

因此，我们在此得出结论：灵魂的一切自然属性都作为其本体的部分而内在于它；并且，从它出生的那一刻起，这些属性就与它一同成长和发展。正如我们常发现站在我们这边的\Person{塞内加}所说：\Quote{在我们内植入了各种技艺和人生各时期的种子。天主，我们的导师，暗中产生我们的心智倾向；}也就是说，从通过婴儿期植入并隐藏在我们内的胚芽中——这些胚芽就是理智；因为我们的自然倾向由此演化而来。同样，甚至植物的种子，每个种类也都有一个形态，但它们的发育却各不相同：有些健康完美地开放和舒展，而另一些则因天气和土壤的条件，以及劳动和照料的运用而改善或退化；也因季节的变迁和偶然情况的发生而改变。同样，灵魂在其种源上可能是一致的，尽管通过出生的过程而变为多样的。在此，也必须考虑当地的影响。据说在\Place{底比斯}出生的人愚钝粗野；而在\Place{雅典}出生的人则最有智慧与辩才，在\Place{科利托斯}地区，婴儿未满月就能说话——他们的舌头如此早熟。事实上，\Person{柏拉图}本人在\WorkTitle{蒂迈欧篇}中告诉我们，当\Person{密涅瓦}准备建立她伟大的城市时，她只考虑了那片土地的性质，它预示了这种心智倾向；因此，他本人在\WorkTitle{法律篇}中教导\Person{墨吉罗斯}和\Person{克利尼亚斯}在选择城址时要小心。而\Person{恩培多克勒}则将智慧或迟钝的原因归于血液的质量，他由此推导出学习和科学中的进步与完善。民族特性的主题此时已变得脍炙人口。喜剧诗人嘲笑\Person{弗里吉亚人}的懦弱；\Person{撒路斯提乌斯}指责\Person{摩尔人}的轻浮，\Person{达尔马提亚人}的残忍；连\Person{宗徒}也斥责\Person{克里特人}为\Quote{说谎者}\BibleRef{铎 1:12}。很可能，也必须将某些原因归于身体状况和健康状态。肥胖妨碍知识，而瘦削的身体则刺激知识；瘫痪使心智萎靡，消耗性的疾病却保存心智。更何况那些偶然的情况——它们除了身体或健康状态外，还会使才智敏锐或迟钝——更应被注意！才智因学识追求、科学、艺术、实验知识、商业习惯和研究而敏锐；因无知、懒惰、不活动、情欲、无经验、无精打采和邪恶追求而迟钝。然后，除了这些影响，或许还必须加上至高的力量。这些至高的力量是：按照我们（基督徒）的观念，是主天主和祂的仇敌魔鬼；但按照人们关于天主眷顾的普遍意见，是命运与必然性；关于运气，则是人的自由意志。连哲学家也承认这些区分；而我们方面，已经按照（基督）信仰的原则，在另一部作品中着手处理它们。显然，那些以如此多样的方式影响灵魂同一本性的力量何其之大，因为它们通常被视为独立的\Quote{\Emph{\OriginalTerm{本性}{natures}}}。然而，它们并非不同的物种，而是同一本性和实体的偶然事件——即天主赋予\Person{亚当}、并使之成为所有（后续）灵魂的模子的那同一本性。它们将永远是偶然事件，但绝不会成为种差。此外，不论现在人类的纷杂多么巨大，但在人类始祖\Person{亚当}那里并非如此。如果这种多样性是由于本性所致，那么所有这些不协调本应存在于他这源头之中，并由此以无损的多样性传递给我们。\par

% structure: p0048 source=h2 target=section reason=relative-heading-rank; base=h2

\section{第21章 正如\OriginalTerm{自由意志}{Free-Will}促使个人行动，其品性亦能改变}

现在，如果灵魂从起初在亚当内就拥有这种统一而单纯的本性，先于后来从它发展出来的众多心理倾向，那么它就不会因这种多样的发展而变为多形，也不会因被预断为\Quote{\Emph{华伦提努派的天主圣三}}的三重形式（以便我们仍可将对该异端的谴责记在心里）而变为多形，因为这种本性在亚当身上也找不到。他有什么是属神的？是因为他预言性地宣告了\BibleQuote{基督与教会的伟大奥迹吗？}\BibleRef{弗 5:32}\BibleQuote{这是我的骨中之骨，肉中之肉；她应称为\Quote{女人}，因为是从男人取出的。为此，人应离开自己的父母，依附自己的妻子，二人成为一体。}\BibleRef{创 2:23}但这（预言的恩赐）只是后来才临到他身上，当时天主把出神或属神品质——预言即在于此——注入他内。再者，如果他身上发展出罪的邪恶，这也不能被算作一种\Emph{自然}倾向；它反而是由（古）蛇的诱惑产生的，与他的本性毫不相干，正如他与\Emph{物质}无关一样，因为我们已排除了对\Quote{物质}的信仰。现在，如果属神的元素，以及异端者所谓的物质元素，都没有恰当地内在于他（因为如果他是由物质创造的，邪恶的种子必然是构成的一部分），那么剩下的便是他本性的唯一原始元素就是所谓的\Emph{魂}（生命的本原，灵魂），我们主张其状态是单纯而统一的。关于这一点，我们还需要探究，既然它被称为自然的，是否应被视为易变的。（我们所提到的异端者）否认本性是可变的，以便他们能够建立并确定其三重理论或\Quote{天主圣三}在其各种本性中的全部特征，因为\BibleQuote{荆棘上收不到无花果，蒺藜上也收不到葡萄。}\BibleRef{路 6:43}如果是这样，那么\BibleQuote{天主就不再能从这些石头给亚巴郎兴起子孙来，也不能使毒蛇的种类结出悔改的果实。}\BibleRef{玛 3:7}如果是这样，那么宗徒在他的书信中也犯了错误，他说：\BibleQuote{你们从前原是黑暗，但现在在主内却是光明；}\BibleRef{弗 5:8}并且\BibleQuote{我们从前也都在本性上身为义怒之子；}\BibleRef{弗 2:3}以及\BibleQuote{你们中有些人从前就是这样，但你们已被洗净了。}\BibleRef{格前 6:11}然而，圣经的陈述永远不会与真理不和谐。一棵坏树永远不会结出好果实，除非更好的本性被嫁接进去；一棵好树也不会结出坏果实，除非经过同样的栽培过程。石头也会成为亚巴郎的子孙，如果他们受教于亚巴郎的信德；毒蛇的种类也会结出悔改的果实，如果他们摈弃那恶毒本性的毒液。这就是天主恩宠的能力，确实比本性更强大，它对我们内在本性之下潜藏的能力——甚至我们的自由意志——行使主权，这种能力被描述为\OriginalTerm{独立自主}{\GreekText{αὐτεξούσιος}}；并且，正因这能力本身也是自然的和可变的，无论它转向何方，都出于自身本性而倾斜。现在，我们内里自然存在着这种\OriginalTerm{独立自主}{\GreekText{τὸ} \GreekText{αὐτεξούσιον}}，我们已经在对玛尔基翁和赫尔谟根尼的驳斥中表明了。那么，如果自然状态必须接受一个定义，它就必须被确定为双重的——有受生和未受生、受造和未受造的范畴。那因受造或受生而获得其存有的，本性上是可以改变的，因为它可以重生和再造；而那未受造和未受生的，将永远保持不变。然而，既然这种状态只适合天主，作为唯一未受生、未受造的存在（因此是不死不灭、不变的），那么其他一切受生和受造的存在本性上必然是可修改和可变的；因此，如果三重状态要归于灵魂，它必须被认为来自其偶性环境的可变性，而非本性的命定。\par

% structure: p0050 source=h2 target=section reason=relative-heading-rank; base=h2

\section{第22章。综述。灵魂的定义}

赫尔谟根尼已经听我们说明了灵魂的其他自然官能是什么，以及它们的辩护和证明；由此可以看出，灵魂是天主的产儿，远胜于物质。这些官能的名称在此简单重复，以免被人遗忘而忽视。我们已经赋予灵魂刚才提到的自由意志，以及它对自然作为的主宰权，还有它偶尔拥有的占卜恩赐，这是独立于专门从天主恩宠而来的预言天赋之外的。我们现在就离开灵魂性向这个主题，以便充分按顺序阐述其各种品质。因此，我们定义灵魂：由天主的嘘气而生，不朽的，拥有形体，具有形式，在其本质上是单纯的，在其本性上具有理智，以各种方式发展其能力，在决定上是自由的，易受偶性变化的影响，在其官能上可变，有理性的，至高的，赋有预感的本能，由一个（原型的灵魂）演化而出。我们现在还需要考虑它如何从这个唯一的原始源头发展出来；换言之，它从何处、何时以及如何产生。\par

% structure: p0052 source=h2 target=section reason=relative-heading-rank; base=h2

\section{第23章。各种异端者的意见，最终源于柏拉图}

有些人以为灵魂从天上降下，其坚信程度就如同他们沉湎于必定回到那里的前景时通常抱持的信念一样。属于\Person{西蒙}派别的\Person{默南德}的门徒\Person{撒图尔尼诺}引入了这一观点：他断言人是由天使造的。起初，它是一件徒劳、不完美的受造物，软弱无力，无法站立，像虫子一样在地上爬动，因为它缺乏保持直立姿态的力量。但后来，由于至高权能的怜悯（人就是按照这权能的形象被笨拙地造成，但这一形象并未被完全理解），获得了微弱的生命火花；这火花唤醒并扶正了他不完美的形体，以更高的活力赋予它生气，并为其在离世后返回本原作了预备。确实，\Person{卡尔波克拉特}为自己索取极其大量的超凡品质，以至于他的门徒立即将自己的灵魂置于与基督同等的地位（更不用说宗徒了）；有时，当他们兴致所至，甚至赋予灵魂优越性——他们认为灵魂确实分享了那种崇高德能，这德能俯视治理这世界的权势。\Person{阿佩莱斯}告诉我们，我们的灵魂被一个火天使（即以色列的天主，也是我们的天主）用人间的诱饵从超天居所引诱下来，然后他将它们牢牢地封闭在我们有罪的肉身之内。\Person{瓦伦廷}学派用\Emph{索菲亚}（即智慧）的胚芽来坚固灵魂；凭借这胚芽，他们在可见事物的形象中，认出他们自己的\OriginalTerm{永世}{Æons}的故事和米利都寓言。我由衷地感到遗憾，\Person{柏拉图}一直充当着所有这些异端者的供应者。因为在\WorkTitle{费多篇}中，他想象灵魂从此世界游荡到彼世界，又从那里返回此地；而在\WorkTitle{蒂迈欧篇}中，他假定被指派去创造可朽受造物的天主儿女们，为灵魂取了不朽的胚芽后，围绕它凝结了可朽的身体——从而表明这世界是某个其他世界的形象。如今，为了让这一切得到相信——即灵魂先前曾与天主一起生活在天上，与祂分享祂的理念，后来降下与我们一同生活在地上，在此地时它回忆先前学过的万物的永恒模式——他精心设计了新的公式，\OriginalTerm{学习就是回忆}{\GreekText{μαθήσεις} \GreekText{ἀναμνήσεις}}，意为\Quote{学习就是回忆}；这暗示着，从那里来到我们这里的灵魂，忘记了它们先前生活所在的事物，但后来，因看到周围对象的教导而回忆起它们。因此，既然异端者从\Person{柏拉图}借用来的教义被这类论证巧妙地辩护，我只要推翻\Person{柏拉图}的论证，就足以驳斥异端者了。\par

% structure: p0054 source=h2 target=section reason=relative-heading-rank; base=h2

\section{第24章 柏拉图的自相矛盾：他以为灵魂自存，却能忘记前世经验}

首先，我不能容许灵魂会丧失记忆；因为他已赋予灵魂如此大量的神圣品质，以致于把它与天主等同。他声称灵魂是\Emph{未生的}，仅此一属性就足以证明其完美的神性；他又补充说灵魂是不朽的、不可朽坏的、无形体的——因为他相信天主也是同一性质——不可见的、不可描画的、单一的、至高的、理性的和理智的。如果他想要称灵魂为天主，他还能赋予灵魂更多什么呢？然而我们，既不允许天主有任何附加物（在平等的意义上），正因此认为灵魂远低于天主：因为我们以为灵魂是受生的，并因此拥有某种稀释过的神性和衰减的福乐，如同天主的嘘气，却非祂的神；且虽然不朽，因这是神性的一种属性，却仍是可受苦难的，因为这是受生之物的伴随特性，因而从起初就能偏离完美与正义，从而易受记忆的丧失。这一点我已与\Person{赫尔摩格内斯}充分讨论过了。但还可以进一步指出，如果灵魂由于其所有品质都与天主的属性相等而可被认作神明，那么它就必须不受任何情欲所左右，因而也不会丧失记忆；因为遗忘这一缺陷对你所断言的对象造成的伤害，正如记忆是其荣耀一样大；\Person{柏拉图}本人认为记忆是感官和理智能力的保障，而\Person{西塞罗}则称之为一切科学的宝库。现在，我们不必提出疑问，即像灵魂这样神圣的能力是否会失去记忆；问题更在于，它是否能够重新找回它所失去的记忆。我无法断定，那本应失去记忆的东西，倘若一旦遭受损失，是否有足够的力量重新记起。实际上，两种可能性都与我的灵魂相合，但与柏拉图的灵魂不合。\newline{}\newline{}其次，我对他的反驳将是这样：\Quote{柏拉图，你是否赋予灵魂一种天然能力，以理解你那著名的\Emph{理型}？当然如此，你将会回答。}好吧，现在没有人会向你承认，你所说是天赋之知、关于自然科学的认识会丧失。但科学的知识会丧失；各门学问和生活技艺的知识会丧失；同样地，我们心灵各种官能和情感的知识也许也会丧失，尽管它们看起来是天赋于我们本性的，但实际上并非如此：因为，正如我们已经说过的，它们受到地点、风俗习惯、身体状况、人的健康状态——受到至高力量之影响和人的自由意志之变化的偶然因素所影响。\newline{}\newline{}如今，对自然事物的直觉知识从不会丧失，即使在野兽中也不例外。狮子无疑会因驯化的软化影响而忘记其凶猛；它可能，带着它美丽的鬃毛，成为某位\Person{贝勒尼基女王}的玩物，并用舌头舔她的脸颊。野兽可能放弃它的习性，但它不会忘记其天然本能。它不会忘记适当的食物，不会忘记天然的资源，也不会忘记天然的恐惧；即使女王给它鱼或饼，它还是会想要肉；如果它生病时给它准备了解药，它仍然需要猴子；如果没有矛枪对着它，它还是会害怕公鸡的啼叫。同样地，在人身上——也许是所有受造物中最健忘的——一切属于他天性的知识将不可磨灭地固定在他内——但仅限于此，因为只有这些是天性本能。他永远不会忘记在饥饿时吃，在口渴时喝；想要看时使用眼睛，听时使用耳朵，闻时使用鼻子，尝时使用嘴巴，触时使用手。这些当然是感官，哲学因偏好理智能力而贬低它们。\newline{}\newline{}但如果感官能力的天然知识是持久的，那么被赋予优越性的理智能力的知识为何会丧失呢？那么，这种先于回忆的遗忘能力本身是从何而来的呢？他说，来自时间的漫长流逝。但这是一种短视的回答。时间的长短不可能属于那按他所说的是未生的东西，因此必然被视为永恒。因为永恒之物，基于其未生的根据，既没有时间的起点也没有终点，不受任何\Emph{时间性}标准的衡量。而时间不能衡量的事物，不会因时间而发生变化；时间的漫长流逝对它也毫无影响。如果时间是一个遗忘的原因，那么为什么从灵魂进入身体之时起，记忆就丧失了，仿佛从此以后灵魂就受时间影响？因为灵魂无疑先于身体，当然并非与时间无关。那么，遗忘是发生在灵魂进入身体的瞬间，还是之后一段时间？如果是瞬间，那么在假设中尚不允许的时间漫长流逝在哪里？举例来说，婴儿的情况。如果是在之后一段时间，那么在遗忘发生前的间隔期间，灵魂难道不是仍然运用其记忆能力吗？而且，灵魂后来又如何会遗忘，而后又重新记起呢？此外，遗忘压抑灵魂的这段时间必须被看作有多长？我认为，人一生的全部时间都不足以抹去灵魂在进入身体之前所经历的那段漫长岁月的记忆。\newline{}\newline{}但是，柏拉图又把责任推给身体，好像一个受生的实体竟能熄灭一个未生的实体的能力，这竟能令人置信。然而，在身体之间存在着许多差异，因其理性、大小、状况、年龄和健康而不同。那么，我们是否应当认为在遗忘上也有类似的差异？然而，遗忘是统一且同一的。因此，身体的特殊性，以其多种多样的变化，不会成为一个恒定不变之效果的原因。同样，根据柏拉图本人的见证，有许多证据表明灵魂具有占卜能力，正如我们先前对\Person{赫尔摩格内斯}所主张的那样。但没有一个活着的人，自己不感到其灵魂拥有对某种预兆、危险或喜悦的预感。那么，如果身体对占卜无害，我想它对记忆也无害。有一点是确定的，处在同一个身体中的灵魂既会遗忘也会记忆。如果某种身体条件导致遗忘，它又如何容许相反的记忆状态？因为遗忘之后的记忆实际上是记忆的复活。那么，最初对记忆怀有敌意的东西，第二次又如何不会对它有害呢？\newline{}\newline{}最后，谁比小孩子有更好的记忆力？他们拥有新鲜、未磨损的灵魂，尚未沉浸于家务和公共事务的忧虑中，而只专注于那些学习，其掌握本身就是一种回忆。确实，为什么我们大家不是同样地记忆，既然我们在遗忘上是平等的？但这只适用于哲学家！甚至不是全体哲学家。在这么多民族中，在如此众多的贤哲中，柏拉图的的确确是唯一一个将\Emph{理型}的遗忘和回忆结合起来的人。既然他这一主要论点根本站不住脚，那么它的整个上层建筑也必然随之倒塌，即：灵魂被认为是未生的，生活在天上领域，并在那里得到神圣奥秘的教导；此外，他们下降到这个世界，在这里回忆起他们之前的存在，当然，目的是为我们的异端者提供适合他们体系的材料。\par

% structure: p0056 source=h2 target=section reason=relative-heading-rank; base=h2

\section{第二十五章 \Person{戴尔都良}从生理学上驳斥灵魂在出生后注入的观点}

我现在要回到这个离题的原因，以便解释所有灵魂如何源自同一个，以及它们在何时、何地、以何种方式产生。关于这个问题，无论它是由哲学家、异端者还是普通人提出，都无关紧要。那些宣称真理的人并不关心他们的反对者，尤其是那些一开始就坚持认为灵魂不是在子宫中受孕，也不是在肉体成形时形成和产生，而是在婴儿完全有生命之前，但在分娩过程之后从外部印入的人。他们还说，人的种子在\OriginalTerm{交合之后}{ex concubiter}被适当安置在子宫中，并通过自然冲动获得生命后，凝聚成单纯的肉体物质，在适当的时候从子宫的炉膛中温暖地诞生，然后从其热量中释放出来。这肉体就像热铁的情况，在那种状态下被浸入冷水中；因为它被（它诞生的）冷空气击中，立即接受生气的能力，并发出声音。这种观点为斯多亚派和\Person{埃涅西德谟斯}所持有，\Person{柏拉图}本人也偶尔持有，当他告诉我们，灵魂是一种完全独立的构成，起源于子宫之外和外部，在新生婴儿第一次呼吸时被吸入，并随着人的最后一息呼出。我们将看到他的这种观点是否仅仅是虚构的。就连医学界也不缺少它的\Person{希塞修斯}，他背叛了自然和自己的职业。我想，这些先生们太过腼腆，不愿与妇女就分娩的秘密达成共识，而这些秘密对后者来说是众所周知的。但还有更多让他们脸红的事，因为最终是妇女驳斥了他们，而不是称赞他们。现在，在这样一个问题上，没有人能比与这个问题密切相关的女性本身更有效的教师、法官或证人了。那么，请你们作证吧，母亲们，无论是怀孕期间还是分娩之后（让不生育的妇女和男人保持沉默），——你们本性的真相正在被质疑，你们受苦的现实正是要决断的点。（那么请告诉我们，）你们是否感到在你们体内的胚胎中有任何不同于你们自己的生命力，你们的肠子因此颤抖，你们的肋旁因此摇动，你们的整个子宫悸动，压迫你们的负担不断改变位置？这些动作对你们来说是喜悦，是真正的消除焦虑，使你们确信你们的婴儿既有生命又享受生命？或者，如果他的不安停止，你们首先为他担忧；他在你们体内会意识到这一点，因为他被新的声音打扰；你们会渴望有害的饮食，甚至厌恶食物——全部是为了他；然后你们和他（在你们同情的亲密中）会分担共同的病痛——以至于你们的瘀伤和挫伤会实际上在他身上留下印记——而他在你们体内，甚至在身体的相同部位，如此专横地将母亲的伤害归于自己！现在，每当瘀青和发红是血液的附带事件时，血液就不会没有生命原理或灵魂；或者当疾病攻击灵魂或生命力时，（这就成为它真实存在的证据，因为没有灵魂或生命原理就没有疾病。此外，既然通过食物维持生命和缺乏食物、生长和衰退、恐惧和运动都是灵魂或生命的条件，那么经历这些的人一定活着。因此，最终当一个人停止经历这些时，他就停止了活着。就这样，婴儿随后成为死胎；但如果不是他们曾有生命，怎么会这样呢？因为一个人怎么能死，如果他没有先活过？但有时由于残酷的必要性，婴儿还在子宫内就被处死，当他横躺在子宫口阻碍分娩时，如果他自己不死，就会杀死母亲。因此，在外科医生的工具中，有一种仪器，它有一个精心调节的柔性框架，用于首先打开\OriginalTerm{子宫}{uterus}并保持打开；它还配备了一个环形刀片，通过它，子宫内的四肢被焦虑但坚定的细心解剖；它的最后一个附件是钝的或覆盖的钩子，通过暴力分娩将整个\OriginalTerm{胎儿}{fœtus}取出。还有另一种仪器，形状像铜针或尖刺，通过它在这个秘密的劫掠生命行为中实际执行死亡：他们根据其杀婴功能称它为\OriginalTerm{杀胎器}{\GreekText{ἐμβρυοσφάκτης}}，即婴儿的杀戮者，婴儿当然是有生命的。\Person{希波克拉底}、\Person{阿斯克勒庇阿德斯}、\Person{埃拉西斯特拉图斯}、\Person{希罗菲卢斯}（他甚至解剖成年人）以及更温和的\Person{索拉努斯}本人都有这样的器械，他们都很清楚一个活生生的存在已经被孕育，并同情这个最不幸的婴儿状态，它必须先被杀死，以避免被活活折磨。对这种严厉处理的必要性，我毫不怀疑甚至\Person{希塞修斯}也深信不疑，尽管他将婴儿的灵魂归入出生后来自冷空气的打击，因为灵魂这个术语在希腊文中恰好对应于这种冷却！那么，野蛮人和罗马民族是通过其他过程获得灵魂的，我想知道吗？因为他们用不同于\OriginalTerm{灵魂}{\GreekText{ψυχή}}的名称来称呼灵魂。有多少民族在炽热的太阳下开始生活，将他们的皮肤晒成黝黑色？他们从哪里获得灵魂，没有寒冷空气的帮助？我不说那些温暖的卧室和产妇所需的所有加热装置，因为一丝冷空气可能会危及她们的生命。几乎就是在浴室里，一个婴儿也会滑入生命，立刻听到他的哭声！然而，如果优质的寒冷空气对灵魂来说是如此不可或缺的宝藏，那么除了日耳曼人和斯基泰人部落，以及阿尔卑斯山和阿格山的高地之外，没人应该出生！但事实是，在东方和西方的温带地区人口更多，人的心智更敏锐；而没有一个萨尔马提亚人不是头脑迟钝和平庸的。人的心智也会因寒冷而变得敏锐，如果灵魂是在刺骨的严寒中产生的；因为物质是什么，其能力也必然是什么。现在，在这些预备性陈述之后，我们也可以提及那些从母腹中被剖出、呼吸并保持生命的人的例子——你们的\Person{巴克斯}和\Person{西庇阿}们。然而，如果有人像\Person{柏拉图}那样，认为两个灵魂不能像两个身体那样共存于同一个人体中，相反地，我可以向他展示，不仅两个灵魂共存于一个人中，就像两个身体共存于同一个子宫中，而且许多其他事物在自然关联中与灵魂结合——例如，魔鬼附身；而且\Emph{那}不仅仅是一个灵魂，就像\Person{苏格拉底}自己的魔鬼那样；而是七个魔鬼，就像\Person{玛利亚玛达肋纳}的情况那样（\BibleRef{谷 16:9}），以及许多的数量，就像加达辣人那样（\BibleRef{谷 6:1}）。现在，一个灵魂自然比一个邪灵更容易与另一个灵魂结合，因为相似的物质，而邪灵由于本质不同更困难。但当同一位哲学家在\WorkTitle{法律篇}第六卷中警告我们要当心，以免种子的缺陷从非法或卑贱的姘居中将污秽注入身体和灵魂时，我几乎不知道他是更与他之前的陈述矛盾，还是与他刚刚做出的陈述矛盾。因为他在这里向我们展示灵魂来自人的种子（并警告我们要对此当心），而不是（如他之前所说）来自新生婴儿的第一次呼吸。请问，如果我们不是从这灵魂的种子中产生，那么根据\Person{克莱安西斯}的见证，我们怎么会与父母在性情上相似呢？为什么古老的占星家从一个人最初受孕时起就为他算命，如果他的灵魂也不是从那一刻开始起源？与此（诞生）同样也属于灵魂的吹入，无论那是什么。\par

% structure: p0058 source=h2 target=section reason=relative-heading-rank; base=h2

\section{第二十六章 唯有圣经对我们所争论的问题提供清晰的知识}

论人的意见，其不确定与不规则，直到我们到达天主所设定的界限，方才止息。我最终要退守到自己的阵地，并坚稳地持守在那里，为了向基督徒证明我回答哲学家和医生的话是确当的。\newline{} 弟兄（在基督内），在你自己基础上建立你的信德。请看最圣善的妇人的子宫，充满内在的生命，她们的婴孩不仅在其中呼吸，甚至拥有先知的直觉。请看黎贝加的肠腹如何不安，\BibleRef{创 25:22} 虽然她离生产尚远，也没有气息的冲动。看哪，双胞胎在母腹中相争，虽然她尚无两族的迹象。或许我们会把这幼小后代的争斗视为奇事——他们在活着之前就相争，在有了生命之前就有敌意——如果他们仅是因在腹中不安而搅扰了母亲的话。但当她的子宫张开，后代的数目显明，他们被预言的状态也被知道时，我们便得到一个证据，不仅证明婴孩的（个别）灵魂，也证明他们敌对的挣扎。\newline{} 那首先出生的，被那出生在前、尚未完全生出、但仅手已生出者威胁要扣留。现在，如果他实际上是在第一次呼吸时才获得生命、以柏拉图式的方式接受灵魂；或者，按照斯多亚的规则，在接触寒冷空气时才初次尝到生命；那么，那急切期盼、仍被扣留在子宫内、并试图扣留（另一个）在外面的另一个婴孩在做什么呢？我想他还没有呼吸就抓住了他兄弟的脚跟，\BibleRef{创 25:26} 并且当他如此强烈地想要第一个离开子宫时，他仍带着母亲的热度。何等的一个婴孩！如此好胜，如此强壮，且已如此好斗；而这一切，我想，是因为他现在就已充满生命！\newline{} 再看那些不寻常的怀孕，更为奇妙的是不孕和童女的怀孕：这些妇人只能违反自然产生不完美的后代，因为其中一个年迈无法孕育种子，另一个则与男人无接触而洁净。如果在这情形中要生产，那么按照（哲学家的说法），那些不正常受孕的婴孩应无灵魂而生。\newline{} 然而，这些婴孩也有生命，各在母腹中。依撒伯尔欢喜雀跃，（因为）若翰在她腹中跳跃；\BibleRef{路 1:41} 玛利亚颂扬主，（因为）基督在她腹中激动。\BibleRef{路 1:46} 母亲们各自认出自己的后代，且也被她们的婴孩认出，这些婴孩当然活着，不仅是灵魂，而且是精神。\newline{} 因此你读到天主对耶肋米亚说的话：\BibleQuote{在我形成你于母腹之前，我已认识了你。}\BibleRef{耶 1:5} 既然天主在子宫中形成我们，祂也向我们嘘气，如同在最初的创造时一样：\BibleQuote{上主天主形成了人，在他鼻孔内吹了生气。}\BibleRef{创 2:7} 天主除非在人的整个本性中，否则不能在子宫中认识人：\Quote{在你从母腹出来以前，我已祝圣了你。} 那么，那时它是个死尸吗？当然不是。因为\Quote{天主不是死人的天主，而是活人的天主。}\par

% structure: p0060 source=h2 target=section reason=relative-heading-rank; base=h2

\section{第二十七章 灵魂与身体在元素中同时被受孕、形成并完成}

那么，生命体是如何受孕的呢？身体和灵魂的本质是同时一起形成的吗？还是其中一个在自然形成上先于另一个？我们坚持认为，两者是同时受孕、形成、并完美地一同出生；在受孕中没有片刻的间隔，以至于不能为两者中的任何一个指定更早的位置。事实上，你可以通过人最终所发生的事来判断他最初存在的事。既然死亡的定义无非是身体和灵魂的分离，那么生命，作为死亡的对立面，也只能被定义为身体与灵魂的结合。如果分离通过死亡在同一时刻发生于两种实体，那么结合法则也应使我们确信，通过生命，这两种实体是同时结合的。现在我们承认生命始于受孕，因为我们主张灵魂也始于受孕；生命与灵魂在同一时刻、同一地点开始。因此，共同导致死亡分离的过程，也同时结合而产生生命。如果我们给其中一个本性赋予（形成）的优先性，而给另一个较晚的时间，那么我们还需根据各自的条件和等级确定播精的精确时间。若是如此，我们应给身体的种子什么时间？又给灵魂的种子什么时间？此外，如果由于时间差异而给播精分配不同的时期，那么实体也会不同。因为虽然我们承认有两种种子——身体的种子和灵魂的种子——但我们仍宣称它们不可分离，因此在起源上是同时的。现在，不要有人因（为真理辩护所必需的）对自然过程的解释而感到冒犯或羞愧。自然对我们来说应是敬畏的对象，而非羞耻的对象。是情欲，而非自然的用途，使两性的交合蒙受了羞耻。是过度，而非正常状态，是不端庄、不贞洁的；正常状况已蒙天主祝福，并受祂祝福：\BibleQuote{你们要生育繁殖，（充满大地。）}\BibleRef{创 1:28} 然而，祂咒诅了过度，即奸淫、放荡和纵欲。如今，在两性寻常的功能中，即男女在共同交合中结合时，我们知道灵魂和肉身共同履行一项职责：灵魂提供欲望，肉身提供满足；灵魂提供刺激，肉身提供实现。整个人因两种本性的共同努力而激动，他的精液被排出，其流动性来自身体，其温暖来自灵魂。如果希腊语中\Emph{灵魂}一词与\Emph{冷}同义，那么为何身体在灵魂离开后会变冷呢？事实上（即使我冒着冒犯谦逊的风险，渴望证明真理），我不禁要问：在极度满足的炽热中，当生殖液射出时，我们是否感到自己的灵魂有一部份离我们而去？我们不是也经历虚弱、虚脱，同时视线模糊吗？那么，这一定是产生灵魂的种子，它来自灵魂的滴落，正如那种液体是产生身体的种子，来自肉身的排泄。最初的创造的例子最为真实。亚当的肉身由泥土形成。泥土岂非绝佳的湿气，生殖液由此而生？灵魂首先来自天主的嘘气。但天主的嘘气除了是灵的蒸气之外，又是什么呢？我们通过生殖液所呼出的东西岂非由此而生？因此，既然这两种不同而分离的实体——泥土和嘘气——在最初的创造中结合形成了单个的人，它们便将各自适当的精原要素融合并混合为一，并从此将人类繁殖的正常方式传递给人类；以至于如今，这两种实体尽管彼此不同，却同时从一条统一的通道流出；它们一起进入指定的精床，用结合的力量，从各自的本性中孕育出人类的果实。在这人类产品中，蕴含着自身的种子，这是按照为所有具有生殖功能的受造物所制定的过程。因此，从这一个（原始的）人涌出了人类灵魂的全部流溢和丰富——自然证实自己忠于天主的诫命：\BibleQuote{你们要生育繁殖。}\BibleRef{创 1:28} 因为在这单一创造的序言中，\Quote{让我们造人}，人类的所有后裔已被用复数短语宣告并描述：\Quote{让\Emph{他们}治理海中的鱼}，等等。这不足为奇：在种子中蕴含着收获的许诺和定金。\par

% structure: p0062 source=h2 target=section reason=relative-heading-rank; base=h2

\section{第二十八章 毕达哥拉斯轮回论概述与批判}

那么，时至今日，\Person{柏拉图}所提及的那个关于灵魂相互迁移的古旧说法究竟是什么意思：灵魂如何从这里离开到那里去，然后又回到这里，经历一生，接着再次离开此生，之后又从死者中复活？有些人认为这是\Person{毕达哥拉斯}的话；\Person{阿尔比努斯}则认为这是一项神圣的宣告，或许出自\Place{埃及}的\Person{墨丘利}。然而，除了独一无二的真天主之外，并无神圣的话语；先知们、宗徒们以及\Person{基督}本人正是藉着天主宣告了他们伟大的信息。\Person{摩西}比\Person{萨图恩}还要古老得多（大约早九百年左右），甚至比他的孙子们更古老；而且他无疑更为神圣，因为他从世界之初开始记述并追溯人类的历程，按名字和年代指明（人类祖先的）各次出生；这样从他的神圣权威和话语中，清楚证明了他工作的神圣性质。事实上，如果\Place{萨摩斯}的诡辩家是\Person{柏拉图}关于灵魂因死与活状态的永恒交替而循环迁移的权威来源，那么著名的\Person{毕达哥拉斯}，尽管在其他方面卓越，为了编造这样的意见，必定依赖于一个不仅可耻而且危险的谎言。你们这些不知情的人，请思考一下，并与我们一起相信。他假装死亡，把自己藏在地下，将自己判罚在那样的忍耐中约七年之久，在此期间，他从他唯一的同谋和照料者——他的母亲那里，学到他应该向世界讲述关于他隐居期间去世之人的事，以使世人相信；当他以为已成功将自己的身体弄得像一个死去的老人那样可怕时，他从隐藏和欺骗的地方出来，假装从死者中返回。谁会怀疑那个被认为已经死去的人又复活了呢？尤其是在听到他讲述关于最近死去之人的事情之后，这些事他显然只能是在\Place{哈得斯}本身发现的！因此，人死后重新活过来，倒是一个古老的说法。但假如它也是近期的说法呢？真理不追求古老，谬误也不避讳新颖。我认为这个著名的说法显然是虚假的，尽管它因古老而尊贵。那依靠谎言作为证据的，怎么可能是真的呢？——我怎能不相信\Person{毕达哥拉斯}是骗子，他为了赢得我的信任而行骗？他如何说服我：在成为\Person{毕达哥拉斯}之前，他曾是\Person{埃塔利德斯}、\Person{欧福尔布斯}、渔夫\Person{皮洛斯}和\Person{赫尔墨提摩斯}，以便让我们相信人死后能再次活着——而他自己后来作为\Person{毕达哥拉斯}却发了假誓？正如我相信他曾亲自复活一次比相信他曾在这人或那人身上多次复活更容易一样，同样他在难以相信的事情上欺骗了我，因为他在易于相信的事情上行了欺诈。好吧，但他认出了从前献在\Place{德尔斐}的\Person{欧福尔布斯}的盾牌，声称那是他自己的，并通过普遍不为人知的标记证明了这一声称。现在，再看看他地下的藏身之处，如果可能的话，相信他的故事吧。因为，对于那个设计了如此诡计损害自己健康的人，欺诈性地浪费生命，在地下饥饿、懒散和黑暗之中折磨自己七年——对浩瀚的天空深感厌恶——他会做出怎样鲁莽的努力，尝试怎样奇妙的发明，以发现这著名的盾牌呢？假设他是在那些隐秘的研究中找到了它；假设他捕捉到了一些关于现已过时的传统的微弱传闻；假设他是通过贿赂看守人而得以查看从而知道了它——我们很清楚巫术技巧在探索隐秘之事上的资源：有 \OriginalTerm{击倒之灵}{\GreekText{καταβολικοὶ} \GreekText{πνεύματα}}，它们击倒受害者；有 \OriginalTerm{陪伴之灵}{\GreekText{πάρεδροι} \GreekText{πνεύματα}}，它们常伴左右纠缠他们；还有 \OriginalTerm{占卜之灵}{\GreekText{πυθωνικοὶ} \GreekText{πνεύματα}}，它们通过占卜和腹语术使他们入迷。难道我们的\Person{毕达哥拉斯}的老师\Person{斐瑞居德斯}不也正是通过这类技艺和计谋行占卜，或者我更愿意说是发狂和做梦吗？难道同一个恶魔不也曾在他身上，当他在\Person{欧福尔布斯}体内时，行过流血之事吗？但最后，为什么那个通过盾牌的证据证明自己是\Person{欧福尔布斯}的人，却没有认出他任何一个\Place{特洛伊}时代的同伴呢？因为此时他们也必定已经恢复了生命，既然人们正从死者中复活。\par

% structure: p0064 source=h2 target=section reason=relative-heading-rank; base=h2

\section{第二十九章 从\Person{毕达哥拉斯}学说的首要原则（即活人由死者形成）反驳该学说}

确实，死人是由活人形成的，这显而易见；但不能由此推论活人是由死人形成的。因为在起初，活人在事物次序中在先，因此同样在起初，死人在次序中在后。但死人除了来自活人之外，并无其他来源。活人的起源来自任何其他来源（随你选择）而非来自死人；而死人除了来自活人之外，并无来源使其开始。那么，如果从一开始活人并非来自死人，为何后来（据说）来自死人？那最初的来源，无论它是什么，是否已经止息？其形式或法则是否令人后悔？那么为何在死人身上又得以保存？难道不是这样：因为起初死人来自活人，所以死人总是来自活人？因为要么起初有效的法则在其两个关系中都继续存在，要么它在两者中都改变了；因此，如果后来活人必须从死人中产生，那么同样，死人也不应该从活人中产生。因为如果对制度的忠实遵守并不打算在每个方面延续，那么相反者就不能恰当地交替地由相反者重新形成。我们也要向你们提出某些相反者：出生与未出生、视力与盲目、青年与老年、智慧与愚昧。然而，不能因为相反者来自相反者，就认为未出生者来自出生者；也不能因为盲目发生在视力之后，就认为视力来自盲目；也不能因为青年之后是衰老的暮年，就认为青年从老年中复活；也不能因为智慧有时可能从愚昧中磨砺出来，就认为愚昧从智慧中产生。\Person{阿尔比努斯}为他的（老师和朋友）\Person{柏拉图}在这些点上感到担忧，并煞费苦心地用巧妙的方法区分不同种类的相反者；仿佛这些例子不像那些他用来阐明其伟大导师原则的例子——即生与死——那样绝对具有相反性的本质。而且，就这件事而言，生命从死亡中恢复的说法并不成立，因为死亡确实发生在生命之后。\par

% structure: p0066 source=h2 target=section reason=relative-heading-rank; base=h2

\section{第三十章 进一步驳斥毕达哥拉斯理论。当代文明的状态}

但是，对接下来要说的我们能回复什么呢？因为，首先，如果活人来自死人，正如死人来自活人，那么人类的数量必然始终保持不变，即最初引入（人类）生命的那个数量。活人在死人之前，之后死人来于活人，然后又活人来自死人。既然这个过程始终发生在同一些人身上，那么他们出于同一来源，数量必然始终相同。因为（进入生命）的人绝不会比（在死亡中）消失的人更多或更少。然而，我们在人类古代纪事中发现，人类种族随着人口逐渐增长而发展，要么作为土著占据地球的不同部分，要么作为游牧部落、流亡者或征服者——如\Place{帕提亚}的\Person{斯基泰人}、\Place{伯罗奔尼撒}的\Person{泰梅尼德人}、\Place{亚洲}的\Person{雅典人}、\Place{意大利}的\Person{弗里吉亚人}和\Place{非洲}的\Person{腓尼基人}；或者通过更常见的移民方式，他们称之为\OriginalTerm{殖民地}{\GreekText{ἀποικίαι}}或\Emph{殖民地}，目的是释放过剩的人口，将过度拥挤的人群倾泻到其他住所。土著仍然留在他们的旧居，并且还通过借贷甚至更大的人口来充实其他地区。显然，如果看整个世界，它每天都比古代得到更好的耕种和更充分的人口。现在所有地方都可到达，所有地方都为人所知，所有地方都对商业开放；最宜人的农场已抹去了曾经荒凉险恶的荒地的所有痕迹；耕地已征服了森林；畜群驱走了野兽；沙地已被播种；岩石已被种植；沼泽已被排干；曾经几乎孤零零的小屋之处，现在是大城市。不再惧怕（野蛮的）岛屿，也不惧怕岩石海岸；到处都是房屋、居民、定居政府和文明生活。最常出现在我们眼前（并引起抱怨）的是我们稠密的人口：我们的数量使世界不堪重负，世界几乎无法从其自然元素中供应我们；我们的需要越来越迫切，所有人口中的抱怨越来越苦涩，而自然无力提供其惯常的滋养。事实上，瘟疫、饥荒、战争和地震必须被视为国家的补救措施，是修剪人类繁茂的手段；然而，当斧头一次砍倒大量人群时，世界迄今为止从未因看到其死人千年放逐后复活归来而惊恐。但如果死人确实重新回到生命，这种景象本应通过死亡损失与生命恢复的平衡而变得相当明显。然而，为什么是在一千年之后，而不是在当下，这种从死亡中的返回要发生？假设损失没有立即得到补充，那么必然有完全灭绝的危险，因为补偿之前的失败。事实上，我们今生这段休假与一千年的时期极不相称；它短得多，因此其火炬熄灭比重新点燃容易得多。既然如此，根据我们讨论的假设，如果活人要由死人形成，本应出现的间隔时期并未实际发生，那么结论就是我们不应相信人从死亡中回到生命（正如这种哲学所假想的那样）。\par

% structure: p0068 source=h2 target=section reason=relative-heading-rank; base=h2

\section{第三十一章 对轮回的进一步揭露，其不可摆脱的困境}

再者，如果这从死者中复生确实发生，那么个体当然必须恢复他们各自的个体性。因此，那些赋予各个身体生命的灵魂，必须已经分别返回到它们各自的躯体。现在，每当两个、三个或五个灵魂（正如它们经常那样）被重新封闭在一个子宫中时，在这种情况下就不算是从死者中复生，因为没有个体应当有的分别恢复。虽然按照这种说法，（无疑地）原始受造的法则被显著地遵守，因为仍然从仅一个（灵魂）产生出多个灵魂！再者，如果灵魂在人类生命的不同年龄离去，它们又如何能以一致的年龄回来呢？因为所有人出生时都被赋予一个\Emph{婴儿}灵魂。但一个年老去世的人怎么能以婴儿身份复活呢？如果灵魂在脱离肉体时，因为年龄的倒退而如此缩小，那么经过一千年后，它应当以更丰富的生命成就进步来恢复生命，岂不是更合理得多！无论如何，它应当返回它去世时所达到的年龄，以便恢复它曾放下的确切生命。但即使按照这种说法，它们在循环中永远以相同面貌重现，它们也应当带回，如果不是完全相同的身体形式，至少它们原本的品格、品味和性情的特质，因为如果它们缺少那些证明其身份的特征，就几乎不可能被视为相同。（然而，你以这个问题来反驳我）：\Quote{你怎么可能知道一切不都是隐秘的过程呢？一千年的作为难道不会夺去你辨认的能力，因为它们以陌生的面貌返回吗？}但我十分确定情况并非如此，因为你自己向我提出\Person{毕达哥拉斯}是（恢复的）\Person{欧福耳布斯}。现在请看\Person{欧福耳布斯}：他显然拥有一个军事好战的灵魂，正如神圣盾牌的名声所证明的。至于\Person{毕达哥拉斯}，他却如此隐退，如此不好战，以致他回避当时\Place{希腊}充满的战绩，宁愿在\Place{意大利}的宁静退隐中致力于几何学、占星学和音乐的研究——在品味和性情上与\Person{欧福耳布斯}完全相反。再者，他所代表的\Person{皮洛士}以捕鱼为生；但\Person{毕达哥拉斯}恰恰相反，从不碰鱼，甚至戒绝尝鱼如同戒绝动物食物。此外，\Person{埃塔利德斯}和\Person{赫尔墨提穆斯}把豆子列为日常餐食，而\Person{毕达哥拉斯}却教导他的门徒甚至不要穿过种有豆子的土地。那么我问，同样的灵魂是如何被恢复的，这些灵魂无法通过性情、习惯或生活方式提供任何身份证明？而现在，毕竟，（我们发现）在\Place{希腊}的众多人群中，只有四个灵魂被提及复生。但如果我们仅仅限于\Place{希腊}，好像没有灵魂转世和身体复元发生，而且\Emph{那}每天都在每个国家、每个年龄、阶层和性别中发生，那么为什么唯独\Person{毕达哥拉斯}经历这些变成另一个人的变化呢？为什么我不应该也经历它们呢？或者如果这是一种被哲学家垄断的特权——而且只是希腊哲学家，好像\Place{斯基泰人}和\Place{印度人}没有哲学家——那么为什么\Person{伊壁鸠鲁}没有回忆起他曾是另一个人，\Person{克吕西波}、\Person{芝诺}，甚至\Person{柏拉图}本人也没有，我们或许会因其甜美口才而猜想他是\Person{涅斯托耳}呢？\par

% structure: p0070 source=h2 target=section reason=relative-heading-rank; base=h2

\section{第32章 \Person{恩培多克勒}通过发展人死后变成各种动物的理论，进一步增加了\Person{毕达哥拉斯}的荒谬性。}

但事实上，恩培多克勒（\Person{Empedocles}）曾梦见自己是一位神祇，我猜想正因如此，他不屑于让人认为他从前仅仅是一位英雄，于是直言不讳地宣称：\Quote{我曾是塔姆努斯（\Person{Thamnus}），也曾是一条鱼。}他为何不说自己是甜瓜——鉴于他是如此愚蠢——或是变色龙——鉴于他如此吹嘘膨胀？无疑，他作为一条鱼（而且是一条怪鱼！）逃脱了某个无名坟墓的腐坏，因为他宁愿跳进埃特纳（\Place{Ætna}）被烤熟；此后，他的\OriginalTerm{转生}{\GreekText{μετενσωμάτωσις}}——即进入另一个身体——便永远终结了（如今只适合）作为烤肉后的一道清淡菜肴。因此，在这一点上，我们也必须驳斥那种更为荒谬的假设：在转生过程中，野兽由人而来，人亦由野兽而来。姑且放下（恩培多克勒的）塔姆努斯们（\Person{Thamnuses}）不提。我们顺带略加提及便已足够：（若详加讨论，会妨碍我们，）恐怕我们不得不诉诸嘲笑和哄笑，而非严肃的教导。我们的立场是：人的灵魂无论如何绝不能转移到野兽身上，即便哲学家认为野兽是从元素物质中产生的。现在，我们假定灵魂或是火，或是水，或是血，或是气息，或是空气，或是光；我们必须记住，各类动物都具有与相应元素相反的特性。有些冷血动物与火相对——水蛇、蜥蜴、火蜥蜴，以及一切从相克元素——水——中产生之物。同样，那些本性干燥无汁的生物与水相对；蝗虫、蝴蝶、变色龙确实喜欢干旱。凡没有紫色血色的生物（如蜗牛、蠕虫和大多数鱼类）则与血相对。凡似乎没有呼吸、没有肺和气管的生物（如蚊虫、蚂蚁、蛾子及此类微小之物）则与气息相对。凡总是生活在地下和水中、从不吸入空气的生物（你对它们的存在比对它们的名字更熟悉）则与空气相对。凡完全失明或只有黑暗视觉的生物（如鼹鼠、蝙蝠和猫头鹰）则与光相对。我举这些例子，是为了通过清晰可感的性质来阐明我的主题。然而，即使我能用手抓住伊壁鸠鲁（\Person{Epicurus}）的\Quote{原子}，或我的眼睛能看见毕达哥拉斯（\Person{Pythagoras}）的\Quote{数}，或我的脚能绊到柏拉图（\Person{Plato}）的\Quote{理型}，或我能抓住亚里士多德（\Person{Aristotle}）的\Quote{隐德莱希}，那么很可能，即使在（这些不可感知的）类型中，我也会找到那样一些动物，我必须根据它们的相反性将它们相互对立。因为我认为，无论人的灵魂由上述哪一种性质构成，它都不可能进入与每种单独性质如此相反的动物，并藉其进入赋予那些生物以起源；相反，由于我们假定灵魂所具有的那种原始相反性，它反而会被它们排斥和拒绝，而非被接纳和接受；这种原始相反性使接纳它的身体物质陷入无休止的冲突；随后又有来自每种单独性质发展所致的后续相反性。须知，人的灵魂是在完全不同的条件下，（在个别身体中）被赋予了居所、滋养、秩序、感觉、情感、性交和生育子女；此外，（在个别身体中，灵魂还获得了特殊的）性情，以及应尽的义务、好恶、恶习、欲望、快乐、疾病、治疗——简言之，它自己的生存方式和死亡出口。那么，那个（人的）灵魂——它紧贴地面，没有惊惶便不能仰望任何高度或深度，登上几级台阶便感到疲惫，潜入鱼塘便窒息——（我说，被这类软弱所困扰的灵魂）又怎能将来某时升入高空作鹰，或潜入海中作鳗？再者，它习惯于精美、细腻而珍贵的食物之后，又岂能刻意以——我不说豆荚——而是以荆棘、苦叶的野生食物、粪堆上的野兽和毒虫为生，倘若它必须转入一只山羊或一只鹌鹑？——不，或许还要在熊或狮子里吃腐肉，甚至人的尸体？然而，当它忆起自己的（本性和尊严）时，又怎能降格至此？同样，你可以将其他所有事例都提交给这个不协调的标准，从而省却我们一一详细审视。现在，无论人的灵魂尺度如何，方式如何，（我们被迫提出的问题是：）它在更大的动物或极小的动物中会怎样？每个个体身体无论大小，都必须被灵魂充满，而灵魂全身又为身体所覆盖。因此，人的灵魂如何能充满一头大象？又如何能收缩到一只蚊子之中？若它如此极度地扩展或收缩，无疑将面临危险。这又引出了另一个问题：如果灵魂完全无法转入那些不适合接受它的动物——无论因其身体习性还是其他存在法则——那么它是否会根据各种动物的特性发生变化，并适应它们的生活，尽管这与人的生活相反——事实上，由于彻底的变化，灵魂已变得与自己相反？然而，真相是：如果灵魂经历这样的转化，并失去它曾经所是，那么人的灵魂便不再是它从前之所是；若它不再是原来的自己，那么\Emph{转生}（即适应另一个身体）便化为乌有，当然不能归于这个将因彻底变化而不再存在的灵魂。因为只有当灵魂在不改变其（原始）状态的情况下经历这一过程时，才能说它经历了这种\Emph{转生}。既然灵魂不允许改变，以免失去同一性；却又无法在原始状态中保持不变，因为它无法接受相反的身体——我仍然想知道，有什么可信的理由能证明我们所讨论的这种转化。因为尽管有些人因其品性、性情和追求而被比作野兽（因为天主曾说过：\BibleQuote{人如同牲畜一样死亡}），但这并不因此意味着贪婪的人变成鸢，淫荡的人变成狗，暴怒的人变成豹，善良的人变成羊，多言的人变成燕子，贞洁的人变成鸽子，仿佛灵魂的同一实体处处在（它所转入的）动物特性中重复自身的本性。此外，实体是一回事，实体的本性是另一回事；因为实体是某一特定事物的专有属性，而本性则可能属于许多事物。试举一二例。石头或铁是实体：石头的硬度和铁的硬度是实体的本性。它们的硬度通过共同性质将物体联合起来；它们的实体则使它们彼此分离。再者，羊毛中有柔软，羽毛中也有柔软：它们的自然性质相似，（使它们等同）；它们的实体性质不相似，（使它们区别）。同样，若一个人被称为野兽或温顺的动物，这并不意味着灵魂的同一性。因为即便实体的不相似性最为显著，本性的相似性仍然被观察到：事实上，正是因你判断一个人像一头野兽，你才承认他们的灵魂并不相同；因为你只说他们\Emph{相似}，而非\Emph{相同}。这也是（我们刚引用的）天主圣言的含义：它把人的本性比作野兽，而非实体。此外，倘若天主知道人只是兽性实体，祂便不会对人作出这样的评论。\par

% structure: p0072 source=h2 target=section reason=relative-heading-rank; base=h2

\section{第33章 以嘲笑驳斥这些迁移的司法报应}

既然此种学说甚至基于司法报应的原则也被辩护，借口人的灵魂与适合其生活和功过的动物种类结合——仿佛灵魂按其不同性格，要么在注定被处决的罪犯中被杀，要么在仆役中被贬为苦力，要么在劳动者中被劳累疲惫，要么在不洁之物中受玷污；或者，同样地，被保留给最卓越的阶层和德行、有用和温柔情感中的人物，以获得尊荣、爱、关怀和殷勤的照顾——在此我还必须指出，如果灵魂经历变形，它们实际上将无法完成和体验它们应得的命运；司法报应的目的和意图将被废除，因为将缺乏功过和报应的感觉与意识。而这种意识必然缺乏，如果灵魂失去其状态；这种失去必然发生，如果它们不恒常不变。但即使它们有足够的恒常性，能保持不变直到审判——这一点埃及的梅库里乌斯（\Person{Mercurius Ægyptius}）曾承认，他说灵魂在脱离身体后，不会消散回宇宙灵魂中，而是永久保留其独特的个体性，\Quote{以便它能够}用他自己的话说，\Quote{向圣父呈报它在身体中所做的事}——（即使如此，我仍要说）我仍想审视这所谓天主审判的公义、严肃、威严和尊荣，看看人间的审判是否在其中占据了过高的宝座——在惩罚和赏报两方面都被夸大，施报时过于严厉，施恩时过于慷慨。你认为凶手的灵魂会怎样？我想，它将进入某些注定被宰杀和卖肉的牲畜中，以便它自己被杀，如同它曾杀人；自己被剥皮，如同它曾剥他人的皮；自己被当作食物，如同它曾将不幸的受害者丢给野兽，这些受害者曾被他杀害于森林和偏僻道路上。如果这就是它要接受的司法报应，这样的灵魂难道不会在事实中得到安慰而非惩罚吗？它接受\Emph{慈悲一击}来自最熟练的行刑者之手——以阿皮奇乌斯（\Person{Apicius}）或卢尔科（\Person{Lurco}）最辛辣的调味方式被埋葬，被引入你们精致的西塞罗（\Person{Ciceros}）们的餐桌，被端上苏拉（\Person{Sylla}）最辉煌的菜肴，在宴席中找到自己的葬仪，被与自己同等的体面之口吞食，而非被鸢与狼吞食，以便所有人都看到它如何以人的身体为坟墓，并在返回其同族后复活——如果它曾经历人间的审判，便在人间的审判面前得意扬扬？因为那些野蛮的死刑判决将犯有谋杀罪的罪犯，甚至在其活着时，就交给各种野兽，这些野兽被选择并训练得违背其本性来执行可怕的职分；甚至通过拖延最后时刻以加重惩罚的计谋，阻止其轻易死亡。但即使他的灵魂在剑的最后打击之前已先离去，他的身体至少不能逃脱刀剑：通过刺穿他的喉咙和胃，并刺透他的肋旁，对他自己罪过的报应仍被索取。之后他被投入火中，使他的坟墓也被欺骗。因为绝不允许他以其他方式被埋葬。毕竟，对他的火葬堆没有太多关注，以至于其他动物会找到他的遗骸。无论如何，对他的骨头没有怜悯，对他的骨灰没有宽容，它们必须受到暴露和赤裸的惩罚。人间对杀人的报复，如同自然施加的报复一样大。谁不会偏爱这世界的公义呢？正如宗徒亲自作证，它\BibleQuote{不是凭空佩剑}，\BibleRef{罗 13:4}，并且当它严厉报复以捍卫人的生命时，它是一种宗教制度。当我们思考其他罪行所受的惩罚——绞刑架、火焚、布袋、鱼叉、悬崖——谁会不认为在毕达哥拉斯（\Person{Pythagoras}）和恩培多克勒（\Person{Empedocles}）的法庭上接受判决更好呢？因为即使那些被他们送入驴和骡的体内，通过苦役和奴隶制受惩罚的可怜人，当他们回忆起矿山、囚犯队伍、公共工程，甚至监狱和黑牢（这些在无所事事的例行公事中可怕），他们又将如何为磨坊和水车的温和劳动而庆幸呢？然后，对于那些在正直生活之后将生命交还给审判者的人，我同样寻找赏报，但我反而发现惩罚。当然，好人不论被恢复成什么动物，都是一种漂亮的收获！荷马——恩尼乌斯（\Person{Ennius}）曾如此梦到——回忆起他曾是一只孔雀；但我本人不能相信诗人，即使他们完全清醒。无疑，孔雀是一种非常美丽的鸟，随意以华丽的羽毛炫耀自己；但它的翅膀不能弥补其声音的刺耳和不悦；而诗人最喜欢的就是好歌。因此，荷马变成孔雀是一种惩罚，而非荣誉。世界的报酬会给他带来更大的喜悦，当它赞扬他为文科之父时；他将更喜欢其名誉的装饰，而非其尾羽的魅力！但没关系！让诗人变成孔雀，或者如果你愿意，变成天鹅，尤其因为天鹅有可尊敬的声音：你将把那位义人埃阿科斯（\Person{Æacus}）置于何种动物中？你将用什么兽皮包裹贞洁而卓越的狄多（\Person{Dido}）？什么鸟将分配给忍耐？什么动物分配给圣洁？什么鱼分配给天真？如今，所有受造物都是人的仆役；都是他的臣民，都是他的依附者。如果不久他将成为这些受造物之一，他便因这种变化而卑贱和降级，他本来因其德行，自由获得雕像、塑像和称号作为公共荣誉和显赫特权，元老院和人民甚至投票给他祭献！哦，这是什么样的审判，神祇在死后作为对人的报酬而颁布！它们比任何人的审判更虚假；作为惩罚是可鄙的，作为赏报是可厌的；最坏的人从不害怕，最好的人也从不渴望；确实，罪犯会希望得到它们而非圣徒——前者为了更快躲避世界的严厉判决，后者为了更迟遭受它。哦，哲学家们，你们如何美好地教导我们，如何有益地劝告我们，死后惩罚和赏报变得更轻！然而，如果灵魂有任何审判等待它们，则更应设想它在生命终结时比在生命过程中更重，因为没有什么比最后到来的更完全——此外，没有什么比特别属于神的更完全。因此，天主的审判将更充分和完全，因为它将在最后被宣布，在一个永恒不可撤销的判决中，包括惩罚和安慰，（人的）灵魂不是要转生为兽，而是要返回它们自己本来的身体。所有这一切一次而永远，并在\BibleQuote{那日子，也只有圣父知道的那日子}，\BibleRef{谷 13:32}（只有他知道），以便信德通过战栗的期待，充分考验其焦虑的真诚，始终凝视那日子，在永恒的无知中，每日害怕那每日又盼望的事。\par

% structure: p0074 source=h2 target=section reason=relative-heading-rank; base=h2

\section{第三十四章 这些游移不定的想法激发了某些亵渎的基督教败坏。谴责西满（魔术师）的亵渎。}

的确，迄今还没有任何异端的面具下爆发的教义，包含如此荒唐的虚构，即人的灵魂进入野兽的身体；但我们认为有必要攻击和驳斥这一概念，作为前述观点的合理延续，以便孔雀中的\Person{荷马}能像欧福尔波中的\Person{毕达哥拉斯}一样被彻底摆脱；并且，通过同一打击摧毁\OriginalTerm{灵魂转世}{metempsychosis}和\OriginalTerm{身体转换}{metensomatosis}，从而铲除为我们的异端者提供了不小支持的基础。在\BibleBook{宗徒}{大事录}中，有（臭名昭著的）撒玛黎雅的\Person{西满}，他曾想用钱购买\Emph{圣神}：在被\Emph{祂}谴责并徒然懊悔他和他的钱财一同灭亡之后，他把精力用于摧毁真理，好像以报复来安慰自己。除了他自己的魔术所提供的支持外，他还诉诸欺骗，用他曾经为圣神提供的那笔钱，从一个妓院买了一个名叫\Person{海伦}的提洛女子——一桩配得上这可怜人的交易。他竟假称自己是至高圣父，并且进一步假装那女人是他自己的原初构想，他曾意图用这构想创造天使和总领天使；在她拥有这意图之后，她从圣父那里跃出，下降到下层空间，在那里预见了圣父的计划，产生了天使权势，这些权势对圣父、这世界的创造者一无所知；她因一种叛逆的动机（与她自己很像）被这些权势囚禁，唯恐她离开后，他们显得是另一个存在的后裔；因此，她遭受了各种侮辱，为了防止她在受辱后离开他们，她被贬低到人的形状，好像被限制在肉体的束缚中。在许多世代里，她以各种女性形象辗转，成了臭名昭著的\Person{海伦}，她对\Person{普里阿摩斯}造成了毁灭性的影响，后来又对\Person{斯特西科罗斯}的眼睛造成了损害，她因他的讽刺诗而弄瞎了他的眼睛，后又因他的颂词而恢复了他的视力。这样从一个身体流浪到另一个身体，她最终在耻辱中变成了一个更卑贱的海伦，成了一个职业妓女。因此，这个婢女就是那亡羊，至高圣父，即\Person{西满}，降到了她身上，他找回她并带她回来——不知是扛在肩上还是腰间——然后眷顾人的救恩，为了满足他的怨恨而将他们从天使权势中解放。此外，为了欺骗这些权势，他自己也取了一个可见的形像；在人间假装成人的外貌，他在\Place{犹达}扮演圣子，在\Place{撒玛黎雅}扮演圣父。啊，不幸的\Person{海伦}，你在诗人和异端者之间命运多舛，他们时而用通奸，时而用卖淫玷污你的名声！唯有将她从\Place{特洛伊}救出，比从妓院解救出来更为光荣。有一千艘船将她从\Place{特洛伊}运走；大概一千个银币就足以将她从妓院赎出。呸，\Person{西满}，你寻她如此迟缓，赎她如此反复无常！与\Person{墨涅拉俄斯}何其不同！他一失去她，就去追她；她刚被掳走，他就开始搜寻；经过十年争战，他勇敢地救回了她：没有潜伏，没有欺骗，没有诡辩。我实在担心，他是一位更好的\Quote{圣父}，为了寻回他的海伦，他更加警觉、勇敢、坚忍地劳苦。\par

% structure: p0076 source=h2 target=section reason=relative-heading-rank; base=h2

\section{第三十五章 卡尔波克拉特的意见，毕达哥拉斯教义的另一个旁支，陈述与驳斥}

然而，这故事并不只是为了你（\Person{西满}）一人而编造的。\Person{加波喀特}也一样利用它，他像你一样是个术士和行淫者，只是他没有一个海伦。他为什么不该有呢？因为他断言灵魂重新穿戴身体，以确保以一切手段推翻神性和人性的真理。因为，依照他那可悲的教义，对任何人来说，此生不会圆满，除非所有被认为玷污它的瑕疵都在其行为中完全展现出来；因为没有什么在本性上被视为恶，仅仅在于人们如何看待它。因此，他认为人的灵魂转生到任何异质的身体中是绝对必要的，每当在生命旅程的早期阶段任何邪恶尚未被充分犯下时。恶行（可以肯定）属于生命。此外，每当灵魂在罪中有所亏损，如同债务人，它就必须被重新召回到存在中，直到它\BibleQuote{交还最后一分钱}（见：\BibleRef{玛 5:26}），一次又一次地被驱逐到身体的监狱中。为此，他篡改了主的那整个比喻，其含义极其清晰明了，并且从一开始就应当按其平实自然的含义来理解。因此，我们的\Quote{仇敌}（那里提到的）是外邦人，他与我们行走在同一条生命之路上，这条路对他和我们都是共同的。现在，我们\BibleQuote{必须离开这个世界}（见：\BibleRef{格前 5:10}），如果不允许我们与他们交往的话。因此，祂命令我们对这样的人表现出友善的态度。\BibleQuote{爱你们的仇敌}，祂说，\BibleQuote{为那诅咒你们的祈祷}（见：\BibleRef{路 6:27}），免得这样的人在任何交易中被你的不义行为激怒，将你交给本族的法官（见：\BibleRef{玛 5:25}），你被投进监狱，被拘禁在狭小的牢房里，直到你还清他的一切债务。再者，如果你倾向于将\Quote{仇敌}一词应用于魔鬼，主（的诫命）劝你，\Quote{当你与他行在路上时}，要与他达成一个与你真信德的要求相一致的契约。现在，你关于他所立的契约就是弃绝他、他的浮华和它的天使。这就是你在此事上的协议。现在，你必须从遵守契约出发来实行这一友好谅解：你绝不能再考虑收回任何你已经宣誓放弃并归还给他的东西，免得他作为一个欺诈者和违约者，在审判者天主面前控告你（因为在另一处经文中，我们读到他是\BibleQuote{控告兄弟的}（见：\BibleRef{默 12:10}），或控告圣人，那里指的是实际的法律诉讼行为）；免得这位审判者将你交给执行判决的天使，而\Emph{他}将你投入地狱的监狱，在那里直到复活之前，即使是你最小的过犯也还清之后，才得释放。有什么比这更合适的含义呢？有什么更真实的解释呢？然而，如果按照\Person{加波喀特}的说法，灵魂被束缚于各种罪行和恶行，那么根据他的体系，我们必须将其\Quote{仇敌}和敌人理解为那个更好的心智，它将强迫它履行某种美德行为，使它从一个身体被驱逐到另一个身体，直到在任何一个身体中都不再欠美德要求之债。这意味着，好树由其坏果子而被认识——换言之，真理的教义从最坏的诫命中被理解。我明白，这一学派的异端者特别贪婪地抓住\Person{厄里亚}的例子，他们认为厄里亚在\Person{若翰}（洗者）身上如此再造，以至于我主的话成了他们转生理论的依据，祂说：\BibleQuote{厄里亚已经来了，而他们却不认识他}（见：\BibleRef{玛 17:12}）；又在另一处：\BibleQuote{如果你们愿意接受，他就是那要来的厄里亚}（见：\BibleRef{玛 11:14}）。那么，犹太人走近\Person{若翰}问他：\BibleQuote{你是厄里亚吗？}（见：\BibleRef{若 1:21}）真的出于毕达哥拉斯式的含义，而不是出于神性预言的\BibleQuote{看，我要派遣提市贝人\Person{厄里亚}到你们这里来}（见：\BibleRef{拉 4:5}）吗？然而事实是，他们的灵魂转世或转生理论，意指召回早已死去的灵魂，并将其送回另一个身体。但是\Person{厄里亚}要再来的，不是在离开生命（以死亡的方式）之后，而是在他被提升（未经死亡而被转移）之后；不是为了恢复他并未脱离的身体，而是为了重访他被提升离去的世界；不是以重拾他已放下的生命的方式，而是以应验预言的方式——真正同样的人，无论就他的名字和称号，还是就他未改变的人性而言。因此，\Person{若翰}如何能是\Person{厄里亚}呢？你在天使的宣告中得到答案：\Quote{他要行在百姓面前，}他说，\Quote{在\Person{厄里亚}的精神和能力中}——请注意，不是在他的灵魂和身体中。这些实体事实上是每个人的自然财产；而精神和能力是由天主的恩宠作为外在恩赐赋予的，因此可以按照全能者的意图和旨意转给另一个人，正如古时\Person{梅瑟}精神的情况（见：\BibleRef{户 12:2}）。\par

% structure: p0078 source=h2 target=section reason=relative-heading-rank; base=h2

\section{第36章 我们作者的主题要点：论人类的两性}

为了讨论这些问题，我们曾离开（如果我没有记错）一处基点，现在必须回到那里。我们已经确立了这样的立场：灵魂是藉着人的中介，以种子的方式被置于人内，并且它的种子从起初就是同一的，如同灵魂的种子对于人类也是同一的；（我们这样确定）是由于哲学家和异端者们的对立意见，以及\Person{柏拉图}所提及的那句古老格言（上面我们曾提到过）。我们现在按顺序探讨从这些前提推出的论点。灵魂与身体同时被播种在子宫中，也同时接受其性别；并且这确实如此同步，以致两种实体都不能单独被视为性别的成因。现在，如果在这些实体的播种中，在受孕时允许有任何时间间隔，以致要么肉身、要么灵魂首先被孕育，那么由于受孕时间的差异，人们就可以将特定的性别归之于其中一个实体，这样要么肉身将其性别印在灵魂上，要么灵魂将其性别印在肉身上；正如\Person{阿佩莱斯}（那位异端者，不是画家）根据\Person{斐鲁美纳}的教导，认为男女的灵魂优先于他们的身体，因此使作为较晚者的肉身从灵魂接受其性别。那些主张灵魂在出生后才附加于肉身的人，当然预先规定了早已形成的灵魂的性别为男性或女性，与肉身的性别一致。但事实是，两种实体的播种在时间上是不可分的，它们的流溢也是同一的，因此它们获得了性别的共同性；以致自然的过程，无论它是什么，将为（不同的性别）划定界线。当然，从这个观点我们得到了最初两次形成方式的证明：当时男性以更完善的方式被塑造和调适，因为\Person{亚当}是先被形成的；女人远远落后于他，因为\Person{厄娃}是后被形成的。因此她的肉身在很长时间内没有特定的形状（正如后来她从\Person{亚当}肋旁被取出来时所具有的形状）；但即使那时她本身也是一个活生生的存在，因为我认为她在灵魂方面甚至当时就是\Person{亚当}的一部分。此外，天主的\Emph{\OriginalTerm{吹气}{afflatus}}本也可以使她拥有生命，若不是在女人身上已经有一种从\Person{亚当}传来的他的灵魂和肉身的话。\par

% structure: p0080 source=h2 target=section reason=relative-heading-rank; base=h2

\section{第三十七章 论胚胎的形成和状态。其与本论文主题的关系}

如今，在母胎中播种、形成并完成人类胚胎的整个过程，无疑由某种力量所掌管，这力量在此服务于天主的旨意，无论它被指定采用何种方法。就连罗马的迷信，也通过细心关注这些细节，设想出女神\OriginalTerm{阿莱摩娜}{Alemona}滋养子宫中的胎儿；以及\OriginalTerm{诺娜}{Nona}和\OriginalTerm{德奇玛}{Decima}，以妊娠最关键的月份命名；还有\OriginalTerm{帕图拉}{Partula}管理和引导分娩；以及\OriginalTerm{卢奇娜}{Lucina}将孩子带到出生和光明之中。而我们则相信，天使在此为天主司职。因此，胚胎一旦形体完成，就在母胎中成为人了。的确，梅瑟法律对导致流产者处以应有的惩罚，因为那里已存在人的雏形，它甚至已经被赋予了生与死的状态，因为它已面临两者的结局；尽管如此，由于它仍活在母腹中，它在很大程度上与母亲共享其状态。我也必须谈谈灵魂诞生的时期，以免遗漏整个过程中任何附带的情况。一般而言，成熟且正常的出生发生在第十个月初。那些对数进行理论思考的人，尊崇十这个数为其他数字之母，并赋予人类诞生以完美。就我个人而言，我更愿意从天主的角度来看待这个时间度量，似乎这十个月更是在引导人进入十诫；因此，完成我们自然诞生所需时间的数字估算，应当与我们重生生活之规则的数字分类相对应。但由于出生也在第七个月完成，我更容易在这个数字中（而非第八个月）认识到与安息时期数字相符的荣誉；这样，有时人类出生中产生天主肖像的月份，其数字就与天主创造完成并被祝圣的日子相符。人类诞生有时被允许是早产的，却仍然与一个\Emph{hebdomad}（七之数）适切完美地一致，作为我们复活、安息和国度的预兆。因此，\Emph{ogdoad}（即八之数）与我们的形成无关；因为在它所代表的时间里，将不再有婚姻。我们已经证明了身体与灵魂的结合，从它们最初种子的凝结直到\Emph{胎儿}的完全形成。我们现在也主张它们从出生起就结合；首先，因为它们共同成长，只是各自以适合其本性的不同方式——肉体在形体上成长，灵魂在智力上成长；肉体在物质状态上，灵魂在感觉上。然而，我们不可认为灵魂的实体有所增长，以免它也被说成实体能够减少，甚至可能被认为会消亡；但其固有的能力（其中包含所有自然特性，原本就被植入其存在中）随着肉体逐渐发展，而不会损伤其最初吹入人内时所获得的实体的胚基。取一定量的金子或银子——还是一块粗糙的团块：它确实具有紧实的状态，而且此刻比将来更为压缩；然而在其外形之内，自始至终是一块金子或银子的团块。当这块团块后来被打成金叶或银叶而延展时，它由于原团块的伸长而变得比以前更大，但并非由于任何添加，因为它是在空间中延展，而非在体积上增加；尽管从某种意义上说，它延展时甚至也增加了：因为它可能在形式上增加，而非状态上。再者，金子或银子的光泽，当金属成块时无疑真实地内在于其中，但还只是隐晦的，现在则发出耀眼的光辉。随后，根据材料可塑性的不同，形状发生各种改变，使材料顺从工匠的加工；但工匠除了赋予团块形状之外，并未在其状态上添加任何东西。同样，灵魂的成长和发展应被理解为，并非扩大其实体，而是唤起其能力。\par

% structure: p0082 source=h2 target=section reason=relative-heading-rank; base=h2

\section{第三十八章 论灵魂的成长。其成熟与人的肉体成熟同时发生。}

我们已经确立了这样的原则：灵魂所有与感觉和理智相关的自然属性都内在于其自身实体，源于其天生构造，但它们在生命的各个阶段逐渐生长，并通过偶然的环境以不同方式发展，取决于人的手段和技艺、风俗习惯、地方境况以及至高能力的影响；但现在我们要考虑灵魂与身体结合的这一方面，我们主张灵魂的\Emph{青春期}与身体的青春期同时发生，二者大约在生命的第十四年一同达到完全成熟——大体而言——前者通过感官的提示，后者通过身体器官的生长；（我们确定这个年龄）并不像\Person{阿斯克勒庇亚得}所假设的那样是因为反思从此开始，也不是因为民法将人生真正事务的起始定在这个时期，而是因为这是从起初就设定的秩序。因为正如\Person{亚当}和\Person{厄娃}在知道善恶之后感到必须遮盖自己的赤身，我们宣称从我们体验到同样的羞耻感之时起，就具有同样的善恶辨识力。从此前提到的年龄（十四岁）起，性别便被一种特殊的感觉所弥漫和包裹，\OriginalTerm{贪欲}{concupiscence}利用眼睛的服务，将它的快乐传递给他人，理解男女之间的自然关系，穿上无花果树的围裙来遮盖它仍引起的羞耻，将人逐出纯真和贞洁的乐园，并在它狂野的\OriginalTerm{淫欲骚动}{pruriency}中坠入罪恶和反常的犯罪刺激；因为它的冲动此时已超越了自然的命定，源于对自然的恶用。但严格自然的贪欲仅仅限于对上帝起初赐予人的那些\OriginalTerm{食物}{aliments}的欲望。\BibleQuote{园中各树上的果子}祂说，\BibleQuote{你可以随意吃；}\BibleRef{创 2:16}然后再次对洪水之后的下一个世代扩大了赐予：\BibleQuote{所有活的动物都可作你们的食物；看，我将这一切都赐给你们，如同蔬菜一样；}\BibleRef{创 9:3}——这里祂更关注身体而非灵魂，虽然也对灵魂有益。因为我们应当消除\OriginalTerm{吹毛求疵者}{caviller}的一切借口，他因灵魂显然需要\OriginalTerm{食物}{ailments}，就坚持据此认为灵魂是可死的，因为它靠饮食维持，当它们被断绝时，过一段时间就失去活力，最终完全断绝时就会衰萎而死。现在我们须注重的要点不仅是哪个特定官能欲望这些（\OriginalTerm{食物}{aliments}），而且还有为何目的；即使是为其自身，仍要问：这种欲望为何、何时感到、持续多久？然后再考虑：出于本能欲望是一回事，出于必要欲望是另一回事；作为存在的属性欲望是一回事，为特定对象欲望是另一回事。所以，灵魂会欲望饮食——确实是为了自身，出于一种特殊的必要；但为了肉体，则出于其属性的本性。因为肉体无疑是灵魂的房屋，而灵魂是肉体的暂居者。因此，居住者的欲望将出于暂时的原因和其称谓所暗示的特殊必要——目的是为了在他暂居期间获益并改善其临时住所的条件；当然，并不是为了自己成为房屋的地基、墙壁、支撑或屋顶，而是仅仅为了得到住所和安居，因为除非是完好且建造良好的房屋，他无法得到这样的安居。（现在，将这些意象应用于灵魂，）如果它没有这样的住所，它将无法离开其居所，并且由于缺乏适当合宜的资源，无法安然无恙地离去，还拥有其自身的支撑以及属于其自身状态的\OriginalTerm{食物}{aliments}——即不死、理性、感觉、理智和自由意志。\par

% structure: p0084 source=h2 target=section reason=relative-heading-rank; base=h2

\section{第三十九章 恶神从出生起就玷污了灵魂的纯洁}

灵魂在出生时所领受的这些所有禀赋，仍然被起初嫉妒地看着它们的邪恶者所遮蔽和败坏，以至于它们从未在自发的行动中显现，也未能得到应有的管理。因为哪个人类个体不会被恶灵所依附，随时准备从他们出生的门坎就攫取他们的灵魂，而恶灵正是通过伴随生育的所有迷信过程被邀请来出席的？因此，所有人生来都有拜偶像的接生婆，而怀胎的子宫本身，仍绑着在偶像前编织的束带，宣布其子女被奉献给魔鬼：因为在分娩时，他们祈求\OriginalTerm{卢喀娜}{Lucina}和\OriginalTerm{狄安娜}{Diana}的帮助；整整一周，在\OriginalTerm{朱诺}{Juno}的荣誉下摆设筵席；最后一天，祈求命星的命运；然后婴儿踏上地面的第一步是献给女神\OriginalTerm{斯塔提娜}{Statina}。这之后，难道有人不把他儿子的整个头部献于偶像崇拜之役，或拔下一根头发，或剃光全部，或绑起来献祭，或封为圣用——为了家族、祖先或公共虔诚吗？正是基于这种早期占有的原则，\Person{苏格拉底}在孩提时就被魔鬼之灵找到。同样，每个人都被分配了他们的\Emph{\OriginalTerm{守护神}{genii}}，这只是\Emph{\OriginalTerm{魔鬼}{demons}}的另一个名称。因此，没有一个出生（当然，我指的是外邦人）是纯洁的，没有偶像崇拜的迷信。正是由于这种情况，宗徒说，如果父母中有一方成圣，子女就是圣洁的；\BibleRef{格前 7:14} 这既出于（基督徒）种子的特权，也出于制度（通过圣洗和基督徒教育）的规训。\Quote{否则，}他说，\Quote{儿女就不洁净}因出生：\BibleRef{格前 7:14} 仿佛他要我们明白，信徒的子女是被定为圣洁，从而得救；以便他通过这种希望的保证来支持婚姻，他决心保持婚姻的完整性。此外，他当然没有忘记主如此明确地宣告的话：\Quote{除非一个人由水和圣神而生，否则不能进入天主的国；}\BibleRef{若 3:5} 换言之，他不能成为圣洁。\par

% structure: p0086 source=h2 target=section reason=relative-heading-rank; base=h2

\section{第四十章 人的身体在作恶时只是灵魂的辅助}

每一个灵魂，因其出生，在亚当内有其本性，直到它在基督内重生；此外，只要它没有这重生，就是不洁的；\BibleRef{罗 6:4} 因为不洁，所以活跃地犯罪，甚至（因它们的结合）用自己的羞耻沾染了肉体。虽然肉体是有罪的，我们被禁止随从它而行，\BibleRef{迦 5:16} 它的作为被谴责为与圣神争战，人因它被斥为属肉体的，\BibleRef{罗 8:5} 但肉体本身并没有这样的耻辱。因为它本身并不会思考或感觉任何事以建议或命令犯罪。它怎么能呢？它只是一个服务的东西，它的服务不像仆人或密友——有生命的人；而是像器皿或类似的东西：它是肉体，不是灵魂。杯子可以向口渴的人提供服务；然而，如果口渴的人不把杯子送到嘴边，杯子就不会提供任何服务。因此，人的区别（\OriginalTerm{区别}{differentia}）或特性绝不是在于他的土性元素；肉体也不是人本身，如同灵魂的某种官能和人格特质；而是一种完全不同实体和不同状况的东西，虽然被附加给灵魂，作为财产或生活职事的工具。所以，圣经责备肉体，因为在贪欲、食欲、酗酒、残忍、拜偶像和其他肉体的作为等运作中，没有肉体，灵魂就什么也做不了——我指的是那些不仅限于感觉，而且产生效果的运作。诚然，当罪的冲动不产生效果时，通常归咎于灵魂：\Quote{凡注视妇女而贪恋她的，已在她心里犯了奸淫。}\BibleRef{玛 5:28} 但肉体单独，没有灵魂，在美德、义德、忍耐或贞洁等运作中，又曾做过什么呢？然而，将罪和罪行归给某一实体，而你不给它任何自己的善行或品行，这是多么荒谬！现在，协助犯罪的同伙被带到审判前，但方式是这样的：实际犯下罪行的主犯应承担刑罚的重担，虽然从犯也不会逃脱控告。当主犯的官员因他的过错而受惩罚时，更大的厌恶落在他身上。煽动和命令犯罪的人受到更多的鞭打，而同时听从这种邪恶命令的人也不能被宣告无罪。\par

% structure: p0088 source=h2 target=section reason=relative-heading-rank; base=h2

\section{第四十一章 尽管人的灵魂因原罪而败坏，仍留下基础，使天主的恩宠能藉属灵的重生为其恢复而工作}

因此，除了因邪灵的介入而附加于灵魂的恶以外，还有一种先于它并在某种意义上自然的恶，源自灵魂的堕落根源。因为，正如我们之前所说，我们本性的败坏是另一种本性，拥有自己的神和父亲，即那败坏的作者。然而，灵魂中仍有一部分善，那是原初的、神圣的、真正的善，是其本性的固有之物。因为出自天主的，与其说是被消灭，不如说是被遮蔽。它确能被遮蔽，因为它不是天主；但它不能被消灭，因为它出自天主。因此，如同光线被不透光的物体阻挡时，虽因如此致密的物体阻隔而不显现，却仍然存在；同样，灵魂中的善被恶所压制，因恶的遮蔽性，要么完全不见光芒，要么只在偶然的出口处挣扎透出一缕微光。因此，有些人极其邪恶，有些人极其良善；但所有人的灵魂仍属同一类型：即使在最坏的人中也有某种善，在最好的人中也有某种恶。因为唯有天主是无罪的；而无罪的人只有基督，因为基督也是天主。如此，灵魂的神性因其原初的善而在预言性的预感中迸发；它意识到自己的起源，便在诸如\Quote{\Emph{善哉，天主！天主知道！}}和\Quote{\Emph{再见！}}的呼喊中为天主（其作者）作证。正如没有灵魂是无罪的，同样也没有灵魂是缺乏善的种子的。因此，当灵魂接受信仰，在第二次诞生中藉水和来自高天的能力而更新时，其先前败坏的幔子便被除去，它便以全部的光明看见那光。它（在第二次诞生中）也被圣神所接纳，正如在第一次诞生中被邪灵所拥抱。肉体如今跟从与圣神结合的灵魂，作为嫁妆的一部分——不再是灵魂的仆人，而是圣神的仆人。哦，幸福的婚姻，若其中不犯婚誓的违犯！\par

% structure: p0090 source=h2 target=section reason=relative-heading-rank; base=h2

\section{第四十二章  睡眠，死亡的镜子，作为讨论死亡的导言}

现在尚待讨论的，是死亡的主题，好让我们的论述随着灵魂的完成而终止；尽管\Person{伊壁鸠鲁}在其广为流传的学说中断言，死亡与我们无关。他说：\Quote{那分解的东西没有知觉；而没有知觉的东西与我们无关。}然而，实际遭受分解、失去知觉的并非死亡本身，而是经历死亡的人。但就连\Emph{他}也承认，承受属于那具有行动的主体。如果人注定要经历死亡，而死亡分解身体、摧毁感官，那么说如此重大的承受不属于人，是何等荒谬！\Person{塞内卡}说得更加精确：\Quote{死后一切皆终，甚至死亡本身也终结。}从他的这一立场必然得出：死亡将属于它自身，因为它自身也终结；而且更属于人，因为在人的终结中，\Quote{一切}也终结，死亡自身亦终结。伊壁鸠鲁说，死亡不属于我们；那么，照此说来，生命也不属于我们。因为，当然，如果那导致我们分解的与我们无关，那么那结合与组成我们的也必定与我们无关。如果我们的知觉的丧失与我们无关，那么知觉的获得也必定与我们无关。事实上，那摧毁灵魂本身的人（如伊壁鸠鲁所为）也不能不摧毁死亡。至于我们（作为基督徒），必须如同讨论死后生命和灵魂的某些其他领域那样讨论死亡，（假设）我们无论如何属于死亡，如果它不属于我们的话。按同样的原则，即使是睡眠——死亡的镜子——也不属于我们的主题之外。\par

% structure: p0092 source=h2 target=section reason=relative-heading-rank; base=h2

\section{第四十三章  睡眠是一种自然功能，由其他考虑和圣经的见证所显示}

因此，让我们首先讨论睡眠的问题，然后探讨灵魂是如何遭遇死亡的。睡眠当然不是一种超自然的事情，正如某些哲学家所主张的那样，他们认为睡眠是由看似超自然的原因造成的。斯多葛派断言睡眠是\Quote{感官活动的暂时中止}；伊壁鸠鲁派将其定义为动物灵的中断；阿那克萨戈拉和色诺芬尼认为是同一灵体的疲惫；恩培多克勒和帕尔梅尼德斯认为是它的冷却；斯特拉托认为是（灵魂）同质灵的分离；德谟克利特认为是灵魂的匮乏；亚里士多德认为是心脏周围热的中断。至于我自己，我可以有把握地说，我从未以这样一种方式睡眠，以至于发现这些情况中的任何一个。事实上，我们不可能相信睡眠是一种疲惫；它恰恰相反，因为它无疑会消除疲惫，人通过睡眠得到恢复，而不是感到疲劳。此外，睡眠并不总是疲劳的结果；即使如此，疲劳也不会持续下去。我也不能认为睡眠是动物热力的冷却或衰退，因为我们的身体通过睡眠获得温暖，以至于食物通过睡眠的规律分散如果因为过多的热而加速，或者因为寒冷而延迟，如果睡眠具有所称的冷却作用，就不会那么容易进行。还有进一步的事实，即出汗表明消化过热；而消化被我们描述为一种煮化的过程，这是一种涉及热而不是冷的活动。同样，灵魂的不朽排除了相信睡眠是动物灵的中断、或灵的匮乏、或（灵魂）同质灵的分离的理论。如果灵魂遭受减少或中断，它就会消亡。事实上，我们唯一的办法是同意斯多葛派，将灵魂确定为感官活动的暂时中止，只为身体提供休息，而不是也为灵魂。因为灵魂，作为永远在运动，永远在活动，从不屈服于休息——这种状态与不朽不相容：没有任何不朽的东西允许其活动终结；但睡眠是活动的终结。实际上，睡眠慈祥地将工作的停止赐予那属于有死的身体，并且仅赐予身体。因此，任何怀疑睡眠是否是自然功能的人，都会有辩证法专家质疑自然与超自然之间的整个区别——这样，他认为是超自然的事物，他就可以（如果他愿意）安全地归因于自然，而自然确实如此安排了事物，以至于它们似乎可以被视为超自然；因此，当然，所有事物都是自然的，或者没有事物是自然的（视情况而定）。然而，在我们（基督徒）看来，只有通过默观天主——我们现在讨论的所有事物的作者——所提示的才能被听取。因为我们相信，如果自然存在的话，它是天主的理性作品。现在理性主宰睡眠；因为睡眠如此适合人、如此有益、如此必要，以至于没有它，就没有灵魂能够提供恢复身体、恢复精力、确保健康、提供工作的中止和劳动的补救，以及白日离去、黑夜通过从所有物体上夺走它们的颜色而提供合法享受的机能。既然睡眠对我们的生命、健康和救助不可或缺，那么它就没有任何不合理、不自然之处。因此，医生将所有与生命、健康和对自然有益的事物相悖的东西驱逐出自然的大门；因为他们认为那些敌视睡眠的疾病——精神和胃的疾病——是违反自然的，并通过这一判断确定了其必然结论，即睡眠是完全自然的。此外，当他们宣称在昏睡状态下睡眠不自然时，他们得出的结论基于这样的事实：当睡眠处于适当和规律的行使时，它是自然的。因为每一种自然状态都会因缺陷或过度而受损，而通过适当的度量和数量得以维持。因此，那种在状态上可能是自然的，如果因缺陷或过度而变得不自然。那么，如果你从自然条件中除去饮食，又会怎样呢？因为饮食中蕴含着睡眠的主要诱因。可以肯定的是，从人的自然之初，人就受到这些（睡眠）本能的印刻。\BibleRef{创 2:21} 如果你接受天主的教导，（你会发现）人类之源亚当在品尝安息之前先尝到了困倦；他在劳作之前，甚至在他进食之前，甚至在他说话之前就睡着了；以便人们可以看到睡眠是一种自然的特征和功能，并且它实际上先于所有自然能力。从这个最初的例子，我们也引导出在睡眠中死亡的图像。因为正如亚当是基督的预像，亚当的睡眠预示了基督的死亡，基督要睡一个致命的睡眠，以便从祂肋旁所受的创伤中，同样（如同厄娃的形成）预表教会，真正的活人之母。这就是为什么睡眠如此有益、如此合理，并且实际上被塑造成那普遍和共通于人类之死的模型。天主确实愿意（可以顺便说一句，祂通常在其安排中不会没有这些预像和影子而成就任何事）以比柏拉图的例子更充分和更完全的方式，通过每日的重复，将人类状况的轮廓摆在我们面前，特别是关于其开始和终结；如此伸手帮助我们的信德更易于通过预像和比喻，不仅言语上，而且事物上。因此，祂在你眼前呈现人的身体被睡眠的友善力量击倒，被休息的友善必要性所征服，静止不动，正如它生命之前所躺卧，也正如它生命之后所躺卧：它躺在那里，作为其最初被塑造时的形态和最终被埋葬时的状况的见证——在两个阶段等待灵魂，前者在其被赋予之前，后者在其刚刚离开之后。与此同时，灵魂所处的情境使其似乎在其他地方活动，通过暂时伪装其在场来学习忍受未来的缺席。我们很快就会知道赫尔莫提摩斯的情况。但在此期间它做梦。那么它的梦从哪里来呢？事实是，它不能完全休息或无所事事，也不将不朽的本性局限于睡眠的安静时刻。它证明自己拥有持续的运动；它穿越陆地海洋，它经商，它激动，它劳作，它游戏，它悲伤，它喜乐，它追求合法和非法的事务；它显示出即使没有身体也具有多么巨大的力量，它自身配备着多么完备的器官，尽管同时暴露出它需要再次将其活动印刻在某个身体上。因此，当身体摆脱昏睡时，它通过自己恢复自然功能，在你眼前宣告死者的复活。因此，睡眠既必须有自然理性，也必须有理性本性。如果你只把它视为死亡的图像，你就开始信德，滋养望德，学习如何死和如何活，甚至在睡眠中学习警醒。\par

% structure: p0094 source=h2 target=section reason=relative-heading-rank; base=h2

\section{第四十四章 \Person{赫尔莫提摩斯}的故事与\Person{尼禄}皇帝的无眠。灵魂与身体直到死亡才分离。}

至于赫尔莫提摩斯（\Person{Hermotimus}）的情况，据说他在睡眠中灵魂会离开身体，仿佛灵魂像度假的人一样四处游荡。他的妻子背叛了这一奇特现象。他的敌人发现他熟睡后，烧毁了他的身体，如同那是一具尸体；当他的灵魂回来得太迟时，它（我想）将谋杀之罪归于自身。然而，克拉佐美奈（\Place{Clazomenæ}）的好公民们用一座庙宇安慰可怜的赫尔莫提摩斯，任何妇女都不得进入该庙，因为那个妻子的恶名。为何讲述这个故事？为了不让赫尔莫提摩斯的例子助长轻信，因为俗众常轻易认为睡眠是灵魂与身体的分离。那必定是一种更沉重的睡眠：可以设想是梦魇，或者是索拉努斯（\Person{Soranus}）针对梦魇提出的那种病态倦怠，又或者是传说加于厄庇美尼德（\Person{Epimenides}）的那种疾病，他睡了大约五十年。然而，苏埃托尼乌斯（\Person{Suetonius}）告诉我们，尼禄（\Person{Nero}）从不做梦；忒奥彭普斯（\Person{Theopompus}）对特拉叙美德斯（\Person{Thrasymedes}）也这么说；但尼禄在生命末期，在极度惊恐之后，确实勉强做了梦。如果赫尔莫提摩斯的例子被信以为真，认为他的灵魂在睡眠中处于实际怠惰状态，并与身体完全分离，那又会怎么说呢？你可以推测这绝非灵魂的放纵，以至于允许它在不死亡的情况下离开身体，且如此频繁地发生，仿佛已成为其状态和本性的习惯。如果确实有人告诉我灵魂曾发生过类似日食或月食那样的事，我确实会认为这由天主的直接干预所致，因为一个人从天主那里接受警告或惊恐的警示，如同被闪电或猝死击中，并非不合理；只是更自然的结论是，这一过程应通过梦发生，因为如果必须认为它不是梦（正如我们反对的假设所主张的那样），那么这事更应发生在人完全清醒的时候。\par

% structure: p0096 source=h2 target=section reason=relative-heading-rank; base=h2

\section{第四十五章 梦，灵魂活动的附带效果。\OriginalTerm{神魂超拔}{ecstasy}。}

我们必须在此阐述基督徒对梦的看法，认为梦是睡眠的附带现象，是灵魂不容忽视的激荡——我们已宣称灵魂因永恒运动而始终忙碌活跃，而这又证明了其神圣品质与不朽。因此，当休息降临于人的身体（这是身体特有的安慰）时，灵魂蔑视对其非自然的休息，永不静止；既然它得不到身体肢体的帮助，便使用自己的肢体。想象一名没有武器或工具的角斗士，或者一名没有马匹的驭手，但仍在比划着他们各自职业的全程动作和努力：有战斗，有挣扎；但努力是徒劳的。然而，整个过程似乎完成了，尽管实际上并未真正实现。有行动，却没有效果。我们将这种能力称为\OriginalTerm{神魂超拔}{ecstasy}，此时感性的灵魂仿佛脱离自身，甚至类似疯狂。因此，从最初开始，睡眠便由神魂超拔开启：\BibleQuote{天主使神魂超拔降在亚当身上，他就睡着了。} \BibleRef{创 2:21} 睡眠临于他的身体使之安息，但神魂超拔降于他的灵魂以消除安息；正因此，通常（从秩序中产生事情的本性）睡眠与神魂超拔相结合。事实上，在我们的梦中，我们带着真实的情感、焦虑和痛苦经历喜悦、悲伤和惊恐！而这些情感本不应触动我们，因为梦不过是纯粹的幻想，除非我们在做梦时是自主的（不受神魂超拔影响）。在这些梦中，善行无效，罪过无害；因为我们不会因幻想的罪行而被定罪，也不会因想象的殉道而加冕。但你会问：灵魂如何记得它的梦，既然据说它对自己的运作没有任何控制？这种记忆必是我们正在讨论的神魂超拔状态的特殊恩赐，因为它并非源于任何健康功能的缺失，而全然出于自然过程；它并不驱逐精神功能——只是暂时撤回。震动是一回事，移动是另一回事；毁灭是一回事，搅动是另一回事。因此，记忆所提供的是心智健全的迹象；而健全的心智在神魂超拔中所经历的（当记忆不受阻碍时）是一种疯狂。因此，在那种状态下，我们不被说成疯狂，而是做梦；如果有时如此，也完全拥有我们的心智官能。因为尽管我们运用这些官能的能力可能变得黯淡，但并未熄灭；除了在神魂超拔以特殊方式在我们内运行之时，它本身似乎缺席，以致它既将健全心智和智慧的影像呈现在我们面前，就像它呈现错乱影像一样。\par

% structure: p0098 source=h2 target=section reason=relative-heading-rank; base=h2

\section{第四十六章 梦与异象的多样性。\Person{伊壁鸠鲁}轻看梦，而多数人却极为重视。梦的实例。}

我们现在被迫就灵魂所受激发的梦之性质发表意见。我们何时才能论及死亡的主题呢？关于这个问题，我要说：当天主允许的时候；这不会拖延太久，因为它无论如何必然发生。\Person{伊壁鸠鲁}主张梦完全是虚妄之事；（但他这样说）是为了将神灵从一切忧虑中解放出来，瓦解整个世界的秩序，使万物呈现纯粹偶然的面貌——其结局是偶然的，其本质是随机的。那么，若事物本性如此，则真理也必有某种偶然性，因为它不可能成为唯一免于万物所应得的命运之物。\Person{荷马}为梦指定了两道门——\Quote{\Emph{角质的}}真理之门，\Quote{\Emph{象牙的}}谬误与欺诳之门。因为，据说透过角可以看见，而象牙则不透明。\Person{亚里士多德}在表达其意见说梦在多数情况下不真实时，仍承认其中有某种真理。忒尔墨索斯人不承认梦在任何情况下都是无意义的，但他们在无法揣测其含义时责怪自己的软弱。如今，有谁如此不谙人事，以致从未在梦中感知过某些真理呢？我只需略举几个更为显著的例子，就会使\Person{伊壁鸠鲁}脸红。\Person{希罗多德}叙述说，\Place{米底}人的国王\Person{阿斯提阿格斯}在梦中看见从他的童贞女儿子宫中涌出一股洪水，淹没了\Place{亚细亚}；又在她婚后次年，看见从她同一部位长出一株葡萄藤，覆盖了整个\Place{亚细亚}。在\Person{希罗多德}之前，\Place{兰普萨库斯}的\Person{卡戎}也讲述了同一个故事。那些解释这些异象的人并没有欺骗母亲，当他们预言她的儿子将成就如此伟业时，因为\Person{居鲁士}确实既淹没了也覆盖了\Place{亚细亚}。\Place{马其顿}的\Person{腓力}在成为父亲之前，曾看见其配偶\Person{奥林匹亚丝}的阴部印有一枚小戒指的形状，上面刻有狮子作为印章。他断定她不可能有子嗣（我想是因为狮子只做一次父亲），这时\Person{阿里斯托德摩}或\Person{阿里斯托丰}碰巧猜断，那印章并非无意义或空洞的征兆，而是预示着一个极其杰出的儿子。任何了解\Person{亚历山大}的人都会认出他就是那枚小戒指上的狮子。\Person{埃福鲁斯}如此记载。此外，\Person{赫拉克利德斯}告诉我们，\Place{希梅拉}的一位妇人在梦中看见了\Person{狄奥尼修斯}对\Place{西西里}的暴政。\Person{欧福里翁}公开记载了一个事实：在生下\Person{塞琉古}之前，其母\Person{劳迪丝}预见到他将成为\Place{亚细亚}帝国的君主。我又从\Person{斯特拉博}那里得知，\Person{米特拉达梯}之所以取得\Place{本都}，也是由于一个梦；我还从\Person{卡利斯提尼}那里得知，\Place{伊利里亚}人\Person{巴拉里里斯}从\Place{摩罗西亚}人领地扩展到\Place{马其顿}边境，也是出于梦的指示。罗马人也熟悉这类梦。\Person{马尔库斯·图利乌斯（西塞罗）}从一个梦得知，一个当时还是小男孩、身份卑微、只是普通的\Person{尤利乌斯·屋大维}、且不为（西塞罗）本人所认识的人，就是命定的奥古斯都，是（罗马）内乱的镇压者和毁灭者。此事记载于\Person{维特里乌斯}的\WorkTitle{评注}中。但这种预言性的异象并不限于对最高权力的预言；它们也指示危险和灾难：例如，\Person{凯撒}因病未参加\Place{腓立比}战役，从而躲过了\Person{布鲁图}和\Person{卡西乌斯}的剑，然后他虽然预料会在战场上遭遇敌人更大的危险，却听从了\Person{阿托里乌斯}的异象而离开帐篷，从而逃脱（逃脱了敌人的俘获，他们不久后占领了帐篷）；又如，\Place{萨摩斯}的\Person{波利克拉特斯}的女儿从太阳的涂油和朱庇特的沐浴中预见到等待他的十字架刑。同样，在睡眠中显露出崇高的荣誉和卓越的才能；还发现疗法，揭露偷窃，指示宝藏。因此，\Person{西塞罗}的卓越在他还是个小男孩时就已被他的乳母预见到。从\Person{苏格拉底}胸中飞出的天鹅安抚众人，即其门徒\Person{柏拉图}。拳击手\Person{勒俄倪摩斯}在梦中被\Person{阿喀琉斯}治愈。悲剧诗人\Person{索福克勒斯}在梦中发现了从\Place{雅典}卫城丢失的金冠。悲剧演员\Person{尼奥普托勒摩斯}，通过睡眠中来自\Person{埃阿斯}本人的暗示，拯救了\Place{特洛伊}前\Place{罗提亚}海岸上英雄的坟墓免遭毁灭；当他移走腐朽的石头时，满载黄金而归。有多少注释家和编年史家为这一现象作证？有\Person{阿尔特蒙}、\Person{安提丰}、\Person{斯特拉托}、\Person{菲洛科鲁斯}、\Person{埃庇卡摩斯}、\Person{塞拉皮翁}、\Person{克拉蒂普斯}、\Place{罗得岛}的\Person{狄奥尼修斯}和\Person{赫尔米普斯}——整代的文献。我只会嘲笑这一切，如果我真该嘲笑那个幻想自己将说服我们\Person{萨图尔努斯}比任何人更先做梦的人；我们只能相信如果\Person{亚里士多德}（他巴不得帮助我们得出这样的意见）活在任何其他人之前。请原谅我发笑。\Person{埃庇卡摩斯}和\Place{雅典}的\Person{菲洛科鲁斯}确实将梦置于占卜中最高之位。全世界充斥着这类神谕：有\Place{奥罗普斯}的\Person{安菲阿拉俄斯}神谕、\Place{马卢斯}的\Person{安菲洛科斯}神谕、\Place{特罗阿德}的\Person{萨耳珀冬}神谕、\Place{维奥蒂亚}的\Person{特罗福尼乌斯}神谕、\Place{奇里乞亚}的\Person{摩普索斯}神谕、\Place{马其顿}的\Person{赫耳弥俄涅}神谕、\Place{拉科尼亚}的\Person{帕西法厄}神谕。此外，还有其他神谕，连同它们的原始基础、仪式和历史学家，以及整个关于梦的文献，\Place{贝里图斯}的\Person{赫尔米普斯}在五卷厚重的著作中会让你了解全部，甚至到厌腻的程度。但斯多葛学派很喜欢说，天主以祂最警觉的眷顾对每一制度的最警觉的眷顾，将梦赐给我们，作为占卜技艺和科学的其他保存手段之一，作为自然神谕的特殊支持。关于那些连我们自己也不得不给予信用的梦，就说这么多吧，尽管我们必须以另一种意义来解释它们。至于所有其他神谕，无人曾在其中做梦，我们还能对它们作何宣示，只能说它们是那些邪灵的阴谋，这些邪灵甚至在那时已居住在那些杰出人物自身之内，或者旨在复兴对他们的记忆，将其作为他们邪恶目的的单纯舞台，竟至仿效神的权能于他们的形状和形式之下，并以同样持久的邪恶，借着他们疗愈、警告和预见之类的恩赐来欺骗人——这恩赐的唯一效果，就是越帮助受害者越伤害他们；而他们提供帮助的手段，则通过将虚假之神的观念灌输到人的心中，使人们远离对真天主的追求。当然，如此有害的影响并不被封闭或限制在圣所和庙宇的范围内：它四处游荡，在空中飞翔，始终自由而不受约束。因此，没有人会怀疑我们的家园对这些邪灵是敞开的，他们不仅在人们的小教堂里，而且在他们的卧室里，用他们的幻象困扰人类猎物。\par

% structure: p0100 source=h2 target=section reason=relative-heading-rank; base=h2

\section{第四十七章 梦的各种分类。有些是来自天主的，如\Person{拿步高}的梦；另一些则仅是自然的产物。}

我们主张，梦主要是由魔鬼施加给我们的，尽管它们有时结果真实且对我们有利。然而，当它们怀着刚才所谈到的蓄意作恶的意图，采用谄媚而迷人的风格时，它们就相应地显得虚幻、欺骗、隐晦、放荡和不洁。无怪乎这些影像沾染了实物的特征。但是，从天主——祂确实曾应许\BibleQuote{将圣神的恩宠倾注在一切有血肉的人身上，并命令祂的仆人和使女们要看见异象且说出预言}\BibleRef{岳 3:1}——出发，所有那些可以与天主实际的恩宠相比拟的异象，即诚实、圣洁、先知性、受感动、教导性、引人向善的异象，其丰饶的本性甚至使它们满溢到不敬神的人身上，因为天主以极大的公正，\BibleQuote{使祂的雨水和阳光降给义人也给不义的人}\BibleRef{玛 5:45}，都必须被视为源自天主。\Person{拿步高}的梦确实是由天主的感动而来的；几乎人类大部分对天主的认识都是从梦中获得的。因此，正如天主的仁慈格外丰盛地赐予外邦人，同样，邪恶者的试探也遇到圣人们，他从未停止用恶毒的企图去偷袭他们，即使在他们的睡梦中，如果不能在醒时攻击他们的话。第三类梦将包括那些灵魂本身似乎因对特定情况的强烈专注而为自己产生的梦。既然灵魂不能自主地做梦（连\Person{厄庇卡玛斯}也持此见），它如何能成为任何异象的原因呢？那么，这一类梦就必须被归于自然的作用，而留给灵魂（即使在神魂超拔的状态中）承受所遇之事的能力吗？此外，那些显然既不出自天主，也不出自魔鬼的感动，也不出自灵魂，超越了普通预期、通常解释或可被清晰叙述的可能性，将必须在另一类别中归于纯粹而简单的神魂超拔状态及其特殊状况。\par

% structure: p0102 source=h2 target=section reason=relative-heading-rank; base=h2

\section{第四十八章 梦的原因与境况。什么最有助于有效做梦。}

据说，梦在夜末发生时更为可靠和清晰，因为那时灵魂的活力出现，沉睡离去。至于季节，春天做梦较为平静，因为夏天使灵魂松弛，冬天则使之坚硬；而秋天在其他方面有害健康，却因其果实的甘美而易于使灵魂软弱。再者，关于睡眠时身体的姿势，人不应仰卧，也不应右侧卧，也不应扭曲内脏，仿佛它们的腔体被反向拉伸：这样会导致心悸，或对肝脏的压力引起心灵的痛苦骚动。然而无论怎样，我认为这些都只是巧妙的推测而非确凿的证明（尽管推测者甚至是\Person{柏拉图}）；而且可能一切都无非是偶然的结果。但一般来说，如果梦是可以被引导的，那么它将受人的意志控制；因为我们不可考察一方面\Emph{意见}，另一方面\Emph{迷信}，在处理梦时，关于区分和改变各种食物所规定的内容。至于\Emph{迷信}，我们有一个例子：对打算接受必要睡眠以领受神谕的人规定禁食，以便这种禁欲产生所需的洁净；而\Emph{意见}的例子是，\Person{毕达哥拉斯}学派的门徒为了达到同样目的，拒绝豆类作为会加重胃并引起消化不良的食物。但三位兄弟，即\Person{达尼尔}的同伴，满足于仅有蔬菜，以避免国王膳肴的污染，从天主那里获得了其他智慧外，特别赐予了深入理解梦的恩赐。就我个人而言，我几乎不知道禁食是否只会让我做梦如此深沉，以至于我不清楚自己是否真的做了梦。那么好，你会问，清醒在这件事上不是有作用吗？当然，它在这件事上如同在整个主题中一样有关：如果它对迷信有贡献，那么对宗教贡献更大。因为甚至魔鬼也要求它们的做梦者如此修炼，以取悦它们的神性，因为它们知道这在天主面前是可悦纳的，正如\Person{达尼尔}（再次引用他）\BibleQuote{没有吃美味的饼}\BibleRef{达 10:2}长达三周。然而，这种禁欲他用来以谦卑取悦天主，而非为了在接收梦和异象之前引发灵魂的敏感和智慧，好像它不是为了在神魂超拔状态中产生这种行动似的。那么，这种\Emph{清醒}（这里我们有问题）与激发神魂超拔无关，而是更推荐它由天主所行。\par

% structure: p0104 source=h2 target=section reason=relative-heading-rank; base=h2

\section{第四十九章 没有灵魂天然免于做梦}

至于那些认为婴儿不做梦的人，因为他们以为灵魂在生命中的所有功能都是根据年龄的能力完成的，他们应当仔细观察婴儿在睡眠中的颤抖、点头和明亮的微笑，从这些事实理解到这些都是灵魂做梦时情感的流露，这些情感通过婴儿身体的娇嫩很容易显露到表面。然而，据说阿非利加的\Person{阿特兰特人}民族在夜晚陷入深沉的昏睡，这使他们受到指责，认为他们灵魂的构造有问题。现在，是那个偶尔诽谤蛮族的传闻欺骗了\Person{赫罗多德}，还是这样一群魔鬼在那些蛮族地区作主。既然\Person{亚里士多德}提到\Place{撒丁岛}的一位英雄，他习惯于夺去那些到他的神庙寻求灵感的人们的异象和做梦的能力，那么，夺取和赋予做梦的能力必定取决于魔鬼的意志和任性；由此可能产生了那个显著的事实（我们已经提及），即\Person{尼禄}和\Person{特拉绪墨得斯}直到晚年才做梦。然而，我们领受梦来自天主。那么，为何阿特兰特人未从天主领受做梦的能力？因为如今没有一个民族是与天主无份的，既然福音将其荣耀的光闪照世界直到地极？那么，难道是传闻欺骗了\Person{亚里士多德}，还是这种任性仍是魔鬼的行径？（让我们任取一种看法，）只是不可想象任何灵魂因其本性而免于做梦。\par

% structure: p0106 source=h2 target=section reason=relative-heading-rank; base=h2

\section{第50章 \Person{伊壁鸠鲁}的荒谬意见和异端\Person{梅南德}的亵渎妄想：论死亡，连\Person{哈诺客}和\Person{厄里亚}也被预留死亡}

到目前为止，我们已经足够论述了睡眠，这死亡的镜子和形象；也论述了睡眠的产物，即梦。现在让我们继续考察我们离世的起因——即死亡的安排和历程——因为我们不能对它不加追问和审视，尽管它本身就是一切问题和探究的终点。根据人类的一般情感，我们宣告死亡是\Quote{自然的债务。}这已由天主的声音确定；这是与一切受造物订立的契约：因此甚至\Person{伊壁鸠鲁}的冰冷自负也被反驳了，他说我们并不欠这样的债务；不仅如此，撒玛黎亚异端\Person{梅南德}的疯狂意见也被拒绝，他硬说死亡不仅与他的门徒毫无关系，而且实际上永远不会临到他们。他声称从上面一位的隐秘权力领受了这样的使命，凡领受他的洗礼的人都成为不朽的、不可朽坏的，并立刻领受复活的生命。我们无疑读到过许多非常奇妙的水：例如，林克斯提斯河的酒质水流使饮者陶醉；在科洛丰，能启发神谕的泉水使人疯狂；亚历山大被阿卡迪亚诺纳克里山的毒水杀死。此外，在基督时代之前的犹大有一个具有疗效的水池。众所周知，诗人曾纪念过斯堤克斯沼泽能使人免死；尽管\Person{德蒂斯}，虽有那防腐剂，仍得哀悼她的儿子。至于此事，即使\Person{梅南德}本人投入这著名的斯堤克斯河，他终归还是要死；因为你必须来到斯堤克斯河，根据一切记述，它位于死者区域。那么，那些有福而迷人的水是什么，在哪里？连\Person{若翰洗者}在他先前的职事中也未曾使用，\Person{基督}在他之后也未曾向祂的门徒启示过？\Person{梅南德}的这个奇妙洗礼是什么？我想他是个滑稽家伙。但为什么（这样的洗礼）如此少有需求，如此隐匿，很少人前往那里清洗？我实在怀疑在如此罕见的圣事中，竟附带着如此大的保障和安全，甚至免除在侍奉天主本身时的死亡通常规律，然而，相反地，万民都要\BibleQuote{上主的山和雅各伯天主的殿，}祂要求祂的圣人藉殉道经历死亡，这死亡祂甚至要求于祂的基督。没有人会将这样的影响归于魔法，即能免除死亡，或更新并赋予生命活力，像葡萄树通过状况的更新那样。这样的能力连伟大的\Person{美狄亚}本人也未获得——至少对人是如此，如果允许她对一头愚蠢的羊施展的话。毫无疑问，\Person{哈诺客}被提升，\Person{厄里亚}也一样；\BibleRef{列下 2:11}他们没有经历死亡：它被推迟了，（而且只是推迟，）确实如此：他们被预留承受死亡，好藉他们的血消灭敌基督。\BibleRef{默 11:3}就连\Person{若望}也经历了死亡，尽管关于他曾有一种无根据的期待，认为他会活着直到主的来临。\BibleRef{若 21:23}异端确实多半从我们自己提供的例子中仓促产生：它们从它们所攻击的同一地方获取防御的盔甲。简言之，整个问题归结为这个挑战：\Person{梅南德}自己所洗礼的那些人在哪里？他投入他的斯堤克斯河的那些人在哪里？让他们出来站在我们面前——他那些被他变成不朽的宗徒？让我（多疑的）\Person{多默}看见他们，让他听见他们，让他触摸他们——他便会相信。\par

% structure: p0108 source=h2 target=section reason=relative-heading-rank; base=h2

\section{第51章 死亡完全将灵魂与身体分离}

死亡的作用是明显易见的：即灵魂与身体的分离。然而，有些人对灵魂不朽的认识很薄弱（因为他们未受天主教导），竟用如此贫乏的论据来维护，以至于他们愿意相信某些灵魂在死后仍依附于身体。正是在这个意义上，尽管\Person{柏拉图}随意将他喜欢的灵魂直接送入天堂，但在他的\WorkTitle{理想国}中，他向我们展示了一具未埋葬的尸体，这尸体由于灵魂（如他所说）未与身体分离而长时间保持不腐。对此，\Person{德谟克利特}也提到，在坟墓中人的指甲和头发还会生长相当一段时间。现在，前述尸体的保存很可能是由于大气性质所致。如果空气特别干燥，土壤含盐，那会如何？如果身体本身也异常干燥，那会如何？再者，如果死亡的方式已从尸体中消除了所有腐败物质，那会如何？至于指甲，由于它们是神经的起点，当神经本身松弛延展时，指甲似乎会变长，并随着肌肉萎缩而越来越突出。头发则从大脑获得滋养，这会使它因大脑的秘密滋养和保护而长久存在。事实上，在活人身上，头发的多少也与大脑的丰盈程度成正比。你们有医生可以作证。但灵魂的任何微粒都不可能留在身体内，因为身体本身也注定要在时间消除了它整个活动场景后消失。然而，即使在有些人的意见中，灵魂的这种部分存留仍占有一席之地；因此他们不愿在葬礼上用火焚烧身体，因为他们想保存灵魂的微小残余。但这种虔敬的处理还有另一种解释：并非为了照顾灵魂的残余，而是为了对身体本身避免残忍的习俗，因为身体既然是人的，就不该遭受甚至加在杀人犯身上的结局。事实上，灵魂是不可分的，因为它是不朽的；（这一事实）迫使我们相信死亡本身也是一个不可分的过程，不可分地加于灵魂——不是因为灵魂不朽，而是因为它不可分。但如果灵魂可分为部分，那么死亡的作用就必须被分割，其中一部分要留待以后的死亡。这样，一部分死亡将不得不为一部分灵魂留下。我并非不知道这种观点的痕迹仍然存在。我是从我认识的一个人的经历中发现的。我知道一个事例，有一位妇女，是基督徒父母的女儿，在她青春和美貌的鼎盛时期，经历了一段极其幸福但短暂的婚姻生活后，在耶稣内安眠。在她被安葬之前，当司铎开始规定的礼仪时，就在他祈祷的第一口气息中，她将双手从身体两侧移开，摆出祈祷的姿势，并在圣礼结束后，又将双手放回身体两侧。此外，在我们中间还有一个广为人知的故事：在一个墓地里，一具尸体自动让出位置，以便另一具尸体靠近它。如果（事实上也是如此）在外教人中也流传着类似的故事，（我们只能断定）天主处处彰显祂自己大能的记号——对自己的人民是安慰，对陌生人则是见证。我宁愿相信这种奇异现象是直接来自天主的作为，而非出于灵魂的残余：因为如果有残余，它们肯定会移动其他肢体；即使它们移动了双手，也不会是为了祈祷。尸体也不会仅仅满足于为邻居让路：它也会通过改变位置而对自己有益。但无论这些现象出于何种原因（你们应将其归入征兆和奇事之中），它们不可能规范自然。死亡，如果一旦在作用上不完全，就不是死亡。如果任何一部分灵魂存留，那就构成了存活状态。死亡不再与生命混合，如同黑夜不再与白昼混合。\par

% structure: p0110 source=h2 target=section reason=relative-heading-rank; base=h2

\section{各种死亡都是对自然的强暴，来自罪——罪是对天主所创造之自然的侵犯}

如此，死亡的工作便是——灵魂与身体的分离。排除命运和偶然情况，按照人的观点，它被区分为两种形式——平常的与异常的。平常的他们归因于本性，在每个人的死亡中行使其平静的影响；异常则被认为违反本性，发生在每一次暴死中。至于我们自己的见解，我们确知人的起源，并且我们大胆断言并坚持：死亡不是作为人的自然后果发生的，而是由于一个并非本性自身的过失和缺陷；虽然，毫无疑问，很容易将\Emph{自然}一词应用于那些似乎（尽管由于外在原因的出现）从一开始就与我们不可分离的过失和境况。如果人直接被指定在受造的条件下死去，那么死亡当然必须归于本性。然而，他并非这样被指定去死，这由那法律本身证明，该法律使其状况取决于一个警告，而死亡则来自人的任意选择。事实上，如果他没有犯罪，他肯定不会死。那不能是本性，因为它是在提出一个选择之后，通过意志的行使而发生的，而不是由必然性——一个不可动摇和不可改变的状态的结果——所致。因此，尽管死亡有各种结局，因为它的原因是多样的，我们不能说最容易的死亡是如此温和，以至于不是通过暴力（对我们的本性）发生的。产生死亡的法律本身，尽管简单，却是暴力。难道还能有其他可能吗？当灵魂和身体如此密切的结合，这两个姐妹实体从成胎起就如此不可分割地共同成长，却被撕裂和分离。因为，虽然一个人可能因喜乐而断气，如\Person{斯巴达人基隆}在拥抱他刚在奥林匹克运动会上获胜的儿子时；或因荣耀，如\Person{雅典人克利德莫斯}在因其历史著作的卓越而接受金冠时；或在梦中，如\Person{柏拉图}；或因一阵大笑，如\Person{普布利乌斯·克拉苏}——然而死亡太过暴力，它借助奇怪而异己的手段降临到我们身上，用其独特的方式驱逐灵魂，在我们可能过着一种喜乐和尊荣、平安和愉悦的快乐生活之时，召唤我们去死。这对于船只同样是一种暴力：尽管远离\Place{卡法瑞亚岩}，没有遭受风暴，没有巨浪击碎它们，顺风，滑行，欢乐的船员，它们却在完全安全中沉没，突然地，由于某种内部冲击。生命的沉船——甚至平静死亡的结局——也并非不同。人的身体之船，无论是以未破损的木板航行还是被风暴摧毁，如果灵魂的航行被颠覆，那都无关紧要。\par

% structure: p0112 source=h2 target=section reason=relative-heading-rank; base=h2

\section{第53章。整全灵魂不可分，直至生命最后行动；绝不会部分或零碎地离开身体。}

然而，灵魂最终将寄居何处，当它赤裸地脱离身体之后？在我们的探究次序中，我们当然不能犹豫，必须跟随它到那里。不过，我们首先必须充分陈述属于当前主题的内容，以免有人因为我们提到了死亡的种种结局，就期望我们特别描述这些结局——这些事本应留给医生，他们才是判断与死亡或其原因相关事件以及人体实际状况的合适专家。当然，为了在讨论这一主题时保持灵魂不朽的真理，我在提及死亡时，将不得不引入关于解散的措辞，其含义似乎暗示灵魂是逐渐地、一块一块地逃逸的；因为它带着衰败的一切迹象撤离（身体），看似在消亡，并使我们产生它因缓慢离去而被消灭的想法。但这一现象的全部原因在于身体，并源自身体。因为无论何种死亡（作用于人），都无疑摧毁了生命力所依赖的物质、区域或通道：物质如胆汁和血液；区域如心脏和肝脏；通道如静脉和动脉。既然如此，身体的这些部分各自因其特有的损伤而遭到破坏，直至生命力的最终毁灭——即自然目的、部位和功能的终结——必然会发生的是，在其工具、居所和空间逐渐衰败的过程中，灵魂本身也被迫放弃一个又一个部分，呈现出似乎缩减至无的样子；就仿佛一个车夫，当他的马匹因疲劳而收回其精力时，人们便以为他本人也失败了。但这种假设仅适用于被剥夺者所处的境况，而非任何真实的受苦状况。同样，身体的御车者——动物精神——因其载具的失败而失败，而非其自身；它放弃了工作，但并未放弃活力；在运作上萎靡，但在本质上并未改变；在偿付能力上破产，但在实体上并未枯竭——因为它不再显现，但并未停止存在。因此，每一种快速死亡——例如斩首、折断脖颈（这会立即为灵魂打开一个巨大的出口）、或突然的崩溃（如内源性崩溃中风），它一举粉碎所有生命活动——都不会拖延灵魂的逃逸，也不会将其分离痛苦地分成若干时刻。然而，当死亡是缓慢拖延时，灵魂便以它自己被遗弃的方式放弃其位置。但它并非在这个过程中被分成碎片：它是被缓慢拉出的；在这样被抽出的过程中，它使最后残留的部分看似只是它的一部分。然而，没有任何一部分应当被认为是可分离的，因为它是最后的；也不应因为它微小而认为它能够解体。与整体序列一致的是其终点，中间部分延伸至极端；残余部分与整体凝聚，被整体等待，但从未被整体抛弃。我甚至敢说，整体的最后一部分就是整体本身；因为它虽然更小且最后，却属于整体并完成整体。因此，确实，许多次发生的是，灵魂在分离时反而更加有力地激动，以更焦虑的目光和更快的言辞；因为它现在处于更高且更自由的位置，通过仍然滞留在肉体中的最后残余，它宣告它所看见、它所听见以及开始知道的事物。按\Person{柏拉图}的话说，身体是监狱，但按宗徒的话，它是\Quote{天主的圣殿}，因为它是在\Person{基督}内。然而，（必须承认）由于其封闭，它阻碍并遮蔽灵魂，并通过肉体的凝结玷污灵魂；因此，照亮物体的光芒以更混乱的方式进入灵魂，仿佛透过角窗。无疑地，当灵魂藉死亡的力量从与肉体的结合中释放时，正是因这释放而洁净并纯净：此外，确实，它从肉体的面纱中逃入开阔的空间，进入其清澈、纯净和内在的光；然后发现自己享受着从物质的解放，并凭借其自由恢复了其神性，如同一个从睡眠中醒来的人从影像走向真理。于是它宣告它所看见的；于是它或喜乐或恐惧，视它发现为自己预备的住所，当它一见到天使的面容——灵魂的指控者，诗人们所说的\Person{墨丘利}。\par

% structure: p0114 source=h2 target=section reason=relative-heading-rank; base=h2

\section{第五十四章 灵魂离开身体后往何处去？哲学家的意见或多或少荒谬。柏拉图的哈德斯}

因此，关于灵魂被撤往何处的问题，我们现在给出答案。几乎所有主张灵魂不死的哲学家，尽管对此有各自特殊的观点，仍然为它主张这一（永恒的状态），如\Person{毕达哥拉斯}、\Person{恩培多克勒}和\Person{柏拉图}，以及那些允许它从离开肉体到万物焚毁之间有某种延迟的人，还有斯多葛派，他们只把自己的灵魂，即智者的灵魂，安置在高处的居所。的确，\Person{柏拉图}并不允许所有灵魂，甚至不是所有哲学家的灵魂，都得到这归宿，而只允许那些因对少年的爱而修养其哲学的人。这就是不洁在哲学家手中获得的特权！在他的体系中，智者的灵魂被提升到高处的以太中：按\Person{阿里乌}的说法，进入空气；按斯多葛派的说法，进入月亮。事实上，我感到惊奇的是，他们把不智者的灵魂遗弃在地上，而他们却肯定这些灵魂也由智者（比他们优越得多）教导。因为，在他们之间相隔的巨大空间中，哪里有学校可以让他们受教导？灵魂门徒是如何到达他们的教师那里的，当它们被如此遥远的间隔分开时？还有，在他们死后的状态中，任何教导能给他们带来什么益处呢，既然他们正濒临被普世大火毁灭？所有其他灵魂他们推进阴府，\Person{柏拉图}在他的\WorkTitle{斐多篇}中描述阴府为大地之胸怀，世上一切污秽在那里积聚、沉降、散发，每一股单独的空气只使沸腾的杂质的密度更大。\par

% structure: p0116 source=h2 target=section reason=relative-heading-rank; base=h2

\section{第五十五章 基督徒对阴府位置的看法；死后立即享有的乐园真福。殉道者的特恩。}

我们并不认为阴府的较低区域是一个空洞，也不是世界的地下排污管，而是地球内部的一个巨大深空，是其深处的一个隐秘凹处；因为我们读到基督在死亡中在地心度过了三天，\BibleRef{玛 12:40}也就是说，在隐秘的内凹处，隐藏在地内，被地包围，并叠加在更深的下层深渊之上。如今，虽然基督是天主，但祂也是人，\Quote{祂照圣经死了，}\BibleRef{格前 15:3} 以及\Quote{也照同一圣经被埋葬了。}祂完全遵守了祂存在的同一法则，以死人的形态和状况留在阴府；祂没有先升上高天，就降到了地下的部分，为在那里使圣祖和先知们分享祂自己。\BibleRef{伯前 3:19}（既然如此），你必须认为阴府是一个地下区域，并远离那些骄傲得不肯相信信友的灵魂配得较低区域的人。这些人，他们是\Quote{仆人超过主人，门徒超过师傅，}\BibleRef{玛 10:24} 如果他们必须在亚巴郎怀中期待复活，无疑会拒绝接受复活的安慰。但他们说，基督下降阴府正是为此，使我们自己不必下到那里。那么，如果死后同样的监狱在等待所有人，外教人和基督徒之间有什么区别？的确，当基督已经坐在圣父的右边时，灵魂如何升到天上，而当还没有按照天主的命令听到总领天使的号筒，——当主来临时在地上找到的那些人还没有被提到空中迎接祂的来临，\BibleRef{得前 4:17} 与在基督内死了的人一同，他们将是首先复活的？\BibleRef{得前 4:16} 没有人为谁敞开天堂；大地对他是安全的，我不会说大地对他关闭。的确，当世界过去后，天国将打开。那么，我们是否必须与\Person{柏拉图}那些爱少年的贤人一起在高处的以太中睡眠？或与\Person{阿里乌}一起在空气中？或与斯多葛派的恩底弥翁们一起在月亮周围？不，你告诉我，而是在乐园中，圣祖和先知们已经在主复活的随从中从阴府迁到了那里。那么，为什么乐园的区域，正如在圣神中启示给若望的，位于祭台下面，\BibleRef{默 6:9} 其中除了殉道者的灵魂之外没有显示其他灵魂？为什么最英勇的殉道者\Person{佩佩图阿}在她受难之日，在她所得到的乐园启示中，只看到她的同伴殉道者在那里，难道不是因为守卫入口的剑不允许任何人进去，除了那些在基督内、不在亚当内死去的人？一种为天主的新死亡，甚至是为基督的非凡死亡，被接纳进入有死的接待室，特别为接纳新来者而改变和适应。那么，注意外教人和基督徒在死亡上的区别：如果你必须为天主舍命，正如护慰者所劝告的，那不是在于温和的发热和柔软的床榻，而是在于殉道的剧痛中：你必须背起十字架跟随你的师傅，正如祂亲自教导你的。\BibleRef{玛 16:24} 打开乐园的唯一钥匙是你自己的鲜血。我们有一篇论著，（论乐园），在其中我们确立了每个灵魂都被安全地拘留在阴府中直到主的日子。\par

% structure: p0118 source=h2 target=section reason=relative-heading-rank; base=h2

\section{第五十六章 驳斥\Person{荷马}关于因身体未被埋葬而灵魂被扣留在阴府的观点。并驳斥那些过早离开身体的灵魂必须等待才能进入阴府的观点。}

有一个问题出现：这是否在灵魂离开身体后立即发生？是否有些灵魂因特殊原因而被暂时留在地上？以及是否允许它们自愿或通过权威干预，在之后某个时间从\OriginalTerm{阴府}{Hades}被移除？甚至这样的意见也绝不缺乏人自信地提出。曾有人相信，未埋葬的死者在其得到适当埋葬之前，不被允许进入\OriginalTerm{阴间}{infernal regions}；正如\Person{荷马}笔下的\Person{帕特罗克洛斯}，他在梦中恳求\Person{阿喀琉斯}为他举行葬礼，理由是他无法通过任何其他入口进入阴府，因为已埋葬的死者的灵魂不断将他推开。我们知道\Person{荷马}在此处展现了超越诗意的许可；他着眼于死者的权利。的确，与他对于坟墓正当尊荣的关心相称的是，他对那种对灵魂有害的埋葬拖延的谴责。（他也意图加上一个警告），即任何人都不应通过将友人的尸体扣留在自己家中，使自己与死者一起，因他所滋养的伴有痛苦与忧伤的慰藉的失常，而遭受更大的损伤和困扰。因此，他在描绘一个未埋葬灵魂的抱怨时，心中持有双重目的：他希望透过及时办理葬礼来维持对死者的尊荣，同时也缓和由对他们的记忆所激起的悲伤情绪。但是，总而言之，设想灵魂能承担身体的仪式和要求，或将其中任何一样带入阴间，这是多么虚妄！而且，如果认为灵魂因本应视为恩惠的埋葬疏忽而遭受损伤，那就更加虚妄了！因为无疑，一个不愿死去的灵魂很可能宁愿尽可能晚地迁往阴府。然而，如果确实灵魂因身体延迟埋葬而遭受损伤——且损伤的要害在于埋葬的疏忽——那么，让那不可能归咎于过失延迟的一方承担全部损伤，是极其不公平的，因为当然所有过失都在死者的至亲身上。他们还声称，那些因夭亡而被夺走的灵魂四处游荡，直到他们完成了若非早逝本会度过的余下岁月。如今，要么每个人的日子被个别规定，如果这样规定，我不能设想它们能被缩短；要么，尽管有这样的规定，它们可能因天主的旨意或其他强大影响而被缩短，那么（我说）这种缩短是没有效力的，如果它们仍能以某种其他方式完成的话。另一方面，如果它们没有被规定，那么对于未规定的时期，就没有任何余数需要完成。我还要再说一点。假设一个婴儿死在襁褓之中；或者可能是一个未成熟的幼童；或者又可能是一个到了青春期的少年：此外，假设每个案例中的生命本应达到完整的八十年，那么任何一个灵魂在死后失去身体后，如何可能在地上度过全部缩短的岁月呢？人的年龄不能没有身体而度过，因为正是借助身体，生命的时期才有其职责和劳作得以履行。此外，让我们自己的人记住，灵魂在复活时要收回他们死去时所拥有的同样身体。因此，我们必须期待我们的身体恢复同样的状况和同样的年龄，因为正是这些特性赋予身体其特殊样式。那么，一个婴儿的灵魂如何能在地上度过其剩余岁月，以致在复活时能够具有八十岁老翁的状态，尽管它几乎只活了一个月？或者，如果必须在地上完成规定的生命日子，那么同样历经沧桑的生命历程（它本身被规定伴随那些日子）是否也必须由灵魂与日子一起经历？它是否必须从婴儿期到幼童期的过程中从事学业？在青春和早期成年的兴奋与活力中当兵？在成熟成年与老年之间的更庄重岁月中承担严肃的司法职责？它是否必须从事贸易以获利，用锄头和犁翻土，航海，提起诉讼，结婚，辛劳工作，生病，以及经历岁月流逝中等待它的各种吉凶祸福？但是，没有身体，所有这些事务如何管理？生活（度过）没有生命？但是（你会告诉我）所说的指定时期没有任何事件，只是单纯地流逝而完成。那么，什么阻止它在阴府中完成呢？在那里你绝对无法应用它。因此，我们坚持每个灵魂，无论离开身体时的年龄如何，都保持不变，直到时间来临，那时所应许的完美将在一种与无与伦比的天使相称的适当状态中实现。因此，那些灵魂必须被视为在阴府中流亡，人们往往认为是被暴力夺走的，尤其是残酷的折磨，如十字架、斧头、刀剑和狮子；但我们不认为那些由正义（那暴力的报应者）所判定的死亡是横死。那么，你会说，所有邪恶的灵魂都被放逐在阴府中。（且慢，这是我的回答。）我必须强迫你确定（你所谓阴府是指什么），即其两个区域中哪一个，善者的区域或恶者的区域。如果你指的是恶者，（我只能说，）即使现在，邪恶的灵魂也配被交付到那些住所；如果你指的是善者，那么你为什么认为婴儿、贞女以及那些因其生活状况而纯洁无罪的灵魂不配有这样的安息之所呢？\par

% structure: p0120 source=h2 target=section reason=relative-heading-rank; base=h2

\section{第五十七章 魔术与巫术只是表面的效果。唯有天主能复活死人}

滞留在这些阴间地域，与那些过早离世的\Emph{Aori}（即被强行带走的灵魂）为伴，要么是件极好的事，要么确实是件极坏的事——与那些横死的\Emph{Biaeothanati}在一起。请允许我使用那些魔术师——一切古怪意见的发明者——所用的实际词语和术语，连同其\Person{Ostanes}、\Person{Typhon}、\Person{Dardanus}、\Person{Damigeron}、\Person{Nectabis}和\Person{Berenice}。有一篇众所周知的通俗作品，它承揽从阴府召出那些确实活满天年、经光荣死亡、甚至以完备仪式和适当礼仪埋葬的灵魂。此后我们该对魔术说什么？当然，说它几乎人人都说它是骗局。但这套骗术并非只有我们基督徒才识破。我们确实发现了这些邪神，不是通过与之合谋，而是通过某种敌对的知识；我们也不是用吸引它们的方式，而是用制服它们的能力来处理（它们那可怜的系统）——那人类心灵的多样祸害、一切谬误的工匠、一举毁灭我们救恩和灵魂者。就这样，即便通过魔术——它其实只是第二号的偶像崇拜，其中他们声称死后变成魔鬼，正如在起初字面的偶像崇拜中他们被认为变成神（何不呢？因为神不过是死物）——前述的\Emph{Aori Biaeothanati}竟被呼召——这并非不合理，如果一个人基于这个原则建立信德，即那些被残忍过早死亡强行带走的灵魂，最易于迷恋暴力与不义，并且热切渴望报复，这明显是可信的。然而，在这些灵魂的掩护下，魔鬼在运作，尤其是那些生前曾居住在他们内、并驱使他们走向最终夺命命运的魔鬼。因为，正如我们已经暗示的，几乎没有哪个人不被魔鬼伴随；众所周知，人们归咎于意外的过早横死，实则由魔鬼造成。我相信，我们可以通过实际事例证明恶神潜伏在死者身上的这种骗局，因为在驱魔时，（恶神）有时自称是附身者的一位亲属，有时是角斗士或\Emph{bestiarius}（斗兽士），有时甚至是神；它一贯的首要关切就是消灭我们正在宣告的真理，使人们不肯轻易相信所有灵魂都迁往阴府，并摧毁对复活和审判的信德。然而尽管如此，魔鬼在试图迷惑旁观者之后，终被天主恩宠的压力所征服，极不情愿地承认全部真相。同样，在另一种魔术中——它被认为从阴府召出现在安息在那里的灵魂，并展示给公众观看——没有比这更有效的骗术了。当然，幻影之所以可见，是因为也有身体附着其上；要蒙蔽一个人的外在视觉并不难，因为要蒙蔽他的心眼更容易。从魔术师杖中出来的蛇，确实在法老和埃及人面前显现为实体。诚然，\Person{Moses}的真实吞没了它们的虚假欺骗。\BibleRef{出 7:12}巫术士\Person{Simon}和\Person{Elymas}也曾对宗徒施过许多攻击，但打击（他们）的失明并非咒术师的把戏。污神伪造真理的努力有什么新奇呢？就在此时，同一\Person{Simon}（Magus）的异端追随者们被他们那艺术的夸张自诩所鼓舞，以至于他们承揽从阴府召出先知们的灵魂。我想他们能通过虚假奇迹做到这一点。因为，确实古时曾允许皮同（或腹语）之神做同样的事——当\Person{Saul}在失去（永生）天主后求问死者时，它甚至代表\Person{Samuel}的灵魂。然而，天主禁止我们以为任何圣人，更不用说先知的灵魂，能被魔鬼从（阴府的安息处）拖出来。我们知道\Quote{\Person{Satan}自己伪装成光明的天使}\BibleRef{格后 11:14}，更不用说变成光明的人，而且最后他将\Quote{甚至显示自己是神}\BibleRef{得后 2:4}，并显出\Quote{大奇事和神迹，以致如果可能，连被选者也要被迷惑}\BibleRef{玛 24:24}。在前述场合，他几乎毫不犹豫地自称是天主的先知，尤其对\Person{Saul}——他那时正实际居住在后者内。你不可想象制造幻影的是一人，求问幻影的是另一人；而是同一个灵——既在女巫中也在背道（君王）中——轻易假装显现它已准备让他们相信为真实的东西——（这灵）通过其邪恶影响，\Person{Saul}的心定在哪里，他的财宝就在哪里，而那里当然没有天主。因此，他看见了他信靠其帮助而将要看见的那位，因为他信靠使他看见的那位。但有人反驳说，在夜间的异象中，死人常被看见，且带有明确目的。例如，\Person{Nasamones}（纳萨莫那人）通过频繁长期拜访亲属的坟墓来咨询私人神谕，如\Person{Heraclides}、\Person{Nymphodorus}或\Person{Herodotus}所述；\Person{Celts}（克尔特人）为同样目的在勇士首领的墓旁过夜，如\Person{Nicander}所确证。那么，我们承认梦中死人的显现并不比活人的显现更真实；但我们用同样的标准衡量所有情形——对死者与生者，乃至一切所见现象。事物并非因为显现为真就是真的，而是因为被充分证明为真。梦的真理由实现而非外观来宣告。此外，阴府绝不为（逃出）任何灵魂而开启，这一事实已由主在\Person{Abraham}身上，通过祂对安息穷人和受苦富人的描绘所牢固确立。（祂说）没有人可能从那些居所被派来向我们报告下界的情况——这种目的（若可能）在那个场合或许是可允许的，以说服人相信摩西和先知。无疑，天主的能力有时曾将人的灵魂召回身体，作为祂超越权利的明证；但绝不能因此认为神圣的信德与巫术士的傲慢自诩、与梦的欺骗、与诗人的放纵之间有任何可能的协议。然而，在一切真正复活的情形中，当天主的能力通过先知、或基督、或宗徒将灵魂召回身体时，我们便从（复活身体的）坚实、可触及且确凿的事实得到一个完全的推定，即其真实形态必须如此，以致迫使人相信一切无形死人的显现都是虚假的。\par

% structure: p0122 source=h2 target=section reason=relative-heading-rank; base=h2

\section{第五十八章 结论。待议要点。所有灵魂在阴府被拘留直至复活，预尝其最终痛苦或福乐。}

因此，所有灵魂都被关在阴府内：你承认这一点吗？（这是真的，无论）你说\Quote{是}或\Quote{否}：而且，在那里已经经历了惩罚与安慰；那里你看到一个穷人和一个富人。现在，我推迟了一些零散问题到本部分，我将在适当的地方提及它们，然后结束。那么，你为什么不能设想，灵魂在阴府中，在等待审判的抉择时，处于某种预尝黑暗或光荣的状态，经受惩罚与安慰呢？你回答：因为在天主的审判中，其事情应当是确定而稳妥的，也不应有任何对祂判决结果的事先预知；此外，灵魂必须先被复活的肉身之衣覆盖，这肉身作为其行为的伙伴，也应分享其报偿。那么，在这间隔中会发生什么呢？我们要睡觉吗？但灵魂在人们活着时也不睡觉：确实是身体睡觉的事务，死亡本身也属于此，并不亚于睡眠这面镜子与仿品。或者，你愿意认为，在整个人类被吸引而去、所有人的期待被保管起来的那里，什么也没有发生？你认为这状态是审判的预尝，还是其实际开始？是过早的侵犯，还是其全面执行的第一道菜？现在，真的，甚至在阴府中，如果罪人在那里仍然一切顺利，而义人却还不顺利，这岂不是最大的不公义？怎么，你希望死后希望更加混乱吗？你希望它用不确定的期待更加嘲弄我们吗？还是它现在成为对过去的回顾和审判的安排，带着不可避免的颤抖恐惧之感？但是，再次，灵魂必须总是等待身体才能经历悲伤或喜乐吗？甚至它自己，独自承受这两种感觉，难道还不够吗？多少次，在没有身体任何痛苦的情况下，灵魂单独被坏脾气、愤怒、疲乏折磨，而且常常甚至自己不知不觉？另一方面，在身体受苦时，灵魂多少次为自己寻找隐秘的喜乐，暂时从身体的纠缠中退隐？如果灵魂不是惯于独处，欢欣且荣耀于身体的酷刑，那我就错了。例如，看穆提乌斯·\Person{斯凯沃拉}的灵魂，当他在火上熔化右手时；也看\Person{芝诺}的灵魂，当\Person{狄奥尼修斯}的折磨掠过它时。野兽的咬伤对年轻英雄是荣耀，如同\Person{居鲁士}身上的熊的伤疤。那么，灵魂即使在阴府中，也完全知道如何无需身体而欢欣和悲伤；因为在肉身中，当它愿意时，它感到痛苦，尽管身体未受伤；当它愿意时，它感到喜乐，尽管身体在痛苦。如果这些感觉在世时随其意愿发生，那么死后怎会不因天主的审判定命而发生呢！此外，灵魂并非所有操作都借助肉身的服务；因为天主的审判甚至追逐单纯的思想和最微不足道的意愿。\BibleQuote{凡注视妇女，有意贪恋她的，他已在心里奸淫了她。}\BibleRef{玛 5:28}因此，即使为此原因，灵魂完全不必等待肉身，就应为它没有肉身合伙所做的事而受罚，是极其恰当的。同样，照此原则，作为对虔诚和善意思想的回报，这些思想没有肉身帮助，灵魂应在没有肉身的情况下得到安慰。不仅如此，即使是通过肉身做的事，灵魂也是最先构思、最先安排、最先授权、最先促进行动的。即使它有时不愿行动，它仍然是首先处理它打算借助身体来实现的对象。事实上，已完成的事实绝不可能先于其心理构思。因此，完全符合这种秩序的是，我们本性的这一部分应首先获得报酬和奖赏，因其优先性。简而言之，既然我们理解福音中指出的\Quote{监狱}是阴府，\BibleRef{玛 5:25}并且我们也解释\Quote{最后的一文钱}意味着最小的过犯，必须在复活前在那里偿还，没有人会犹豫相信，灵魂在阴府中接受某种补偿性的训诫，并不妨碍复活的全过程，届时报偿还将通过肉身施行。这一点，圣神也通过最频繁的劝诫让我们注意，只要我们中有人从祂应许的属灵启示的知识中承认祂话语的力量。现在最后，我相信，我已经遭遇了关于灵魂的每一种人类意见，并通过（我们神圣信仰的）教导检验了其性质，我们满足了那种只是合理且必要的好奇心。至于那些过度和无聊的，其信息将永远有同样大的缺陷，就如研究和探索中的夸大和任性一样大。\par

\SourceRecord{Translated by Peter Holmes.；Ante-Nicene Fathers；Vol. 3.；Edited by Alexander Roberts, James Donaldson, and A. Cleveland Coxe.；Buffalo, NY: Christian Literature Publishing Co.,；1885.；<http://www.newadvent.org/fathers/0310.htm>.}
